% !TEX program = xelatex
\IfFileExists{main.tex}{\documentclass[main.tex]{subfiles}}{
  \IfFileExists{../../main.tex}{\documentclass[../../main.tex]{subfiles}}{
    \IfFileExists{../main.tex}{\documentclass[../main.tex]{subfiles}}{
      \documentclass[Book\space NL/main.tex]{subfiles}
    }
  }
}
\begin{document}
\ifSubfilesClassLoaded{
  \ifcsname appendixchapterparent\endcsname
  \else
    \standaloneappendixchaptersetup{D}{6}
  \fi
}{}
\edef\reviewnotesimagepath{\booknlroot appendices/appendix 4-Review Notes/images/}
\chapter{Renal, Urinary, and Reproductive Systems}
\labch{appendix-d6-renal-urinary-reproductive-review}
\margintoc
\noindent\textit{Transcribed from the supplied renal, urinary, and reproductive study notes. The content is grouped by topic rather than lecture number; spelling and grammar are lightly copyedited, while source claims are retained and accuracy concerns are flagged in bold where identified.}\par\medskip

\section{Renal Structure and Physiology}

\subsection{Histology of urinary system}
\begin{itemize}
\item LOs:
\begin{itemize}
\item Major organs of urinary system \(\rightarrow\) microscopic structure (nephron, collecting duct, renal pelvis)
\item Renal blood supply
\item Podocyte filtration slit diaphragm
\item Microscopic structure of ureters, urinary bladder, urethra (M, F)
\end{itemize}
\item Overview (ไม่น่าจะออก?)
\begin{itemize}
\item Urinary system organs: kidney \(\rightarrow\) ureters \(\rightarrow\) urinary bladder \(\rightarrow\) urethra
\item Work to regulate: water, ions, pH, metabolites + excrete other substances
\item Kidney: RAAS (renin angiotensin I/II, aldosterone system) - regulates blood pressure, hematopoietin (new RBC stimulation)
\end{itemize}
\end{itemize}
\begin{itemize}
\item Gross structure:
\end{itemize}
\begin{center}\includegraphics[width=0.82\textwidth,height=0.68\textheight,keepaspectratio]{\reviewnotesimagepath d6-figure-histology-urinary-gross-structure.png}\par\captionof{figure}{The illustration relates to Gross structure:, within the topic of Histology of urinary system.}\labfig{appendix-figure-d6-histology-urinary-gross-structure}\end{center}
\begin{itemize}
\item Kidney \(\rightarrow\) ‘bean shaped’ organ, bilateral at T12-L3 (R lower than L)
\item Retroperitoneal organs (see SADPUCKER mnemonic for reference)
\item Internal (gross) structure of the kidney
\end{itemize}
\begin{center}\includegraphics[width=0.82\textwidth,height=0.68\textheight,keepaspectratio]{\reviewnotesimagepath d6-figure-histology-urinary-internal-gross-structure-kidney.png}\par\captionof{figure}{The illustration relates to Internal (gross) structure of the kidney, within the topic of Histology of urinary system.}\labfig{appendix-figure-d6-histology-urinary-internal-gross-structure-kidney}\end{center}
\begin{itemize}
\item Drainage of urine - renal papilla/papillary duct/duct of Bellini \(\rightarrow\) minor calyx \(\rightarrow\) major calyx \(\rightarrow\) renal pelvis \(\rightarrow\) ureter \(\rightarrow\) urinary bladder \(\rightarrow\) urethra
\item Regions: cortex (outer, dark brown part), medulla (red part with drainage) \(\rightarrow\) forms renal pyramids with renal columns (blood vessel passage) in between
\item Structure of medulla: 6-12 pyramid-shaped regions with base towards cortex and apex towards hilum \(\rightarrow\) medullary pyramid + covering cortical region = renal lobe
\item Renal papilla: sieve-like region (w/ small openings of the ducts of Bellini). Papilla surface = area cribrosa
\end{itemize}
\begin{center}\includegraphics[width=0.82\textwidth,height=0.68\textheight,keepaspectratio]{\reviewnotesimagepath d6-figure-histology-urinary-renal-papilla-sieve-like.png}\par\captionof{figure}{The illustration relates to Renal papilla: sieve-like region (w/ small openings of the ducts of Bellini). Papilla surface = area cribrosa, within the topic of Histology of urinary system.}\labfig{appendix-figure-d6-histology-urinary-renal-papilla-sieve-like}\end{center}
\begin{itemize}
\item Microscopic structures of the kidney - cortical region
\begin{itemize}
\item Renal corpuscle: blood-filtering component of kidney (composed of capillary tuft + bowman’s capsule
\item Proximal, distal convoluted tubules \(\rightarrow\) loop of henlee = longitudinal striations (aka medullary rays)
\item Renal lobules are symmetrical around medullary rays, and are aggregated around them
\item Nephrons around a lobule drain to one collecting duct
\end{itemize}
\end{itemize}
\begin{center}\includegraphics[width=0.82\textwidth,height=0.68\textheight,keepaspectratio]{\reviewnotesimagepath d6-figure-histology-urinary-nephrons-around-lobule-drain.png}\par\captionof{figure}{The illustration relates to Nephrons around a lobule drain to one collecting duct, within the topic of Histology of urinary system.}\labfig{appendix-figure-d6-histology-urinary-nephrons-around-lobule-drain}\end{center}
\begin{itemize}
\item Vasculature of the kidney
\begin{itemize}
\item 1,200 mL/min blood flow from renal artery, supplying nephrons
\item Interlobar\emph{ }artery \(\rightarrow\) arcuate arteries \(\rightarrow\) interlobular arteries \(\rightarrow\) afferent arterioles \(\rightarrow\) glomerulus \(\rightarrow\) efferent arterioles \(\rightarrow\) peritubular capillary, vasa recta \(\rightarrow\) venule \(\rightarrow\) interlobular vein \(\rightarrow\) arcuate vein \(\rightarrow\) interlobar vein
\item Vasa recta: supply loop of henle, collecting duct of juxtamedullary nephron w/ blood (for countercurrent exchange). Either arterial or venous
\end{itemize}
\item Microscopic structure of the nephron
\end{itemize}
\begin{center}\includegraphics[width=0.82\textwidth,height=0.68\textheight,keepaspectratio]{\reviewnotesimagepath d6-figure-histology-urinary-microscopic-structure-nephron.png}\par\captionof{figure}{The illustration relates to Microscopic structure of the nephron, within the topic of Histology of urinary system.}\labfig{appendix-figure-d6-histology-urinary-microscopic-structure-nephron}\end{center}
\begin{itemize}
\item Functional unit of the kidney. Contains: renal corpuscle/glomerulus (bowman’s capsule + capillary tuft), proximal + distal convoluted tubules, loop of henle
\item Nephrons are divided into 2 types depending on location of renal corpuscles: juxtamedullary nephrons (at the cortex near the medulla) and cortical nephrons (at the top near the capsule)
\item Cortical nephrons: 85\% of all nephrons, shorter loop of henlee, almost entirely within cortex
\item Juxtamedullary nephrons: nephrons with longer loop of henle, renal corpuscles located near cortex–medua junction
\item Bowman’s capsule
\begin{itemize}
\item Two layers, ‘double layered cup like’
\item Formed by 2 layers of epithelium w/ space between them (Bowman’s space) containing plasma ultrafiltrate
\item Parietal wall (outer wall) made of simple squamous epithelium continuous w/ simple cuboidal of proximal convoluted tubule
\item Visceral wall: podocyte epithelial cell
\item Capsule has 2 poles: vascular pole (for afferent, efferent arteriole) and urinary pole (proximal convoluted tubule)
\item Inside of bowman’s capsule = glomerulus: glomerular capillaries, mesangium (mesangial cells + matrix), podocytes (visceral layer of bowman’s capsule)
\end{itemize}
\item Mesangium, mesangial cells
\begin{itemize}
\item Mesangial cells = specialized pericytes (smooth muscle, macrophage characteristics)
\item Regulate glomerular filtration rate
\item Interglomerular mesangial cells = structural between glomerular capillaries
\item Also in the form of extraglomerular mesangial cells \(\rightarrow\) transmit information between macula densa and afferent arteriole (gap junctions)
\item Key properties of mesangial cells: contractile, phagocytotic, capable of proliferation
\item Mesangial cell synthesis: matrix, collagen + secretion of bioactive substances: prostaglandins, endothelins (endothelins \(\rightarrow\) aff+eff glomerular arteriole constriction)
\item Mesangial cell = site for immunoglobulin, complement molecule entry (into mesangial matrix) \(\rightarrow\) entry results in cytokine production by mesangial cells \(\rightarrow\) immune response, glomerular occlusion
\item Mesangial cell actions: mechanical support for glomerular capillaries, control turnover of glomerular basal lamina (by phagocytic activity), regulate blood flow (contractile ability), prostaglandin/endothelin secretion, angiotensin II responsive
\end{itemize}
\item Pathologies of the mesangium
\begin{itemize}
\item Anti glomerular basement membrane antibodies
\begin{itemize}
\item NC1 domain of type IV collagen on basal lamina targeted by immunoglobulins
\item Whole length of lamina affected \(\rightarrow\) granular/linear pattern of immunofluorescence
\item Anti-GBM antibody \(\rightarrow\) anti-GBM nephritis \(\rightarrow\) glomerular damage evolving into renal failure
\end{itemize}
\item Antigen-immunoglobulin complex trapping in mesangium \(\rightarrow\) immunoglobulin-compliment interaction \(\rightarrow\) mesangial damage (mesangiolysis)
\item Circulating immunoglobulin-antigen complexes (eg. SLE) deposit in glomerular basal lamina (between endothelium, GBM \(\rightarrow\) “subendothelial deposits” / in mesangium \(\rightarrow\) granular deposits). Causes hematuria, proteinuria
\item Congenital nephrotic syndrome
\begin{itemize}
\item Antibody to podocyte \(\rightarrow\) pedicel detachment
\item Failure to form foot processes (integrin mutation) \(\rightarrow\) flattened podocytes dissociated from basal lamina
\item Both cause slit diaphragm malfunction \(\rightarrow\) massive proteinuria in utero, nephritic syndrome at birth \(\rightarrow\) hypoalbuminemia, abdominal distension, hyperlipidemia, edema
\end{itemize}
\item Alport’s syndrome: X-linked recessive mutation of COL4A5 (male > female) \(\rightarrow\) irregular thinning/thickening/splitting of glomerular basal lamina \(\rightarrow\) loss of kidney function, progressive glomerulonephritis \(\rightarrow\) progressive renal failure; often associated with hearing loss, lens capsule defect
\item Goodpasture syndrome: anti-GBM(V) antibodies \(\rightarrow\) binding to basal lamina, glomerular lamina \(\rightarrow\) progressive glomerulonephritis, pulmonary hemorrhage
\end{itemize}
\item Structure of glomerulus
\begin{itemize}
\item Filters fluid across capillary walls, retaining RBCs
\item Network of capillary endothelium (fenestrated) over glomerular basement membrane (w/ endothelium, podocytes) and distal (foot) processes \(\rightarrow\) interdigitating network
\end{itemize}
\item Podocyte
\begin{itemize}
\item Epithelial cells with long body and branching (foot) processes.Encircles glomerular capillary
\item Produces glomerular endothelial growth factor \(\rightarrow\) endothelium development and fenestration maintenance
\item Underneath interdigitating (interlaced) foot processes = fenestrated endothelial cell, basal lamina (endothelium cell + podocyte)
\item Act as glomerular filtration barrier
\end{itemize}
\item “Glomerular filtration barrier”
\begin{itemize}
\item Fenestration of glomerular capillary
\item Basal lamina of endothelial cell, podocyte
\item Interdigitation of podocyte foot processes \(\rightarrow\) forms filtration slits between pedicels w/ slit diaphragm
\item Overall \(\rightarrow\) functions to filter blood and excrete fluid/waste while retaining blood cells, serum proteins
\item Structure: glomerular capillary (lumen) \(\rightarrow\) endothelial cell w/ amniotic coat (heparin sulfate) \(\rightarrow\) basal lamina \(\rightarrow\) podocytes, filtration slit diaphragm
\item Basal lamina - type IV collagen, laminin, fibronectin, glycosaminoglycan, heparin sulfate
\item Endothelial cells w/ amniotic coat (heparin sulfate, to slow down filtration of anionic protein) - fenestrated and permeable (to water, sodium, urea, glucose, small proteins)
\item Filtration slit diaphragm - links adjacent pedicels w/ nephrin (slit diaphragm protein) and F-actin (in pedicel). Inside cell connects nephrin-F-actin by podocin, ZO-1, CD2AP proteins
\end{itemize}
\end{itemize}
\begin{center}\includegraphics[width=0.82\textwidth,height=0.68\textheight,keepaspectratio]{\reviewnotesimagepath d6-figure-histology-urinary-filtration-slit-diaphragm-links.png}\par\captionof{figure}{The illustration relates to Filtration slit diaphragm - links adjacent pedicels w/ nephrin (slit diaphragm protein) and F-actin (in pedicel). Inside cell connects nephrin-F-actin by podoci, within the topic of Histology of urinary system.}\labfig{appendix-figure-d6-histology-urinary-filtration-slit-diaphragm-links}\end{center}
\begin{itemize}
\item Juxtaglomerular apparatus
\begin{itemize}
\item Endocrine structure w/ macula densa, extraglomerular mesangial cells, juxtaglomerular cells
\item Macula densa = cuboidal epithelium at initial DCT. Adjacent to afferent arteriole and mesangial cells
\item JG apparatus function: blood pressure, electrolyte, body fluid regulation
\item Macula densa cells- sense NaCl level changes \(\rightarrow\) autoregulatory to increase/decrease ion, water resorption \(\rightarrow\) alter blood volume, return BP to normal
\item JG cells: smooth muscle cells in wall of afferent arterioles (and some efferent) \(\rightarrow\) synthesize renin, RAAS regulation, renal autoregulation (bloodflow, filtration)
\end{itemize}
\item Proximal convoluted tubule
\begin{itemize}
\item Adjacent to renal corpuscles
\item Cuboidal epithelium w/ tight junctions
\item Apical domain = brush border + vesicles
\item Apical tubulovesicles, lysosomes \(\rightarrow\) endocytosis, breakdown of small protein to amino acids
\item Basolateral domain = plasma membrane interdigitations
\item Resorb 70\% of filtered water by glucose-NaCl osmotic gradient. Also resorbs glucose, ions bicarbonate, peptides
\end{itemize}
\item Loop of henle
\begin{itemize}
\item Between PCT, DCT
\item Resorbs 15\% of filtered water, 25\% of filtered NaCl, Ca, HCO
\item Descending, ascending limbs w/ thick (low cuboidal), thin (simple squamous) segments
\item Thin segment length depends on type (juxtamedullary > cortical nephrons loop length)
\end{itemize}
\item Distal convoluted tubule
\begin{itemize}
\item Between loop of henle and collecting duct
\item Simple cuboidal epithelium
\item Many mitochondria, basal infoldings, lateral membrane interdigitations
\item Distal DCT, collecting ducts permeable to water when ADH/vasopressin present
\item DCT-afferent arteriole of renal corpuscle = macula densa
\end{itemize}
\item PCT and DCT identification
\begin{itemize}
\item DCT has shorter cuboidal cells than PCT
\item Both adjacent renal corpuscles, abundant basal mitochondria
\item PCT = abundant lysosomes (dark in light microscope)
\item The apical domain of PCT has a brush border and vesicles. DCT has few of either
\end{itemize}
\item Collecting duct / tubule
\begin{itemize}
\item Series of tubules, ducts connecting nephrons to ureter
\item Cuboidal epithelium w/ principal (light) cell and intercalated (dark) cell
\item Reabsorption, excretion processes \(\leftarrow\) regulated by aldosterone and antidiuretic hormone
\item Aldosterone = Na absorption at collecting tubule \(\rightarrow\) water retention (over 24\%)
\item Principal cells
\begin{itemize}
\item Reabsorb Na+ (4-5\% of total resorption), water (5\%) + secrete K+
\item Apical nonmotile cilia, basolateral domain w/ moderate infolding + mitochondria
\end{itemize}
\item Intercalated cells
\begin{itemize}
\item Secrete H+ or bicarbonate, resorb K+ \(\rightarrow\) regulate acid-base balance
\item Possess apical microvilli, abundant mitochondria
\end{itemize}
\end{itemize}
\item Processes involved in formation of urine
\end{itemize}
1. Filtration blood to lumen

2. Reabsorption: lumen to blood

3. Secretion: blood to lumen

4. Excretion: lumen to external environment

\begin{itemize}
\item Ureter - long, straight muscle–walled tube
\end{itemize}
\begin{center}\includegraphics[width=0.82\textwidth,height=0.68\textheight,keepaspectratio]{\reviewnotesimagepath d6-figure-histology-urinary-ureter-long-straight-muscle.png}\par\captionof{figure}{The illustration relates to Ureter - long, straight muscle–walled tube, within the topic of Histology of urinary system.}\labfig{appendix-figure-d6-histology-urinary-ureter-long-straight-muscle}\end{center}
\begin{itemize}
\item Inner layer = mucosa: epithelium (urothelium/transitional epithelium) \(\rightarrow\) thick fibroelastic lamina propria
\item (Upper ⅔) middle layer: smooth muscle. Layers - inner longitudinal \(\rightarrow\) outer circular
\item (lower ⅓) middle layer: inner longitudinal, middle circular \(\rightarrow\) outer longitudinal
\item Outer layer = adventitia: fibroelastic connective tissue, blood vessels, lymphatics, nerves
\item Urinary bladder
\begin{itemize}
\item Distensible, collapsable organ w/ capacity of 700-800 mL
\item Located behind pubic symphysis
\item Lining = urothelium \(\rightarrow\) fibroelastic CNT in lamina propria
\item Consisted of 3 interlaced layers of smooth muscle (‘detrusor muscle’)
\begin{itemize}
\item Forms a sphincter (internal sphincter) in the neck region
\item Inner longitudinal \(\rightarrow\) middle circular \(\rightarrow\) outer longitudinal
\end{itemize}
\item Covered by adventitia (loose CNT), serosal layer in superior surface (visceral peritoneum)
\end{itemize}
\item Urethra
\begin{itemize}
\item Extends from base of bladder to outside
\item 4 cm in female, 14 cm in male
\item Transition of epithelium: (proximal) transitional epithelium \(\rightarrow\) stratified columnar epithelium \(\rightarrow\) pseudostratified columnar epithelium \(\rightarrow\) stratified squamous epithelium (external orifice / distal)
\item Female urethra: proximal ⅔ lined by transitional epithelium, distal ⅓ lined by stratified squamous epithelium
\item Male urethra: prosthetic urethra, transitional epithelium (prostate gland area, joined by 2 ejaculatory ducts) \(\rightarrow\) membranous urethra \(\rightarrow\) penile urethra
\end{itemize}
\item Interstitial cells
\begin{itemize}
\item Fibroblasts + dendritic cells
\item Extend from loop of henle to vasa recta (‘rungs of a ladder’)
\item Fibroblasts - act as the skeleton of the kidney and produce erythropoietin in the renal cortex. Prominent F-actin cytoskeleton w/ abundant rough endoplasmic reticulum
\item Dendritic cells - mononuclear phagocyte system
\end{itemize}
\end{itemize}
\begin{center}\includegraphics[width=0.82\textwidth,height=0.68\textheight,keepaspectratio]{\reviewnotesimagepath d6-figure-histology-urinary-dendritic-cells-mononuclear-phagocyte.png}\par\captionof{figure}{The illustration relates to Dendritic cells - mononuclear phagocyte system, within the topic of Histology of urinary system.}\labfig{appendix-figure-d6-histology-urinary-dendritic-cells-mononuclear-phagocyte}\end{center}

\subsection{Urine formation}
\begin{itemize}
\item LOs:
\begin{itemize}
\item Describe functional organization of nephron in urine formation
\item Explain process of urine formation (glomerular filtration, tubular reabsorption, tubular secretion)
\item Identify, explain factors influencing/regulating urine formation
\item Describe hormonal regulation of urine concentration, volume
\end{itemize}
\item Overview of the kidney
\begin{itemize}
\item Paired, bean shaped retroperitoneal organs
\item Located at the T12-L3 level of the posterior abdominal wall
\item Behind peritoneum, below diaphragm (with ribs 11, 12 behind)
\item Filter blood \(\rightarrow\) produce urine \(\rightarrow\) transport by ureters to bladder \(\rightarrow\) bladder stores urine \(\rightarrow\) urethra conveys urine out of the body
\item 20-25\% of cardiac output (1.0–1.25 L/min). Has autoregulation mechanisms
\end{itemize}
\item Functions of the kidney
\begin{itemize}
\item Excretion of waste (metabolic, foreign substances) \(\rightarrow\) urine formation
\item Regulation of:
\begin{itemize}
\item Body fluid volume, osmolarity
\item Electrolyte balance (Na+, K+, Ca2+)
\item Arterial blood pressure (RAAS)
\item Acid-base balance (H+, HCO3- secretion)
\item Secretion (renin, erythropoietin), metabolism, excretion of hormones
\item Gluconeogenesis during fast
\end{itemize}
\end{itemize}
\item The nephron
\begin{itemize}
\item “Functional unit” of the kidney \(\rightarrow\) 1-1.5 million in each
\item Filtration of blood, reabsorbing substances, forming urine
\item Renal corpuscle = glomerulus + bowman’s capsule in cortex
\item Renal tubule = proximal convoluted tubule \(\rightarrow\) loop of henle \(\rightarrow\) distal convoluted tubule
\item Collecting duct = receives filtrate, carries from medulla to renal papilla
\item Two types - cortical nephron and juxtamedullary nephron
\item Cortical nephron
\begin{itemize}
\item At outer cortex w/ short loop of henle (barely into medulla)
\item 80–85\% of all nephrons
\item Surrounded by peritubular capillaries
\item Majority of filtration, resorption. Minor role in urine concentration
\end{itemize}
\item Juxtamedullary nephron
\begin{itemize}
\item At renal corpuscle near corticomedullary junction
\item Larger glomerulus
\item Long loop of henle into medulla w/ long thin ascending limb
\item Surrounded by vasa recta
\item Establishes osmotic gradient, water/salt balance
\end{itemize}
\end{itemize}
\item Juctaglomerular apparatus
\begin{itemize}
\item Composed of: juxtaglomerular/granular, macula densa and extraglomerular mesangial cells
\item Secretes renin \(\rightarrow\) control of blood pressure by fluid balance, vasoconstriction
\end{itemize}
\item Renal epithelium \(\rightarrow\) specialized for each location
\begin{itemize}
\item Proximal tubule: cuboidal epithelium w/ brush border and mitochondria dense. Responsible for bulk reabsorption
\item Loop of henle: water permeable in descending thin limb, active Na+ transport and water impermeable thick limb
\item Distal tubule: electrolyte regulation. Has macula densa
\item Collecting duct: ADH, aldosterone regulated water reabsorption, acid–base balance
\end{itemize}
\item Flow of blood in kidney: abdominal aorta \(\rightarrow\) renal artery \(\rightarrow\) interlobar arteries \(\rightarrow\) arcuate arteries \(\rightarrow\) interlobular arteries \(\rightarrow\) afferent arterioles \(\rightarrow\) glomerular capillaries \(\rightarrow\) efferent arterioles \(\rightarrow\) peritubular capillaries, vasa recta \(\rightarrow\) venules \(\rightarrow\) interlobular veins \(\rightarrow\) arcuate veins \(\rightarrow\) interlobar veins \(\rightarrow\) renal vein \(\rightarrow\) inferior vena cava
\item Renal nerve supply: efferent renal sympathetic nerves from rostral ventrolateral medulla, paraventricular nucleus
\item Renal afferents
\begin{itemize}
\item Renal pelvis wall, around glomeruli, in tubulointerstitium
\item Express CGRP (calcitonin–gene related peptide)
\item Responses: chemosensitive, mechanosensitive, nociceptive (pain)
\item Activation of afferents \(\rightarrow\) increased sympathetic drive \(\rightarrow\) renal sympathetic activation
\item Excessive activity \(\rightarrow\) hypertension, heart failure, diabetic kidney disease
\end{itemize}
\item Noradrenaline (norepinephrine) \(\rightarrow\) \(\alpha\) and β adrenoreceptors causes:
\begin{itemize}
\item Increased sodium reabsorption
\item \(\alpha\)1B-adrenoreceptor \(\rightarrow\) increased Na-H exchanger 3 in proximal tubule
\item \(\alpha\)1-adrenoreceptor \(\rightarrow\) increased Na-K ATPase activity in thick ascending limb
\item Na-Cl cotransporter expression increased in distal convoluted tubule
\item Epithelium sodium channel increased activity in collecting duct
\item \(\alpha\)1-adrenoreceptor induced vasoconstriction
\item Renin release by β1-adrenoreceptor \(\rightarrow\) β1–adrenergic receptor on JG cells \(\rightarrow\) RAAS activation \(\rightarrow\) increased blood pressure
\end{itemize}
\item Urine in renal pelvis \(\rightarrow\) peristaltic contraction (amplify by parasympathetic) generated by pacemaker cell in calices \(\rightarrow\) bladder \(\rightarrow\) micturition (emptying of bladder)
\item Micturition reflex
\begin{itemize}
\item As urine volume increases, pressure remains low
\item As volume approaches max capacity, pressure increases \(\rightarrow\) urge to urinate
\item “Sigmoid shaped graph”
\item Process of urination: bladder filling \(\rightarrow\) activation of stretch receptor \(\rightarrow\) impulses sent to sacral spinal cord \(\rightarrow\) parasympathetic efferents \(\rightarrow\) detrusor contraction, relaxed sphincter \(\rightarrow\) external urethral sphincter relaxation \(\rightarrow\) urination occurs
\item Abnormalities: destruction of sensory fibers, lack of brain inhibitory signalling (spinal/brainstem damage)
\end{itemize}
\item Urine formation
\begin{itemize}
\item Glomerular filtration: blood filtered in glomerulus, enters bowman’s capsule
\item Tubular reabsorption to blood: glucose, water, ions
\item Tubular secretion: waste products, H+, K+, drugs
\item Urine excretion rate = filtration rate - reabsorption rate + secretion rate
\end{itemize}
\item Filtration process in kidney
\begin{itemize}
\item Nonselective, driven by blood pressure. RBCs and large proteins remain in lumen
\item Forces opposing filtration: blood colloid osmotic pressure (30 mmHg), capsular hydrostatic pressure (15 mmHg)
\item Net filtration = 10 mmHg into bowman’s capsule
\item Filtration barrier prevents protein loss, allows ultrafiltrate to be formed
\item Filtrate passes through: fenestrated endothelium \(\rightarrow\) glomerular basement membrane (w/ 3 layer collagen network of 30-60 nm pores) \(\rightarrow\) podocyte layer (foot processes and glycocalyx filter by charge and size), slit diaphragm
\item Freely permeable to Mw < 10,000, slender + flexible. Cation > neutral > anion
\item Measured by glomerular filtration rate (GFR) - volume of filtrate formed in kidneys per minute \(\rightarrow\) indicator of kidney function, disease stage indicator
\item GFR depends on ultrafiltration coefficient, renal bloodflow, blood pressure, filtration surface area, membrane permeability and plasma oncotic pressure
\item Ultrafiltration coefficient = surface area and fluid permeability. Decreased in CKD (decreased surface area)
\item Factors controlling/influencing GFR
\begin{itemize}
\item Glomerular hydrostatic pressure
\item Glomerular capillary colloid osmotic pressure
\item Sympathetic nervous system, hormones, autacoids (local vasoactives), intrinsic kidney feedback mechanisms
\item Reduced by sympathetic activity, vasoconstrictors. Vasodilators increase
\item Angiotensin II efferent arteriole constriction: GFR maintained at low perfusion
\item Autoregulation, intrinsic controls
\item tubuloglomerular feedback: increased NaCl delivery to distal tubule measured by macula densa \(\rightarrow\) afferent arteriole constrict
\item Myogenic mechanism: pressure \(\rightarrow\) stretch \(\rightarrow\) vasoconstriction \(\rightarrow\) muscle tension regulation
\item Extrinsic controls \(\rightarrow\) angiotensin II decreases GFR, prostaglandins increase GFR by selective dilation/constriction of vessels
\item NSAIDS decrease prostaglandin \(\rightarrow\) GFR decreased
\item High protein intake \(\rightarrow\) increased GFR, kidney growth if sustained (because amino acids resorbed with sodium in proximal tubule \(\rightarrow\) macula densa sees less sodium \(\rightarrow\) tubuloglomerular feedback \(\rightarrow\) afferent arteriole dilation)
\item Increased blood glucose \(\rightarrow\) increased filtered load of glucose \(\rightarrow\) Na-Glucose cotransporter \(\rightarrow\) macula densa, tubuloglomerular feedback
\item Other factors: age, hydration, body surface area, renal blood flow restriction, intrarenal tissue pressure change (urethral obstruction, renal edema, perinephritis), reduced surface area
\end{itemize}
\end{itemize}
\end{itemize}
\begin{center}\includegraphics[width=0.82\textwidth,height=0.68\textheight,keepaspectratio]{\reviewnotesimagepath d6-figure-urine-formation-other-factors-age-hydration.png}\par\captionof{figure}{The illustration relates to Other factors: age, hydration, body surface area, renal blood flow restriction, intrarenal tissue pressure change (urethral obstruction, renal edema, perinephri, within the topic of Urine formation.}\labfig{appendix-figure-d6-urine-formation-other-factors-age-hydration}\end{center}
\begin{center}\includegraphics[width=0.82\textwidth,height=0.68\textheight,keepaspectratio]{\reviewnotesimagepath d6-figure-urine-formation-urine-formation.png}\par\captionof{figure}{The illustration relates to Urine formation, within the topic of Urine formation.}\labfig{appendix-figure-d6-urine-formation-urine-formation}\end{center}
\begin{itemize}
\item Intrinsic regulation - maintain consistent GFR. Extrinsic regulation - maintain blood volume, pressure
\item Tubular reabsorption
\begin{itemize}
\item Active and passive transport of: glucose, amino acids, electrolytes, water
\item Non–reabsorbed substances: urea, uric acid, creatinine
\item Factors affecting absorption - charge, size, transtubular potential
\item Parathormone, calcitonin \(\rightarrow\) PTH increases Ca2+ reabsorption, reduces phosphate reabsorption, calcitonin does inverse
\item ADH \(\rightarrow\) affects permeability of distal convoluted tubules, collecting tubules to water. Increases water reabsorption in DCT, collecting tubules
\item Aldosterone \(\rightarrow\) Increases Na+/water absorption, K+ excretion (RAAS regulated)
\item Atrial natriuretic peptide \(\rightarrow\) decreases water Na+ and water reabsorption (decreases blood volume)
\item Transcellular transport - through cells, either active or passive (solutes)
\item Paracellular transport - between cells, by diffusion (solutes)
\item Water transport - osmosis following solute gradient. Tubular lumen \(\rightarrow\) tubular epithelial cells \(\rightarrow\) renal interstitium \(\rightarrow\) peritubular capillaries
\item Bulk flow - ultrafiltration into peritubular capillaries (‘solvent drag’)
\item Chemical/electrical gradient transport - water, chloride, urea
\item Active transport - Na+, K+, PO43-, amino acids, glucose, creatinine, sulfate, uric acid, ketone bodies
\item Protein reabsorption - pinocytosis, receptor mediated uptake (with carriers for amino acids)
\end{itemize}
\end{itemize}
\begin{center}\includegraphics[width=0.82\textwidth,height=0.68\textheight,keepaspectratio]{\reviewnotesimagepath d6-figure-urine-formation-protein-reabsorption-pinocytosis-receptor.png}\par\captionof{figure}{The illustration relates to Protein reabsorption - pinocytosis, receptor mediated uptake (with carriers for amino acids), within the topic of Urine formation.}\labfig{appendix-figure-d6-urine-formation-protein-reabsorption-pinocytosis-receptor}\end{center}
\begin{center}\includegraphics[width=0.82\textwidth,height=0.68\textheight,keepaspectratio]{\reviewnotesimagepath d6-figure-urine-formation-urine-formation-02.png}\par\captionof{figure}{The illustration relates to Urine formation, within the topic of Urine formation.}\labfig{appendix-figure-d6-urine-formation-urine-formation-02}\end{center}
\begin{center}\includegraphics[width=0.82\textwidth,height=0.68\textheight,keepaspectratio]{\reviewnotesimagepath d6-figure-urine-formation-urine-formation-03.png}\par\captionof{figure}{The illustration relates to Urine formation, within the topic of Urine formation.}\labfig{appendix-figure-d6-urine-formation-urine-formation-03}\end{center}
\begin{itemize}
\item Glucose reabsorption - Na+K+ ATPase electrochemical gradient \(\rightarrow\) SGLT2 cotransport of Na-Glucose into cell \(\rightarrow\) GLUT facilitates diffusion out of cell
\item Glucose reabsorption capacity is limited \(\rightarrow\) beyond threshold, glucose is excreted
\item Sodium reabsorption \(\rightarrow\) diffusion from lumen to epithelial cells \(\rightarrow\) Na+-K+ ATPase gradient on basolateral membrane \(\rightarrow\) interstitial fluid to capillaries by ultrafiltration (bulk flow) w/ hydrostatic, colloid osmotic pressure gradients
\item Urea: 70\% excreted, absorbed less at high urine flow (low GFR \(\rightarrow\) low flow \(\rightarrow\) urea retained \(\rightarrow\) blood urea nitrogen increased)
\end{itemize}
\begin{center}\includegraphics[width=0.82\textwidth,height=0.68\textheight,keepaspectratio]{\reviewnotesimagepath d6-figure-urine-formation-urea-70-excreted-absorbed.png}\par\captionof{figure}{The illustration relates to Urea: 70\% excreted, absorbed less at high urine flow (low GFR \(\rightarrow\) low flow \(\rightarrow\) urea retained \(\rightarrow\) blood urea nitrogen increased), within the topic of Urine formation.}\labfig{appendix-figure-d6-urine-formation-urea-70-excreted-absorbed}\end{center}
\begin{center}\includegraphics[width=0.82\textwidth,height=0.68\textheight,keepaspectratio]{\reviewnotesimagepath d6-figure-urine-formation-urine-formation-04.png}\par\captionof{figure}{The illustration relates to Urine formation, within the topic of Urine formation.}\labfig{appendix-figure-d6-urine-formation-urine-formation-04}\end{center}
\begin{itemize}
\item Na+ reabsorbed \(\rightarrow\) negative charge in lumen \(\rightarrow\) Cl follows passively
\item Water reabsorption concentrates Cl, urea in lumen \(\rightarrow\) Cl, urea reabsorption occurs passively
\end{itemize}
\begin{center}\small\begin{tabularx}{\linewidth}{>{\raggedright\arraybackslash}p{1.55cm}>{\raggedright\arraybackslash}X>{\raggedright\arraybackslash}X>{\raggedright\arraybackslash}X>{\raggedright\arraybackslash}X}
\toprule
Segment & Main Function & Key Transporter & Hormone & Diuretic \\
\midrule
PCT & Isosmotic reabsorption & Na/H exchanger (NHE3) & PTH decreases Pi, Ang II increases Na+ & Acetazolamide, SGLT2 inhibitors \\
Thick Asc Limb & NaCl reabs (no H2O) & Na-K-2Cl (NKCC2) & ADH increases NKCC & Loop diuretics (Furosemide) \\
Early DCT & NaCl reabs (no H2O) & Na-Cl (NCC) & PTH increases Ca2+ & Thiazides \\
Principal (Late DCT/CD) & Na+ reabs, K+ sec, H2O & ENaC, ROMK & Aldo increases Na+/K+, ADH increases H2O & K-sparing diuretics (Amiloride) \\
alpha-Intercalated & H+ sec, K+ reabsorption & H+ ATPase & Aldo increases H+ & K-sparing diuretics (Spironolactone) \\
\bottomrule
\end{tabularx}\captionof{table}{Comparison of transport functions, regulatory hormones, and diuretics across nephron segments.}\end{center}
\begin{itemize}
\item Tubular secretion
\begin{itemize}
\item Moves substances from peritubular capillaries to DCT for excretion
\item Substances secreted: H+, K+, creatinine, drugs, weak acids/bases (eg. ammonia, salicylic acid, phenobarbital, penicilin) \(\rightarrow\) acid–base homeostasis, electrolyte balance, unfiltered waste removal \(\rightarrow\) fine tunes volume of blood, blood pressure (by ion balance)
\item Types of secretion
\begin{itemize}
\item Active with transport maximum - transport limited by carrier capacity (secretion increases until Tm is reached). Can be used to measure renal plasma flow
\item Active limited to gradient and time - requires energy. H+ secretion (mostly proximal tubule, some in collecting tubule)
\item Passive secretion - driven by concentration/electrical gradient (may depend on energy for gradient creation)
\end{itemize}
\end{itemize}
\end{itemize}
\begin{center}\includegraphics[width=0.82\textwidth,height=0.68\textheight,keepaspectratio]{\reviewnotesimagepath d6-figure-urine-formation-passive-secretion-driven-concentration.png}\par\captionof{figure}{The illustration relates to Passive secretion - driven by concentration/electrical gradient (may depend on energy for gradient creation), within the topic of Urine formation.}\labfig{appendix-figure-d6-urine-formation-passive-secretion-driven-concentration}\end{center}
\begin{center}\includegraphics[width=0.82\textwidth,height=0.68\textheight,keepaspectratio]{\reviewnotesimagepath d6-figure-urine-formation-urine-formation-05.png}\par\captionof{figure}{The illustration relates to Urine formation, within the topic of Urine formation.}\labfig{appendix-figure-d6-urine-formation-urine-formation-05}\end{center}
\begin{center}\includegraphics[width=0.82\textwidth,height=0.68\textheight,keepaspectratio]{\reviewnotesimagepath d6-figure-urine-formation-urine-formation-06.png}\par\captionof{figure}{The illustration relates to Urine formation, within the topic of Urine formation.}\labfig{appendix-figure-d6-urine-formation-urine-formation-06}\end{center}
\begin{itemize}
\item Measurement of renal function
\begin{itemize}
\item Measures the rate of toxin filtration in glomerulus
\item Clearance - volume of plasma with substance removed by kidneys per unit of time (mL/min). Applies for non-reabsorbed substances (e.g. inulin) where clearance=GFR
\item Clearance = amount excreted per minute / plasma concentration
\item Clearance > GFR indicates secretion, clearance < GFR indicates absorption
\item Creatinine (can be used to calculate GFR): produced by skeletal muscle metabolism \(\rightarrow\) constant release to bloodstream \(\rightarrow\) freely filtered at glomerulus w/o reabsorption (small amount of secretion, slight overestimation)
\end{itemize}
\end{itemize}
\begin{center}\includegraphics[width=0.82\textwidth,height=0.68\textheight,keepaspectratio]{\reviewnotesimagepath d6-figure-urine-formation-creatinine-be-calculate-gfr.png}\par\captionof{figure}{The illustration relates to Creatinine (can be used to calculate GFR): produced by skeletal muscle metabolism \(\rightarrow\) constant release to bloodstream \(\rightarrow\) freely filtered at glomerulus w/o reabsorp, within the topic of Urine formation.}\labfig{appendix-figure-d6-urine-formation-creatinine-be-calculate-gfr}\end{center}
\begin{itemize}
\item Factors influencing urine concentration
\begin{itemize}
\item Gradient generation (Na reabsorption in countercurrent multiplication in thick ascending limb, urea reabsorption in medullary collecting duct)
\item Gradient utilization (ADH-dependant water reabsorption) in collecting duct
\item Gradient maintenance (countercurrent exchange in vasa recta)
\item Failures of urine concentration
\begin{itemize}
\item Poor solute delivery to loop of henle (CKD)
\item Impaired ascending limb function (loop diuretic)
\item Lack of ADH (central DI)
\item ADH insensitivity in kidney (nephrogenic DI)
\end{itemize}
\item Osmotic gradient of medullary interstitium
\begin{itemize}
\item Length of loops of henle (longer = greater ability to concentrate)
\item Rate of NaCl reabsorption in thick ascending limb (higher delivery, better concentrating ability)
\end{itemize}
\item Dietary protein content (more protein = more urea = increased concentrating ability)
\item Medullary blood flow - low flow leads to solvent accumulation \(\rightarrow\) higher concentration of urine (and vice versa)
\item Osmotic permeability of collecting ducts and tubules - ADH promotes permeability \(\rightarrow\) increased reabsorption \(\rightarrow\) concentrated urine (and vice versa)
\item Luminal flow in loop of henle, collecting ducts - higher flow reduces countercurrent efficiency (less time for water–urea equilibrium) \(\rightarrow\) less water absorbed
\item Pathophysiologies: central diabetes insipidus: low plasma ADH = less water absorbed, nephrogenic diabetes insipidus = nonresponsive kidney, dilute urine
\end{itemize}
\end{itemize}

\subsection{Fluid and Electrolyte Balance}
\begin{itemize}
\item LOs:
\begin{itemize}
\item Describe the role of kidneys in fluid, electrolyte balance
\item Identify how kidneys regulate water and key electrolytes in homeostasis
\item Explain renal mechanisms and hormonal control of fluid, electrolyte balance
\item Applied kidney function knowledge to understand and manage disturbances in fluid/electrolytes
\end{itemize}
\item Fluid, electrolyte balance
\begin{itemize}
\item Body water composition: water and solutes (proteins, electrolytes, metabolites)
\item Electrolytes = charged particles
\begin{itemize}
\item Cations (positive charge) - Na+, K+, Ca2+, H+
\item Anions (negative charge - Cl-, HCO3-, PO43-
\end{itemize}
\item Non-electrolytes, uncharged particles: proteins, urea, glucose, O2, CO2
\item Divided into intracellular and extracellular fluid (ICF, ECF)
\item Regulated by kidneys (volume, composition regulation): water, sodium regulation \(\rightarrow\) fluid balance \(\rightarrow\) blood pressure and volume
\item ECF volume = total Na+ content (water follows Na+)
\item Movement of fluid: plasma \(\rightarrow\) interstitial fluid \(\rightarrow\) lymphatic capillaries
\item Movement of water: by cell membranes, bloodstream and lymphatic system by capillary walls
\item Mechanisms: diffusion, osmosis, active transport
\item Fluid distribution determined by solute concentration (water moves towards higher concentration to balance osmolarity)
\end{itemize}
\item Mechanisms regulating osmolarity
\begin{itemize}
\item Effector = ADH (sympathetic nervous system) \(\rightarrow\) aquaporin 2 expressed on late DCT/collecting duct luminal epithelium (see below)
\end{itemize}
\end{itemize}
\begin{center}\includegraphics[width=0.82\textwidth,height=0.68\textheight,keepaspectratio]{\reviewnotesimagepath d6-figure-fluid-electrolyte-balance-effector-adh-sympathetic-nervous.png}\par\captionof{figure}{The illustration relates to Effector = ADH (sympathetic nervous system) \(\rightarrow\) aquaporin 2 expressed on late DCT/collecting duct luminal epithelium (see below), within the topic of Fluid and Electrolyte Balance.}\labfig{appendix-figure-d6-fluid-electrolyte-balance-effector-adh-sympathetic-nervous}\end{center}
\begin{center}\includegraphics[width=0.82\textwidth,height=0.68\textheight,keepaspectratio]{\reviewnotesimagepath d6-figure-fluid-electrolyte-balance-fluid-electrolyte-balance.png}\par\captionof{figure}{The illustration relates to Fluid and Electrolyte Balance, within the topic of Fluid and Electrolyte Balance.}\labfig{appendix-figure-d6-fluid-electrolyte-balance-fluid-electrolyte-balance}\end{center}
\begin{center}\includegraphics[width=0.82\textwidth,height=0.68\textheight,keepaspectratio]{\reviewnotesimagepath d6-figure-fluid-electrolyte-balance-fluid-electrolyte-balance-02.png}\par\captionof{figure}{The illustration relates to Fluid and Electrolyte Balance, within the topic of Fluid and Electrolyte Balance.}\labfig{appendix-figure-d6-fluid-electrolyte-balance-fluid-electrolyte-balance-02}\end{center}
\begin{center}\includegraphics[width=0.82\textwidth,height=0.68\textheight,keepaspectratio]{\reviewnotesimagepath d6-figure-fluid-electrolyte-balance-fluid-electrolyte-balance-03.png}\par\captionof{figure}{The illustration relates to Fluid and Electrolyte Balance, within the topic of Fluid and Electrolyte Balance.}\labfig{appendix-figure-d6-fluid-electrolyte-balance-fluid-electrolyte-balance-03}\end{center}
\begin{itemize}
\item Causes of fluid imbalance
\begin{itemize}
\item Addison’s disease (aldosterone, cortisol deficiency / adrenocortical insufficiency???) \(\rightarrow\) hypoosmotic ECF volume loss (water flows from ECF \(\rightarrow\) ICF)
\item Secretory diarrhea \(\rightarrow\) isoosmotic ECF volume loss (no net ICF-ECF flow)
\item Diabetes insipidus \(\rightarrow\) water loss > solute loss \(\rightarrow\) hyperosmotic volume loss from ECF (water flows from ICF \(\rightarrow\) ECF)
\end{itemize}
\item States of imbalance (and return to euvolemia)
\begin{itemize}
\item Hypovolemia \(\rightarrow\) feedback by RAAS, ADH (thirst and compensatory actions). Results from inadequate intake/fluid loss (e.g. vomiting, diarrhea, fever, sweating, inadequate intake) \(\leftarrow\) oral rehydration, IV fluids
\item Hypervolemia \(\rightarrow\) excess fluid accumulation (renal/cardiovascular diseases) \(\rightarrow\) edema if extracellular, cerebral edema if intracellular fluid gain
\end{itemize}
\end{itemize}
\begin{center}\includegraphics[width=0.82\textwidth,height=0.68\textheight,keepaspectratio]{\reviewnotesimagepath d6-figure-fluid-electrolyte-balance-hypervolemia-excess-fluid-accumulation.png}\par\captionof{figure}{The illustration relates to Hypervolemia \(\rightarrow\) excess fluid accumulation (renal/cardiovascular diseases) \(\rightarrow\) edema if extracellular, cerebral edema if intracellular fluid gain, within the topic of Fluid and Electrolyte Balance.}\labfig{appendix-figure-d6-fluid-electrolyte-balance-hypervolemia-excess-fluid-accumulation}\end{center}
\begin{itemize}
\item Negative water balance \(\rightarrow\) water conserved \(\rightarrow\) hyperosmotic urine (and vice versa)
\item Syndrome of inappropriate antidiuretic hormone secretion (SIADH)
\end{itemize}
\begin{center}\includegraphics[width=0.82\textwidth,height=0.68\textheight,keepaspectratio]{\reviewnotesimagepath d6-figure-fluid-electrolyte-balance-syndrome-inappropriate-antidiuretic-hormone.png}\par\captionof{figure}{The illustration relates to Syndrome of inappropriate antidiuretic hormone secretion (SIADH), within the topic of Fluid and Electrolyte Balance.}\labfig{appendix-figure-d6-fluid-electrolyte-balance-syndrome-inappropriate-antidiuretic-hormone}\end{center}
\begin{itemize}
\item Electrolyte balance
\begin{itemize}
\item Na+ determines ECF volume \(\rightarrow\) effective circulating volume (tissue perfusion)
\item Increased Na+ excretion results in decreased ECF volume
\end{itemize}
\end{itemize}
\begin{center}\includegraphics[width=0.82\textwidth,height=0.68\textheight,keepaspectratio]{\reviewnotesimagepath d6-figure-fluid-electrolyte-balance-increased-na-excretion-results.png}\par\captionof{figure}{The illustration relates to Increased Na+ excretion results in decreased ECF volume, within the topic of Fluid and Electrolyte Balance.}\labfig{appendix-figure-d6-fluid-electrolyte-balance-increased-na-excretion-results}\end{center}
\begin{itemize}
\item Feedback of circulating volume
\begin{itemize}
\item Effective circulating volume is measured by low + high pressure baroreceptors \(\rightarrow\) short and long term responses
\item Short term control (seconds-minutes) - autonomic, humoral responses \(\rightarrow\) heart and vessels maintain BP
\item Long term control (hours-days) - autonomic, hormonal and renal mechanisms regulating Na+ excretion (kidney as main effector responding to ECF volume changes)
\end{itemize}
\end{itemize}
\begin{center}\includegraphics[width=0.82\textwidth,height=0.68\textheight,keepaspectratio]{\reviewnotesimagepath d6-figure-fluid-electrolyte-balance-long-term-control-hours.png}\par\captionof{figure}{The illustration relates to Long term control (hours-days) - autonomic, hormonal and renal mechanisms regulating Na+ excretion (kidney as main effector responding to ECF volume changes), within the topic of Fluid and Electrolyte Balance.}\labfig{appendix-figure-d6-fluid-electrolyte-balance-long-term-control-hours}\end{center}
\begin{itemize}
\item Mechanisms of volume control by the kidney: hemodynamic changes and Na+ transport changes in renal tubular cells (Na+ influences water movement)
\item RAAS axis: RAAS drives the conservation of water, increased water intake and vasoconstriction to compensate for hypovolemia
\end{itemize}
\begin{center}\includegraphics[width=0.82\textwidth,height=0.68\textheight,keepaspectratio]{\reviewnotesimagepath d6-figure-fluid-electrolyte-balance-raas-axis-raas-drives.png}\par\captionof{figure}{The illustration relates to RAAS axis: RAAS drives the conservation of water, increased water intake and vasoconstriction to compensate for hypovolemia, within the topic of Fluid and Electrolyte Balance.}\labfig{appendix-figure-d6-fluid-electrolyte-balance-raas-axis-raas-drives}\end{center}
\begin{itemize}
\item Angiotensin II drives Na+ reabsorption and hemodynamic changes in the kidney resulting in reduced water excretion. Aldosterone from renal cortex secretion is also stimulated
\end{itemize}
\begin{center}\includegraphics[width=0.82\textwidth,height=0.68\textheight,keepaspectratio]{\reviewnotesimagepath d6-figure-fluid-electrolyte-balance-angiotensin-ii-drives-na.png}\par\captionof{figure}{The illustration relates to Angiotensin II drives Na+ reabsorption and hemodynamic changes in the kidney resulting in reduced water excretion. Aldosterone from renal cortex secretion is al, within the topic of Fluid and Electrolyte Balance.}\labfig{appendix-figure-d6-fluid-electrolyte-balance-angiotensin-ii-drives-na}\end{center}
\begin{itemize}
\item Sodium, chloride and the kidneys
\begin{itemize}
\item Filtered load of sodium: GFR \textasciitilde{}180 L/day × [Na⁺] 142 mM \(\approx\) 25,500 mmol/day \textasciitilde{}99.6\% Na⁺ reabsorbed
\item Dietary Na⁺ intake matches excretion at \textasciitilde{}120 mmol/day
\item Cl- is resorbed with Na+ by transcellular and paracellular pathways (varying by segment)
\begin{itemize}
\item Early proximal tubule: paracellular Cl- resorption driven by electrical gradient
\item Late proximal, thick ascending limb, distal, collecting duct: transcellular transport by exchangers, cotransporters
\end{itemize}
\item Transport of Na+ and Cl- are regulated by
\begin{itemize}
\item Renal hemodynamics/glomurolotubular balance: renal blood flow/circulation, proximal tubule resorption rate
\item Effective circulating volume (from hormonal, neural controls)
\item Natriuretic humoral factors \(\rightarrow\) reduced excretion of Na+, Cl-
\end{itemize}
\end{itemize}
\end{itemize}
\begin{center}\includegraphics[width=0.82\textwidth,height=0.68\textheight,keepaspectratio]{\reviewnotesimagepath d6-figure-fluid-electrolyte-balance-natriuretic-humoral-factors-reduced.png}\par\captionof{figure}{The illustration relates to Natriuretic humoral factors \(\rightarrow\) reduced excretion of Na+, Cl-, within the topic of Fluid and Electrolyte Balance.}\labfig{appendix-figure-d6-fluid-electrolyte-balance-natriuretic-humoral-factors-reduced}\end{center}
\begin{center}\includegraphics[width=0.82\textwidth,height=0.68\textheight,keepaspectratio]{\reviewnotesimagepath d6-figure-fluid-electrolyte-balance-fluid-electrolyte-balance-04.png}\par\captionof{figure}{The illustration relates to Fluid and Electrolyte Balance, within the topic of Fluid and Electrolyte Balance.}\labfig{appendix-figure-d6-fluid-electrolyte-balance-fluid-electrolyte-balance-04}\end{center}
\begin{itemize}
\item Factors increasing Na+ reabsorption:
\begin{itemize}
\item Angiotensin II, sympathetic nervous system, ADH (by AT1 receptors in proximal tubule, thick ascending limb); stimulates NHE3 (Na+/H+ exchanger), Na+ channel stimulated \(\rightarrow\) increased Na+ resorption
\item Aldosterone \(\rightarrow\) principal cell mineralocorticoid receptors \(\rightarrow\) ENaC upregulation, Na+,K+ ATPase, K+ channels \(\rightarrow\) Na+ retained \(\leftarrow\)11β-HSD2 blocks receptor sensing of cortisol
\item Norepinephrine release \(\rightarrow\) reduced renal blood flow, GFR \(\rightarrow\) Reduced Na+ load to tubules; direct activation of \(\alpha\)-adrenergic receptors on proximal tubules; NHE3 on apical, Na-K ATPase on basolateral stimulated
\item ADH from posterior pituitary \(\rightarrow\) V2 receptor binding (Gs-cAMP pathway) \(\rightarrow\) water retention by aquaporin 2 channels; Na+ resorption by NKCC2 and K+ channel stimulation in thick ascending limb; increased number of open ENaC channels
\end{itemize}
\end{itemize}
\begin{itemize}
\item Factor decreasing Na+ resorption (increased natriuresis):
\begin{itemize}
\item Atrial natriuretic peptide (ANP from atrial myocytes in increased ECF volume) \(\rightarrow\) guanylyl cyclase activation \(\rightarrow\) cGMP signalling \(\rightarrow\) renal vasodilation, increased GFR, Na+ delivered to nephron, washout of medullary gradient
\item Aldosterone deficiency (eg. Addison’s disease) \(\rightarrow\) Na+ loss, low ECF volume
\item Aldosterone inhibition, e.g. licorice, carenoxolone
\end{itemize}
\item The phosphate balance
\begin{itemize}
\item Total plasma phosphate = 0.8-1.5 mM (2.5-4.5 mg/dL) \(\rightarrow\) 80\% as HPO42- and 20\% as H2PO4- (when at pH 7.4)
\item 7,000 mg filtered per day
\item 80\% reabsorbed at proximal tubule by Na+/phosphate cotransporters NaPi-IIa aka SLC34A1 (electrogenic - 3Na+ for 1HPO42-) and IIc aka SLC34A3 (electroneutral, minor in adults)
\item 10\% reabsorbed in distal tubule, excreted in urine
\item Regulated by pH and Na+ \(\rightarrow\) intracellular H+ stimulates NaPi activity (allosteric), luminal H+ inhibits Na+ binding, acidic lumen decreases transport (most is in the form H2PO4-, less transportable)
\end{itemize}
\end{itemize}
\begin{center}\includegraphics[width=0.82\textwidth,height=0.68\textheight,keepaspectratio]{\reviewnotesimagepath d6-figure-fluid-electrolyte-balance-regulated-ph-na-intracellular.png}\par\captionof{figure}{The illustration relates to Regulated by pH and Na+ \(\rightarrow\) intracellular H+ stimulates NaPi activity (allosteric), luminal H+ inhibits Na+ binding, acidic lumen decreases transport (most is in, within the topic of Fluid and Electrolyte Balance.}\labfig{appendix-figure-d6-fluid-electrolyte-balance-regulated-ph-na-intracellular}\end{center}
\begin{center}\includegraphics[width=0.82\textwidth,height=0.68\textheight,keepaspectratio]{\reviewnotesimagepath d6-figure-fluid-electrolyte-balance-fluid-electrolyte-balance-05.png}\par\captionof{figure}{The illustration relates to Fluid and Electrolyte Balance, within the topic of Fluid and Electrolyte Balance.}\labfig{appendix-figure-d6-fluid-electrolyte-balance-fluid-electrolyte-balance-05}\end{center}
\begin{center}\includegraphics[width=0.82\textwidth,height=0.68\textheight,keepaspectratio]{\reviewnotesimagepath d6-figure-fluid-electrolyte-balance-fluid-electrolyte-balance-06.png}\par\captionof{figure}{The illustration relates to Fluid and Electrolyte Balance, within the topic of Fluid and Electrolyte Balance.}\labfig{appendix-figure-d6-fluid-electrolyte-balance-fluid-electrolyte-balance-06}\end{center}
\begin{itemize}
\item The calcium balance
\begin{itemize}
\item Regulated by parathyroid hormone, vitamin D
\item Plasma calcium
\begin{itemize}
\item Influenced by acid-base, protein/anion concentrations
\item filterable Ca2+ 60\% (45\% of all = ionized \(\rightarrow\) 1.0-1.3 mM; Ca2+ anion complexes 15\% of total)
\item Non-filterable Ca2+ - 40\%, bound to plasma proteins (mainly albumin)
\end{itemize}
\item Mechanisms of calcium movement
\begin{itemize}
\item Apical entry - TRPV5/V6 (ECaC1/2) channels
\item Intracellular buffering (by Ca2+ binding proteins) for low Ca2+i form
\item Basolateral extrusion by PMCA (Ca2+ ATPase) at low ionic calcium concentrations; NCX (Na+/Ca2+ exchanger) at higher concentrations
\end{itemize}
\item Reabsorption of Ca2+
\begin{itemize}
\item Proximal tubule: bulk Ca²⁺ reabsorption, passive/paracellular
\item TAL: mixed paracellular/transcellular, voltage-dependent, PTH sensitive
\item DCT: active transcellular, hormonally controlled.
\item Ca2+ sensing receptor in interstitium: feedback regulation on extracellular Ca²⁺
\end{itemize}
\end{itemize}
\end{itemize}
\begin{center}\includegraphics[width=0.82\textwidth,height=0.68\textheight,keepaspectratio]{\reviewnotesimagepath d6-figure-fluid-electrolyte-balance-ca2-sensing-receptor-interstitium.png}\par\captionof{figure}{The illustration relates to Ca2+ sensing receptor in interstitium: feedback regulation on extracellular Ca²⁺, within the topic of Fluid and Electrolyte Balance.}\labfig{appendix-figure-d6-fluid-electrolyte-balance-ca2-sensing-receptor-interstitium}\end{center}
\begin{itemize}
\item Magnesium balance
\begin{itemize}
\item Plasma magnesium
\begin{itemize}
\item Total plasma Mg2+ = 0.8-1 mM (1.8-2.2 mg/dL)
\item Filterable Mg2+ =70\% (ionized 10\% of total, anion complex <10\% of total)
\item Protein bound Mg2+ = 30\%, non–filterable
\item Disturbance \(\rightarrow\) neuromuscular, cardiac, vascular effects
\end{itemize}
\item Factors influencing absorption of Mg2+
\begin{itemize}
\item Mg2+/Ca2+ depletion, PTH, ADH glucagon, calcitonin \(\rightarrow\) increase Mg2+ absorption
\item Diuretics (loop-furosemide, osmotic-mannitol) \(\rightarrow\) increase excretion of Mg2+
\end{itemize}
\item Renal magnesium handling
\begin{itemize}
\item TAL: main resorptive, regulatory site
\item Proximal CT: minor passive resorption
\item DCT, collecting duct: fine–tuning (by hormones)
\item Calcium sensing receptor, lumen positive voltage \(\rightarrow\) Mg2+, Ca2+ handling (linked)
\end{itemize}
\end{itemize}
\end{itemize}
\begin{center}\includegraphics[width=0.82\textwidth,height=0.68\textheight,keepaspectratio]{\reviewnotesimagepath d6-figure-fluid-electrolyte-balance-calcium-sensing-receptor-lumen.png}\par\captionof{figure}{The illustration relates to Calcium sensing receptor, lumen positive voltage \(\rightarrow\) Mg2+, Ca2+ handling (linked), within the topic of Fluid and Electrolyte Balance.}\labfig{appendix-figure-d6-fluid-electrolyte-balance-calcium-sensing-receptor-lumen}\end{center}
\begin{itemize}
\item Potassium balance
\begin{itemize}
\item K+ in the body
\begin{itemize}
\item Most abundant intracellular cation, 98\% inside cells
\item Normal plasma level 3.5-5.0 mM
\item Functions: maintain cell volume, pH, enzyme activity, DNA/protein synthesis, cell growth, cell membrane potential (maintained by low serum K+)
\item Majority located in muscle cells \(\rightarrow\) 1\% shift can double or halve
\end{itemize}
\item Homeostasis of K+
\begin{itemize}
\item Excretion: 90-95\% by kidney, 5-10\% by colon
\item Hormonal regulation: insulin, epinephrine, aldosterone promotes K+ cellular uptake
\item Acidemia \(\rightarrow\) K+ shifts out of cells \(\rightarrow\) hyperkalemia (and vice versa)
\end{itemize}
\item Renal potassium handling
\begin{itemize}
\item K+ filtered, resorbed and secreted by kidneys (normal GFR = 800 mmol/day of K+)
\item Dietary intake K+ = 80-120 mmol/day
\item Aldosterone promotes renal secretion of K+
\item Proximal tubule, loop of henle: resorbs 90\% of filtered K+
\item Distal nephron (distal CT to medullary collecting duct \(\rightarrow\) resorption or secretion of K+ depending on intake, hormones
\item Alkalosis \(\rightarrow\) increased K+ secretion \(\rightarrow\) hypokalemia (and vice versa)
\end{itemize}
\end{itemize}
\end{itemize}
\begin{itemize}
\item Factors affecting K+ secretion
\begin{itemize}
\item Increased (distal nephron) urine flow: better secretion of K+, increased luminal Na+ = Na-K pump stimulated; K+ secreted
\item Low luminal Cl: increases K-Cl cotransporter
\item Aldosterone: increased secretion (Na-K pump, apical ENaCs, apical K+ channels upregulated),
\item Glucocorticoids: increased excretion by increased distal flow (indirect)
\item Epinephrine: decreased excretion (direct inhibition, shifts K+ to cells)
\item ADH: mixed effect (increased apical Na+ conductance, K+ permeability, reduced urine flow)
\item Loop diuretics (furosemide, bumetanide): increased distal Na+ delivery \(\rightarrow\) enhanced K+ secretion
\item K+ sparing diuretics (amelioride): blocked ENaCs \(\rightarrow\) reduced + secretion
\item High K+ intake: stimulates distal K+ secretion (aldosterone dependant + independent)
\item Low K+ intake: resorption by H- pumps in distal nephron
\end{itemize}
\end{itemize}
\begin{center}\includegraphics[width=0.82\textwidth,height=0.68\textheight,keepaspectratio]{\reviewnotesimagepath d6-figure-fluid-electrolyte-balance-low-intake-resorption-pumps.png}\par\captionof{figure}{The illustration relates to Low K+ intake: resorption by H- pumps in distal nephron, within the topic of Fluid and Electrolyte Balance.}\labfig{appendix-figure-d6-fluid-electrolyte-balance-low-intake-resorption-pumps}\end{center}
\begin{center}\includegraphics[width=0.82\textwidth,height=0.68\textheight,keepaspectratio]{\reviewnotesimagepath d6-figure-fluid-electrolyte-balance-fluid-electrolyte-balance-07.png}\par\captionof{figure}{The illustration relates to Fluid and Electrolyte Balance, within the topic of Fluid and Electrolyte Balance.}\labfig{appendix-figure-d6-fluid-electrolyte-balance-fluid-electrolyte-balance-07}\end{center}
\begin{itemize}
\item K+ influences on membrane potential
\begin{itemize}
\item Low K+ lowers membrane potential \(\rightarrow\) harder to depolarize
\item High K+ increases membrane potential \(\rightarrow\) hyperexcitable
\end{itemize}
\end{itemize}
\begin{center}\includegraphics[width=0.82\textwidth,height=0.68\textheight,keepaspectratio]{\reviewnotesimagepath d6-figure-fluid-electrolyte-balance-high-increases-membrane-potential.png}\par\captionof{figure}{The illustration relates to High K+ increases membrane potential \(\rightarrow\) hyperexcitable, within the topic of Fluid and Electrolyte Balance.}\labfig{appendix-figure-d6-fluid-electrolyte-balance-high-increases-membrane-potential}\end{center}
\begin{center}\includegraphics[width=0.82\textwidth,height=0.68\textheight,keepaspectratio]{\reviewnotesimagepath d6-figure-fluid-electrolyte-balance-fluid-electrolyte-balance-08.png}\par\captionof{figure}{The illustration relates to Fluid and Electrolyte Balance, within the topic of Fluid and Electrolyte Balance.}\labfig{appendix-figure-d6-fluid-electrolyte-balance-fluid-electrolyte-balance-08}\end{center}
\begin{itemize}
\item Role of diuretics
\begin{itemize}
\item Main mechanism: reversible inhibition of Na+ reabsorption at sites in nephron \(\rightarrow\) negative Na+ balance \(\rightarrow\) water excreted
\item Other effects: increased excretion of Cl-, K+, Ca2+, H+ \(\rightarrow\) metabolic alkalosis
\item Increase Na+, water excretion \(\rightarrow\) reduce ECF volume in hypertension, edema (e.g. heart failure, cirrhosis, nephrotic syndrome)
\item Long term use results in decreased ECF volume \(\rightarrow\) increased sympathetic tone, aldosterone, vasopressin, adaptive distal nephron hypertrophy and decreased ANP
\item Classes, site of action
\end{itemize}
\end{itemize}
\begin{center}\includegraphics[width=0.82\textwidth,height=0.68\textheight,keepaspectratio]{\reviewnotesimagepath d6-figure-fluid-electrolyte-balance-classes-site-action.png}\par\captionof{figure}{The illustration relates to Classes, site of action, within the topic of Fluid and Electrolyte Balance.}\labfig{appendix-figure-d6-fluid-electrolyte-balance-classes-site-action}\end{center}
\begin{itemize}
\item Note: may have synergistic effects (e.g. loop+thiazide)
\item Delivery into tubule: organic anion/cation transporters
\item Effect of diuretics on electrolytes
\end{itemize}
\begin{center}\includegraphics[width=0.82\textwidth,height=0.68\textheight,keepaspectratio]{\reviewnotesimagepath d6-figure-fluid-electrolyte-balance-effect-diuretics-electrolytes.png}\par\captionof{figure}{The illustration relates to Effect of diuretics on electrolytes, within the topic of Fluid and Electrolyte Balance.}\labfig{appendix-figure-d6-fluid-electrolyte-balance-effect-diuretics-electrolytes}\end{center}
\begin{center}\includegraphics[width=0.82\textwidth,height=0.68\textheight,keepaspectratio]{\reviewnotesimagepath d6-figure-fluid-electrolyte-balance-fluid-electrolyte-balance-09.png}\par\captionof{figure}{The illustration relates to Fluid and Electrolyte Balance, within the topic of Fluid and Electrolyte Balance.}\labfig{appendix-figure-d6-fluid-electrolyte-balance-fluid-electrolyte-balance-09}\end{center}

\section{Urinary Assessment and Disorders}

\subsection{Acid–Base Balance}
\begin{itemize}
\item Outline:
\begin{itemize}
\item Normal acid-base balance (and maintenance of balance)
\item Abnormal balance, steps to interpret
\item Compensatory responses in abnormal balance, analysis + interpretation
\end{itemize}
\item Physiologic buffer systems
\end{itemize}
\begin{center}\includegraphics[width=0.82\textwidth,height=0.68\textheight,keepaspectratio]{\reviewnotesimagepath d6-figure-acid-base-balance-physiologic-buffer-systems.png}\par\captionof{figure}{The illustration relates to Physiologic buffer systems, within the topic of Acid–Base Balance.}\labfig{appendix-figure-d6-acid-base-balance-physiologic-buffer-systems}\end{center}
\begin{center}\includegraphics[width=0.82\textwidth,height=0.68\textheight,keepaspectratio]{\reviewnotesimagepath d6-figure-acid-base-balance-acid-base-balance.png}\par\captionof{figure}{The illustration relates to Acid–Base Balance, within the topic of Acid–Base Balance.}\labfig{appendix-figure-d6-acid-base-balance-acid-base-balance}\end{center}
\begin{itemize}
\item Generation of bicarbonate
\begin{itemize}
\item CO2 and H2O are converted to/from bicarbonate by carbonic anhydrase (proximal tubule)
\end{itemize}
\end{itemize}
\begin{center}\includegraphics[width=0.82\textwidth,height=0.68\textheight,keepaspectratio]{\reviewnotesimagepath d6-figure-acid-base-balance-co2-h2o-converted-bicarbonate.png}\par\captionof{figure}{The illustration relates to CO2 and H2O are converted to/from bicarbonate by carbonic anhydrase (proximal tubule), within the topic of Acid–Base Balance.}\labfig{appendix-figure-d6-acid-base-balance-co2-h2o-converted-bicarbonate}\end{center}
\begin{itemize}
\item Processing of bicarbonate in the kidney
\begin{itemize}
\item Filtered by glomerulus = 4320 mmol/day (24 mM \(\times\) 180 L)
\item Reabsorption: 80\% at proximal tubule, 10\% at thick ascending limb, 10\% at distal nephron
\item Lost to urine = 0.01\% (0.36 mmol/day)
\item Newly formed bicarbonate sources
\begin{itemize}
\item generation of titratable acid \(\rightarrow\) proximal tubule, cortical collecting ducts each 15 mmol/day
\item Ammonium excretion \(\rightarrow\) proximal tubules excrete glutamine
\end{itemize}
\end{itemize}
\end{itemize}
\begin{center}\includegraphics[width=0.82\textwidth,height=0.68\textheight,keepaspectratio]{\reviewnotesimagepath d6-figure-acid-base-balance-ammonium-excretion-proximal-tubules.png}\par\captionof{figure}{The illustration relates to Ammonium excretion \(\rightarrow\) proximal tubules excrete glutamine, within the topic of Acid–Base Balance.}\labfig{appendix-figure-d6-acid-base-balance-ammonium-excretion-proximal-tubules}\end{center}
\begin{center}\includegraphics[width=0.82\textwidth,height=0.68\textheight,keepaspectratio]{\reviewnotesimagepath d6-figure-acid-base-balance-acid-base-balance-02.png}\par\captionof{figure}{The illustration relates to Acid–Base Balance, within the topic of Acid–Base Balance.}\labfig{appendix-figure-d6-acid-base-balance-acid-base-balance-02}\end{center}
\begin{itemize}
\item Acid-base disturbances
\begin{itemize}
\item Primary organ failure \(\rightarrow\) Acid/alkalosis (the processes tending to increase/decrease pH) \(\rightarrow\) acid/alkalemia (pH < 7.35 or pH > 7.45)
\item Acidosis can be with compensation (no acidemia) or with insufficient compensation (acidemia) \(\leftarrow\) also applies for alkalemia
\end{itemize}
\end{itemize}
\begin{center}\includegraphics[width=0.82\textwidth,height=0.68\textheight,keepaspectratio]{\reviewnotesimagepath d6-figure-acid-base-balance-acidosis-be-compensation-no.png}\par\captionof{figure}{The illustration relates to Acidosis can be with compensation (no acidemia) or with insufficient compensation (acidemia) \(\leftarrow\) also applies for alkalemia, within the topic of Acid–Base Balance.}\labfig{appendix-figure-d6-acid-base-balance-acidosis-be-compensation-no}\end{center}
\begin{itemize}
\item Respiratory acidosis
\begin{itemize}
\item Mechanism: alveolar hypoventilation (less than 15,000 mEq/day) \(\rightarrow\) increased CO2 retention \(\rightarrow\) increased pCO2 \(\rightarrow\) increased H2CO3 \(\rightarrow\) decreased pH
\item Causes: COPD/asthma/airflow restriction, pneumonia, respiratory failure \(\leftarrow\) neuromuscular
\item Symptoms: cyanosis (low O2), increased HR, RR \(\rightarrow\) respiratory distress (shallow, fast breathing)
\item Renal compensation of respiratory acidosis
\begin{itemize}
\item Proximal tubule: increased H+ secretion, reabsorption of HCO3-, ammoniagenesis (new bicarbonate)
\item Collecting duct: increased H⁺ secretion (α-intercalated cells), HCO3- retention in blood
\item Monitoring: level of consciousness, respiratory rate and depth, arterial blood gas (ABG), ECG
\item Treatment: positive pressure ventilation, bronchodilators (asthma/COPD)
\end{itemize}
\end{itemize}
\item Respiratory alkalosis
\begin{itemize}
\item Low blood PCO2 (< 35 mmHg, pH > 7.45)
\item CO2 + H2O \(\Leftrightarrow\) H2CO3 \(\Leftrightarrow\) H+ + HCO3-
\item Causes
\begin{itemize}
\item Hyperventilation
\item High altitude (low O2, hyperventilation)
\item Fever/sepsis \(\rightarrow\) increased respiratory rate
\end{itemize}
\item Consequences
\begin{itemize}
\item Decreased serum potassium, phosphate
\item Calcium fixed in albumin binding (replacing H+) \(\rightarrow\) lowered membrane threshold \(\rightarrow\) hyperexcitability
\item H+ shifts out of cells, K+, PO43- shifts into cells \(\rightarrow\) hypokalemia, hypophosphatemia \(\rightarrow\) decreased 2,3-DPG \(\rightarrow\) increased O2 affinity of RBC
\end{itemize}
\end{itemize}
\end{itemize}
\begin{center}\includegraphics[width=0.82\textwidth,height=0.68\textheight,keepaspectratio]{\reviewnotesimagepath d6-figure-acid-base-balance-shifts-out-cells-po43.png}\par\captionof{figure}{The illustration relates to H+ shifts out of cells, K+, PO43- shifts into cells \(\rightarrow\) hypokalemia, hypophosphatemia \(\rightarrow\) decreased 2,3-DPG \(\rightarrow\) increased O2 affinity of RBC, within the topic of Acid–Base Balance.}\labfig{appendix-figure-d6-acid-base-balance-shifts-out-cells-po43}\end{center}
\begin{itemize}
\item Symptoms
\begin{itemize}
\item Dizziness, tingling, shortness of breath
\item Hyperventilation
\item Trosseau sign (carpopedal spasm)
\item Paresthesia (tingling, numbness), tetany
\item Lightheadedness, confusion, syncope
\end{itemize}
\item Compensatory responses: body buffering system, shifting H+ into cells (exchanged for K+, phosphate)
\item Treatment
\begin{itemize}
\item Respiratory mask with bag or paper bag, underlying condition
\item For ventilator \(\rightarrow\) ventilation adjustment, PO2 and PCO2 monitoring
\end{itemize}
\item Metabolic acidosis
\begin{itemize}
\item Increased acid production
\begin{itemize}
\item Lactic acidosis, ketoacidosis
\item Chronic renal failure (sulfate, phosphate, urea accumulation)
\item Intoxications (organic acids \(\rightarrow\) salicylates, ethanol, methanol, formaldehyde, isoniazid)
\item Massive rhabdomyolysis
\end{itemize}
\item Increased bicarbonate loss
\begin{itemize}
\item Longstanding diarrhea \(\rightarrow\) bicarbonate loss
\item Renal tubular acidosis
\end{itemize}
\item Impaired renal acid secretion \(\rightarrow\) renal failure prevents H+ secretion, HCO3- regeneration
\item Respiratory response:
\begin{itemize}
\item Deep laboured gasping, increased respiratory rate (reduce CO2) \(\rightarrow\) Kussmaul breathing
\item Decreased pH \(\rightarrow\) central chemoreceptor reflex \(\rightarrow\) increased respiratory rate \(\rightarrow\) decreased PaCO2 \(\rightarrow\) increased pH
\end{itemize}
\item Renal response: increased secretion of H+ in proximal tubule \(\rightarrow\) increased proximal tubule ammonia production (in chronic)
\item Treatment: treat underlying disorder, infuse sodium bicarbonate IV, correct ECF volume (esp. renal disease), diabetic ketoacidosis \(\rightarrow\) glucose + insulin
\item The anion gap in acidosis
\begin{itemize}
\item Quantifies the amount of unmeasured anions in circulation
\item Measured cation: Na
\item Measured anions: Cl-, bicarbonate
\item Unmeasured anions: albumin, phosphate, organic acids \(\rightarrow\) “anion gap”
\item AG = [Na+] - [Cl-] - [HCO3-] \(\rightarrow\) normally between 8-12 mEq/L
\item Normal range for patients with reduced albumin = lower gap (-2.5 for every 1g/dL albumin lost)
\end{itemize}
\item High anion gap metabolic acidosis (HAGMA)
\begin{itemize}
\item Implies creation of new acid which dissociates into cation and anion. The cation is neutralized, new anion (unmeasured) is created
\item Common causes: MUDPILES \(\rightarrow\) methanol, uremia (renal failure), diabetic ketoacidosis, propylene glycol, isoniazid/iron, lactic acidosis, ethylene glycol, salicylates
\end{itemize}
\item Normal anion gap metabolic acidosis (NAGMA)
\begin{itemize}
\item Implies loss of base without addition of new acid
\item Common causes: GI loss (diarrhea \(\rightarrow\) bicarbonate loss), renal loss (renal tubular acidosis), Addison’s disease
\end{itemize}
\end{itemize}
\item Metabolic alkalosis
\begin{itemize}
\item High HCO3- in blood with pH > 7.4
\item Renal compensation = HCO3- secreted into lumen, H+ absorbed
\item Causes
\begin{itemize}
\item Vomiting (loss of H+)
\item Hyperaldosteronism (reabsorbed Na, excrete H+)
\item Massive transfusion (citrate \(\rightarrow\) bicarbonate)
\end{itemize}
\item Symptoms: decreased respiratory rate, numbness, muscle weakness (hypocalcemia)
\item Treatment: by symptom, hemodialysis
\end{itemize}
\item Common assessment in acid–base disorders
\begin{itemize}
\item pH of blood (normal range 7.35-7.45)
\item PaCO2, HCO3-, ABG \(\rightarrow\) identify the underlying disturbance
\end{itemize}
\end{itemize}
\begin{center}\includegraphics[width=0.82\textwidth,height=0.68\textheight,keepaspectratio]{\reviewnotesimagepath d6-figure-acid-base-balance-paco2-hco3-abg-identify.png}\par\captionof{figure}{The illustration relates to PaCO2, HCO3-, ABG \(\rightarrow\) identify the underlying disturbance, within the topic of Acid–Base Balance.}\labfig{appendix-figure-d6-acid-base-balance-paco2-hco3-abg-identify}\end{center}
\begin{center}\includegraphics[width=0.82\textwidth,height=0.68\textheight,keepaspectratio]{\reviewnotesimagepath d6-figure-acid-base-balance-acid-base-balance-03.png}\par\captionof{figure}{The illustration relates to Acid–Base Balance, within the topic of Acid–Base Balance.}\labfig{appendix-figure-d6-acid-base-balance-acid-base-balance-03}\end{center}

\subsection{Micturition and Abnormality}
\begin{itemize}
\item LOs:
\begin{itemize}
\item Anatomy of lower urinary tract and pelvic floor muscle
\item Normal physiology of micturition reflex
\item Abnormal micturition
\begin{itemize}
\item Urinary incontinence
\item Urinary retention
\item Infection of lower urinary tract
\end{itemize}
\item Diagnosis of abnormal micturition
\item Treatment for abnormal micturition
\end{itemize}
\item Anatomy of lower urinary tract
\begin{itemize}
\item Urinary bladder
\begin{itemize}
\item Detrusor muscle (smooth muscle) - at wall of urinary bladder
\item Neck - connects between bladder trigone and posterior urethra
\item Internal urethral sphincter - multiunit smooth muscle
\item External urethral sphincter - striated muscle
\end{itemize}
\item Urethra
\item Prostate gland (male)
\item Pelvic floor muscles
\begin{itemize}
\item Levator ani muscles: pubococcygeus, iliococcygeus, puborectalis \(\rightarrow\) support pelvic organs, control flow of defecation and urination
\end{itemize}
\end{itemize}
\end{itemize}
\begin{center}\includegraphics[width=0.82\textwidth,height=0.68\textheight,keepaspectratio]{\reviewnotesimagepath d6-figure-micturition-abnormality-levator-ani-muscles-pubococcygeus.png}\par\captionof{figure}{The illustration relates to Levator ani muscles: pubococcygeus, iliococcygeus, puborectalis \(\rightarrow\) support pelvic organs, control flow of defecation and urination, within the topic of Micturition and Abnormality.}\labfig{appendix-figure-d6-micturition-abnormality-levator-ani-muscles-pubococcygeus}\end{center}
\begin{itemize}
\item Control of bladder emptying:
\begin{itemize}
\item Holding of urge
\begin{itemize}
\item Activated external urethral sphincter
\item Suppression of detrusor muscle contraction (sympathetic nerves)
\item Increased pelvic floor activity
\end{itemize}
\item Voiding without urge
\begin{itemize}
\item Conscious relaxation of external urethral sphincter
\item Activation of parasympathetic pathways (detrusor contraction, internal urethral sphincter reflex
\end{itemize}
\end{itemize}
\end{itemize}
\begin{itemize}
\item Central control of micturition
\begin{itemize}
\item Micturition is controlled by pontine micturition center in medial dorsal pons \(\rightarrow\) coordinates detrusor muscle, urinary sphincters
\item Sympathetic innervation of urinary bladder
\begin{itemize}
\item preganglionic neurons in T10-L2 of spine \(\rightarrow\) synapse in hypogastric nerve \(\rightarrow\) hypogastric nerve to bladder
\item Inhibits detrusor (bladder) contraction \(\rightarrow\) relaxation and filling (by norepinephrine on β3 receptors)
\item Stimulates internal urethral sphincter, urethra, bladder neck contraction \(\rightarrow\) no urine released (norepinephrine on \(\alpha\) receptors)
\end{itemize}
\item Parasympathetic innervation of urinary bladder
\begin{itemize}
\item Originating from S2-S4 spinal segments \(\rightarrow\) pelvic splanchnic nerve (cholinergic) \(\rightarrow\) synapse on postganglionic neurons (muscarinic) on bladder wall
\item Stimulates detrusor contraction, bladder emptying
\item Relaxes internal urethral sphincter \(\rightarrow\) allows urine flow
\end{itemize}
\end{itemize}
\end{itemize}
\begin{center}\includegraphics[width=0.82\textwidth,height=0.68\textheight,keepaspectratio]{\reviewnotesimagepath d6-figure-micturition-abnormality-relaxes-internal-urethral-sphincter.png}\par\captionof{figure}{The illustration relates to Relaxes internal urethral sphincter \(\rightarrow\) allows urine flow, within the topic of Micturition and Abnormality.}\labfig{appendix-figure-d6-micturition-abnormality-relaxes-internal-urethral-sphincter}\end{center}
\begin{itemize}
\item Afferent signalling through pelvic splanchnic nerve to pontine micturition center
\item Efferent signalling through presynaptic parasympathetic neurons (S1-4) \(\rightarrow\) bladder, sphincters
\item Centers responsible for control of micturition: prefrontal cortex, insula
\begin{itemize}
\item Voluntary control (conscious regulation of reflex) \(\rightarrow\) suppresses urge during inappropriate times/not wanting to void. Achieved by overriding pontine micturition center (PMC) temporarily
\item Impairment of voluntary control (stroke, dementia) \(\rightarrow\) urinary incontinence
\end{itemize}
\item Periaqueductal grey in midbrain
\begin{itemize}
\item Integrates sensory information and motor response for micturition
\item Interacts with the PMC and spinal cord \(\rightarrow\) controls detrusor
\end{itemize}
\end{itemize}
\begin{center}\includegraphics[width=0.82\textwidth,height=0.68\textheight,keepaspectratio]{\reviewnotesimagepath d6-figure-micturition-abnormality-interacts-pmc-spinal-cord.png}\par\captionof{figure}{The illustration relates to Interacts with the PMC and spinal cord \(\rightarrow\) controls detrusor, within the topic of Micturition and Abnormality.}\labfig{appendix-figure-d6-micturition-abnormality-interacts-pmc-spinal-cord}\end{center}
\begin{itemize}
\item Control of Micturition
\begin{itemize}
\item Somatic controls from S2-3 \(\rightarrow\) pudendal nerve \(\rightarrow\) external urethral sphincter (contraction/inhibition)
\item Storage phase (of micturition): detrusor muscle contraction inhibited, bladder filling promoted. Sphincters remain closed
\item Voiding phase (of micturition): pontine micturition center relaxes detrusor inhibition \(\rightarrow\) bladder contracts \(\rightarrow\) sphincters relax, urine flows
\item Urine storage reflex
\begin{itemize}
\item Bladder stretches during filling \(\rightarrow\) mechanoreceptors detect expansion
\item Afferent pathway: signal travels by pelvic splanchnic nerves \(\rightarrow\) sacral spinal cord \(\rightarrow\) PMC \(\rightarrow\) higher center receives message
\item Efferent pathway: hypogastric nerve activated \(\rightarrow\) internal urethral sphincter contraction + detrusor relaxation, pudendal nerve activated \(\rightarrow\) external urethral sphincter contraction. Urine retained, desire to void inhibited
\end{itemize}
\end{itemize}
\end{itemize}
\begin{center}\includegraphics[width=0.82\textwidth,height=0.68\textheight,keepaspectratio]{\reviewnotesimagepath d6-figure-micturition-abnormality-efferent-pathway-hypogastric-nerve.png}\par\captionof{figure}{The illustration relates to Efferent pathway: hypogastric nerve activated \(\rightarrow\) internal urethral sphincter contraction + detrusor relaxation, pudendal nerve activated \(\rightarrow\) external urethral sphi, within the topic of Micturition and Abnormality.}\labfig{appendix-figure-d6-micturition-abnormality-efferent-pathway-hypogastric-nerve}\end{center}
\begin{center}\includegraphics[width=0.82\textwidth,height=0.68\textheight,keepaspectratio]{\reviewnotesimagepath d6-figure-micturition-abnormality-micturition-abnormality.png}\par\captionof{figure}{The illustration relates to Micturition and Abnormality, within the topic of Micturition and Abnormality.}\labfig{appendix-figure-d6-micturition-abnormality-micturition-abnormality}\end{center}
\begin{itemize}
\item Detection of bladder volume
\begin{itemize}
\item Urge to void begins at 150 mL. Full at 400-500 mL
\item Increased filling \(\rightarrow\) increased micturition reflex activity
\end{itemize}
\item Process of voiding
\begin{itemize}
\item Bladder filled
\item Stretch signals sent from pelvic nerve to cortex
\item PMC releases inhibitory signalling
\item Internal (by pelvic nerve) and external (by pudendal nerve) sphincters relax
\item Bladder contraction, urine expelled (with help from abdominal wall, pelvic floor)
\end{itemize}
\end{itemize}
\begin{center}\includegraphics[width=0.82\textwidth,height=0.68\textheight,keepaspectratio]{\reviewnotesimagepath d6-figure-micturition-abnormality-bladder-contraction-urine-expelled.png}\par\captionof{figure}{The illustration relates to Bladder contraction, urine expelled (with help from abdominal wall, pelvic floor), within the topic of Micturition and Abnormality.}\labfig{appendix-figure-d6-micturition-abnormality-bladder-contraction-urine-expelled}\end{center}
\begin{itemize}
\item Post-void state
\begin{itemize}
\item Typically, residual urine < 5-10 mL
\item < 50 mL residual normal
\item > 200 mL residual = urinary retention
\end{itemize}
\item Abnormalities in micturition
\begin{itemize}
\item Urinary incontinence = involuntary, unintentional discharge of urine (drops/leak)
\end{itemize}
\end{itemize}
\begin{center}\includegraphics[width=0.82\textwidth,height=0.68\textheight,keepaspectratio]{\reviewnotesimagepath d6-figure-micturition-abnormality-urinary-incontinence-involuntary-unintentional.png}\par\captionof{figure}{The illustration relates to Urinary incontinence = involuntary, unintentional discharge of urine (drops/leak), within the topic of Micturition and Abnormality.}\labfig{appendix-figure-d6-micturition-abnormality-urinary-incontinence-involuntary-unintentional}\end{center}
\begin{itemize}
\item Stress urinary incontinence
\begin{itemize}
\item Occurrence: F > M
\item Physical stimuli may induce, such as exertion, coughing, sneezing, laughing, lifting heavy loads
\item Caused by increased abdominal pressure, impaired sphincter function (from pelvic floor weakness, nerve malfunction)
\item Factors driving weak pelvic floor: age, pregnancy, vaginal birth, surgery, weak connective tissue
\end{itemize}
\item Urge urinary incontinence
\begin{itemize}
\item Occurence: most common in elderly
\item Sudden, irresistible urge to urinate
\item Causes: small bladder, bladder infections, unstable bladder muscle, tumors
\item Characteristics: occurs at short intervals (< 30 min interval), often without full bladder
\end{itemize}
\item Overflow urinary incontinence
\begin{itemize}
\item Occurance: M > F
\item Pressure in bladder exceeds bladder sphincter pressure \(\rightarrow\) urine drip
\item Causes: drainage disorder (e.g. urinary stones, tumors, nerve damage, benign prostate hypertrophy)
\item Characteristics: bladder overfilling \(\rightarrow\) continuous release of small amounts of urine \(\rightarrow\) absent urge to urinate
\end{itemize}
\item Neurogenic bladder
\begin{itemize}
\item Neurologic dysfunction/insult \(\rightarrow\) bladder dysfunction
\item Causes: internal/external trauma, disease, injury
\item Common causes: autonomic/peripheral nerve injury (autonomic neuropathy; diabetes mellitus; Guillain-Barré syndrome; herpes zoster virus), brain (cerebrovascular disease; concussion; multiple sclerosis; neoplasm or tumor), spinal cord (multiple sclerosis; spinal cord trauma; spinal stenosis; disk herniation)
\end{itemize}
\item Drug induced urinary incontinence
\end{itemize}
\begin{center}\includegraphics[width=0.82\textwidth,height=0.68\textheight,keepaspectratio]{\reviewnotesimagepath d6-figure-micturition-abnormality-drug-induced-urinary-incontinence.png}\par\captionof{figure}{The illustration relates to Drug induced urinary incontinence, within the topic of Micturition and Abnormality.}\labfig{appendix-figure-d6-micturition-abnormality-drug-induced-urinary-incontinence}\end{center}
\begin{itemize}
\item Urinary retention (overview)
\begin{itemize}
\item Incomplete bladder emptying following urination
\item Causes: obstruction, infection, neurologic (neurogenic bladder), medication
\item Symptoms: severe abdominal pain, inability to urinate, asymptomatic
\end{itemize}
\item Obstructive urinary retention
\begin{itemize}
\item Male causes: benign prostate hypertrophy (most common), urolithiasis, prostate cancer, meatal stenosis, paraphimosis (foreskin stuck back) /phimosis (foreskin stuck forwards)
\item Female causes: organ prolapse (cysto/rectocele), pelvic mass (malignancy), urolithiasis, fecal impaction
\end{itemize}
\item Infectious urinary retention
\begin{itemize}
\item Male causes: balanitis (inflamed glans of penis), prostate abscess
\item Female causes: acute vulvovaginitis, urinary tract infection
\end{itemize}
\item Neurologic urinary retention: spinal cord injuries: Cauda equina syndrome (compressed nerve roots), cord compression, transverse myelitis (immune autoresponse), spinal trauma
\item Pharmacological urinary retention
\begin{itemize}
\item Anticholinergic drugs, tricyclic antidepressants
\item Opioid analgesics \(\rightarrow\) CNS depression + lowered bladder reflexes
\item Antiparkinson agents
\item Antidepressants with strong anticholinergic effect
\item Sympathomemetrics: ephedrine, phenylepherine
\end{itemize}
\item Urinary tract infection
\begin{itemize}
\item Inflammation of urinary epithelium in response to bacteria
\item Common infections: acute cystitis, acute/chronic pyelonephritis
\item Response depends on virulence of pathogen, host defences
\item Risk factors—elderly with: Immobility, constipation, fecal/urinary incontinence, urinary retention, altered metal status or systemic disease
\item Common pathogen: \emph{Escherichia coli}
\item Acute cystitis = inflammation of the bladder \(\rightarrow\) frequency, dysuria (urination pain), urgency, lower abdominal/suprapubic pain, hematuria
\item Inflammation response to pathogen at bladder wall \(\rightarrow\) stimulation of stretch receptors, stretch receptor hypersensitivity (overactive bladder) \(\rightarrow\) urgency, frequent urination
\item Treatment of UTI
\begin{itemize}
\item Antimicrobial (broad spectrum antibiotic, e.g. fluoroquinolone)
\item Increased fluid intake to 3-4 L/day
\item Avoidance of bladder irritants (alcohol, caffeine)
\item Urinary analgesics for pain reduction
\item Self–care for recurrent UTI
\end{itemize}
\item Treatment of UTI (cont. pharmacological)
\begin{itemize}
\item First–line for mild cases: Nitrofurantoin 100 mg bid
\item Fluoroquinolone (norfloxacin 400 mg bid = twice daily, ofloxacin 200 mg bid, ciprofloxacin 250 mg bid)
\item Second line drug (e.g. in pregnancy): 3rd generation oral cephalosporin (cefdinir, cefixime, cefditoren)
\item Moderate to severe: 3rd generation IV cephalosporin
\end{itemize}
\item Interstitial cystitis = bladder pain syndrome, caused by breakdown of lining
\item Nonbacterial infectious cystitis
\begin{itemize}
\item Common in women aged 20-30
\item Symptoms: bladder fullness, frequency, small urine volume, chronic pelvic pain with negative urine culture
\item Treatment: pain management, improve coping (no single treatment)
\end{itemize}
\end{itemize}
\item Diagnostic tests in disorders of micturition
\begin{itemize}
\item Urodynamic studies
\begin{itemize}
\item Uroflowmetry: measures speed and flow of urine during urination (non-invasive) \(\rightarrow\) mL/sec \(\rightarrow\) post void residual volume
\item Cystometry: measures the pressure inside the bladder (during external filling with water), bladder’s ability to empty \(\rightarrow\) determines compliance and capacity of bladder
\item Electromyography: measures electrical activity of muscles involved in urination
\item Pressure–flow studies: measures pressure inside urethra and bladder during urination
\item Videourodynamic studies: X-ray, video visualization of urinary tract during urination
\end{itemize}
\end{itemize}
\item Treatment of abnormal micturition (urine retention)
\begin{itemize}
\item Bladder retraining
\item Kegel exercise to strengthen pelvic floor, muscles surrounding bladder, urethra
\item Pharmacological treatment (benign prostatic hyperplasia) \(\rightarrow\) alpha blockers, 5-alpha reductase inhibitors (block testosterone, dihydrotestosterone)
\item Prostate removal (partial)
\item Urethral dilation
\end{itemize}
\end{itemize}

\subsection{Principle of Urine Specimen collection, Urinalysis, and Urine Culture}
\begin{itemize}
\item Urine analysis (urinalysis)
\begin{itemize}
\item Purposes: general health evaluation, diagnosis of disease/disorder of kidneys/urinary tract, diagnosis of other systemic diseases affecting kidney function, monitoring of diabetes, screening for drug abuse
\item Examinations: macroscopic, chemical, microscopic
\end{itemize}
\item Urine specimen types, collection, preservation
\begin{itemize}
\item Separated into 3 main variables: collection method, timing, handling
\item Urine containers, methods depend on test ordered
\item Notes:
\begin{itemize}
\item Single random specimen may be inaccurate if too dilute
\item First morning specimen can also be used for bacterial culture
\item Timed specimens also include 2 hour postprandial specimens. Used for substances with diurnal variation (e.g. creatinine, urine urea nitrogen, glucose, sodium, potassium, hormones)
\item Catheterized specimens are used when the patient cannot urinate independently
\end{itemize}
\end{itemize}
\end{itemize}
\begin{center}\includegraphics[width=0.82\textwidth,height=0.68\textheight,keepaspectratio]{\reviewnotesimagepath d6-figure-principle-urine-specimen-catheterized-specimens-when-patient.png}\par\captionof{figure}{The illustration relates to Catheterized specimens are used when the patient cannot urinate independently, within the topic of Principle of Urine Specimen collection, Urinalysis, and Urine Culture.}\labfig{appendix-figure-d6-principle-urine-specimen-catheterized-specimens-when-patient}\end{center}
\begin{center}\includegraphics[width=0.82\textwidth,height=0.68\textheight,keepaspectratio]{\reviewnotesimagepath d6-figure-principle-urine-specimen-principle-urine-specimen-collection.png}\par\captionof{figure}{The illustration relates to Principle of Urine Specimen collection, Urinalysis, and Urine Culture, within the topic of Principle of Urine Specimen collection, Urinalysis, and Urine Culture.}\labfig{appendix-figure-d6-principle-urine-specimen-principle-urine-specimen-collection}\end{center}
\begin{itemize}
\item Specimen collection, transport guidelines
\begin{itemize}
\item Transport by clean container free of particles/interfering substances
\item Containers should be leak–resistant with secure lid, not reused
\item Labels: patient name, identification, collection date + time, method of collection
\item Analyzed within 2 hours from collection
\end{itemize}
\end{itemize}
\begin{center}\includegraphics[width=0.82\textwidth,height=0.68\textheight,keepaspectratio]{\reviewnotesimagepath d6-figure-principle-urine-specimen-analyzed-hours-collection.png}\par\captionof{figure}{The illustration relates to Analyzed within 2 hours from collection, within the topic of Principle of Urine Specimen collection, Urinalysis, and Urine Culture.}\labfig{appendix-figure-d6-principle-urine-specimen-analyzed-hours-collection}\end{center}
\begin{center}\includegraphics[width=0.82\textwidth,height=0.68\textheight,keepaspectratio]{\reviewnotesimagepath d6-figure-principle-urine-specimen-principle-urine-specimen-collection-02.png}\par\captionof{figure}{The illustration relates to Principle of Urine Specimen collection, Urinalysis, and Urine Culture, within the topic of Principle of Urine Specimen collection, Urinalysis, and Urine Culture.}\labfig{appendix-figure-d6-principle-urine-specimen-principle-urine-specimen-collection-02}\end{center}
\begin{itemize}
\item Routine urinalysis
\begin{itemize}
\item Macroscopic, physical examination
\item Chemical examination
\item Microscopic examination: colour, appearance, volume, specific gravity, odor, pH
\end{itemize}
\end{itemize}
\begin{center}\includegraphics[width=0.82\textwidth,height=0.68\textheight,keepaspectratio]{\reviewnotesimagepath d6-figure-principle-urine-specimen-microscopic-examination-colour-appearance.png}\par\captionof{figure}{The illustration relates to Microscopic examination: colour, appearance, volume, specific gravity, odor, pH, within the topic of Principle of Urine Specimen collection, Urinalysis, and Urine Culture.}\labfig{appendix-figure-d6-principle-urine-specimen-microscopic-examination-colour-appearance}\end{center}
\begin{itemize}
\item Appearance (clarity)
\begin{itemize}
\item Clear = no visible particulates, transparent
\item Hazy = slightly cloudy, few particulates, print easily seen through urine
\item Cloudy = many particulates, print blurred
\item Turbid = print not visible \(\rightarrow\) crystalized mineral salts and bacteria (from old urine), pathologies producing blood/pus or acute infection
\end{itemize}
\end{itemize}
\begin{center}\includegraphics[width=0.82\textwidth,height=0.68\textheight,keepaspectratio]{\reviewnotesimagepath d6-figure-principle-urine-specimen-turbid-print-not-visible.png}\par\captionof{figure}{The illustration relates to Turbid = print not visible \(\rightarrow\) crystalized mineral salts and bacteria (from old urine), pathologies producing blood/pus or acute infection, within the topic of Principle of Urine Specimen collection, Urinalysis, and Urine Culture.}\labfig{appendix-figure-d6-principle-urine-specimen-turbid-print-not-visible}\end{center}
\begin{itemize}
\item Urine specific gravity: measures urine concentration (wt of urine / wt of equivalent volume of water)
\end{itemize}
\begin{center}\includegraphics[width=0.82\textwidth,height=0.68\textheight,keepaspectratio]{\reviewnotesimagepath d6-figure-principle-urine-specimen-urine-specific-gravity-measures.png}\par\captionof{figure}{The illustration relates to Urine specific gravity: measures urine concentration (wt of urine / wt of equivalent volume of water), within the topic of Principle of Urine Specimen collection, Urinalysis, and Urine Culture.}\labfig{appendix-figure-d6-principle-urine-specimen-urine-specific-gravity-measures}\end{center}
\begin{center}\includegraphics[width=0.82\textwidth,height=0.68\textheight,keepaspectratio]{\reviewnotesimagepath d6-figure-principle-urine-specimen-principle-urine-specimen-collection-03.png}\par\captionof{figure}{The illustration relates to Principle of Urine Specimen collection, Urinalysis, and Urine Culture, within the topic of Principle of Urine Specimen collection, Urinalysis, and Urine Culture.}\labfig{appendix-figure-d6-principle-urine-specimen-principle-urine-specimen-collection-03}\end{center}
\begin{center}\includegraphics[width=0.82\textwidth,height=0.68\textheight,keepaspectratio]{\reviewnotesimagepath d6-figure-principle-urine-specimen-principle-urine-specimen-collection-04.png}\par\captionof{figure}{The illustration relates to Principle of Urine Specimen collection, Urinalysis, and Urine Culture, within the topic of Principle of Urine Specimen collection, Urinalysis, and Urine Culture.}\labfig{appendix-figure-d6-principle-urine-specimen-principle-urine-specimen-collection-04}\end{center}
\begin{itemize}
\item Normal urine pH = 5.0-6.0 (range 4.5-8.0) \(\rightarrow\) abnormal can develop renal stones
\begin{itemize}
\item Urine pH determined by dipstick, pH meter or pH paper
\item Acidic = xanthine, cystine uric acid, and calcium oxalate stones,
\item Alkaline = calcium carbonate, calcium phosphate, and magnesium phosphate stones
\end{itemize}
\end{itemize}
\begin{center}\includegraphics[width=0.82\textwidth,height=0.68\textheight,keepaspectratio]{\reviewnotesimagepath d6-figure-principle-urine-specimen-alkaline-calcium-carbonate-calcium.png}\par\captionof{figure}{The illustration relates to Alkaline = calcium carbonate, calcium phosphate, and magnesium phosphate stones, within the topic of Principle of Urine Specimen collection, Urinalysis, and Urine Culture.}\labfig{appendix-figure-d6-principle-urine-specimen-alkaline-calcium-carbonate-calcium}\end{center}
\begin{itemize}
\item Chemical examination of urine
\end{itemize}
\begin{center}\includegraphics[width=0.82\textwidth,height=0.68\textheight,keepaspectratio]{\reviewnotesimagepath d6-figure-principle-urine-specimen-chemical-examination-urine.png}\par\captionof{figure}{The illustration relates to Chemical examination of urine, within the topic of Principle of Urine Specimen collection, Urinalysis, and Urine Culture.}\labfig{appendix-figure-d6-principle-urine-specimen-chemical-examination-urine}\end{center}
\begin{center}\includegraphics[width=0.82\textwidth,height=0.68\textheight,keepaspectratio]{\reviewnotesimagepath d6-figure-principle-urine-specimen-principle-urine-specimen-collection-05.png}\par\captionof{figure}{The illustration relates to Principle of Urine Specimen collection, Urinalysis, and Urine Culture, within the topic of Principle of Urine Specimen collection, Urinalysis, and Urine Culture.}\labfig{appendix-figure-d6-principle-urine-specimen-principle-urine-specimen-collection-05}\end{center}
\begin{itemize}
\item Leukeocytes
\begin{itemize}
\item Normal value = 0-5 WBCs/HPF, negative
\item Measured by leukocyte esterase (component of most WBCs) + microscopic examination
\item Negative \(\rightarrow\) trace \(\rightarrow\) + to +++
\item Leukocyturia \(\rightarrow\) sign for inflammatory disease \(\rightarrow\) UTI, kidney inflammation
\end{itemize}
\item Nitrites
\begin{itemize}
\item Present in infection with nitrate reducing bacteria (Escherichia coli, Enterobacter, Klebsiella, Citrobacter and Proteus)
\item Positive urine test \(\rightarrow\) possible UTI (- does not mean no UTI)
\item To determine UTI also see leukocyte esterase, WBC count
\end{itemize}
\item Bilirubin/urobilinogen
\begin{itemize}
\item Product of hemoglobin degradation
\item Conjugated bilirubin = direct product of bilirubin breakdown \(\rightarrow\) not present in normal urine. Presence indicates liver dysfunction (hepatitis, cirrhosis, others) or biliary obstruction
\item Urobilinogen = bilirubin reabsorbed from intestines. Increase indicates early hepatitis, hemolytic jaundice, impaired liver function, hepatocellular jaundice
\end{itemize}
\item Protein
\begin{itemize}
\item Usual finding < 150 mg/24 hours or 10 mg/100 mL for any specimen
\item Proteinuria > 150 mg/ 24 hrs \(\rightarrow\) renal abnormality
\item Microalbuminuria = 30-150mg/24 hrs \(\rightarrow\) early renal disease, esp. in diabetic patients
\item Strip tests - sensitive to albumin, insensitive to low concentrations of: microglobulin, \(\gamma\)-globulins and other proteins
\item Range = -ve to ++++
\item Low molecular weight protein (< 40,000) \(\rightarrow\) typically pass through glomerular filtration and are reabsorbed (low concentration = total resorption)
\item Moderate Mw protein (albumin) \(\rightarrow\) high plasma concentration, but small amounts present in normal urine
\item High Mw proteins (> 90,000) \(\rightarrow\) unable to penetrate healthy glomerular filtration barrier
\end{itemize}
\item Urinary tract proteins (contd. of protein)
\begin{itemize}
\item Uromodulin (Tamm-Horsfall protein) \(\rightarrow\) mucoprotein synthesized by distal tubular cells \(\rightarrow\) involved in cast formation
\item Urokinase \(\rightarrow\) fibrinolytic enzyme secreted by tubular cells
\item Secretory immunoglobulin A \(\rightarrow\) synthesized by renal tubular epithelial cells
\end{itemize}
\end{itemize}
\begin{center}\includegraphics[width=0.82\textwidth,height=0.68\textheight,keepaspectratio]{\reviewnotesimagepath d6-figure-principle-urine-specimen-secretory-immunoglobulin-synthesized-renal.png}\par\captionof{figure}{The illustration relates to Secretory immunoglobulin A \(\rightarrow\) synthesized by renal tubular epithelial cells, within the topic of Principle of Urine Specimen collection, Urinalysis, and Urine Culture.}\labfig{appendix-figure-d6-principle-urine-specimen-secretory-immunoglobulin-synthesized-renal}\end{center}
\begin{center}\includegraphics[width=0.82\textwidth,height=0.68\textheight,keepaspectratio]{\reviewnotesimagepath d6-figure-principle-urine-specimen-principle-urine-specimen-collection-06.png}\par\captionof{figure}{The illustration relates to Principle of Urine Specimen collection, Urinalysis, and Urine Culture, within the topic of Principle of Urine Specimen collection, Urinalysis, and Urine Culture.}\labfig{appendix-figure-d6-principle-urine-specimen-principle-urine-specimen-collection-06}\end{center}
\begin{center}\includegraphics[width=0.82\textwidth,height=0.68\textheight,keepaspectratio]{\reviewnotesimagepath d6-figure-principle-urine-specimen-principle-urine-specimen-collection-07.png}\par\captionof{figure}{The illustration relates to Principle of Urine Specimen collection, Urinalysis, and Urine Culture, within the topic of Principle of Urine Specimen collection, Urinalysis, and Urine Culture.}\labfig{appendix-figure-d6-principle-urine-specimen-principle-urine-specimen-collection-07}\end{center}
\begin{itemize}
\item Glucose
\begin{itemize}
\item Filtered at glomerulus, reabsorbed at proximal tubule
\item Normal blood glucose = 70-100 mg/dL \(\rightarrow\) when high blood glucose occurs, glycosuria (glucose in urine) occurs
\item Glycosuria = diabetes, Cushing’s syndrome, liver/pancreatic disease, fanconi’s syndrome
\end{itemize}
\item Ketones
\begin{itemize}
\item Product of fat metabolism. Typically absent in urine
\item Ketones in urine: 20\% acetoacetic acid, 78\% beta-hydroxybutyric acid, and 2\% acetone
\item Nondiabetic ketonuria \(\rightarrow\) increased catabolism of adipose tissues under limited food intake
\item Diabetic ketonuria \(\rightarrow\) diabetic acidosis (ketosis) \(\rightarrow\) high probability of coma, other issues. Assoc. w/ uncontrolled diabetes
\end{itemize}
\item Blood
\begin{itemize}
\item Hematuria = presence of RBCs in urine \(\rightarrow\) indicates urinary tract bleeding, glomerular damage (visible on microscopic examination)
\item Hemoglobinuria = presence of dissolved hemoglobin in urine \(\rightarrow\) can indicate destruction of circulating WBCs
\item Myoglobinuria = red or brown urine \(\rightarrow\) results from rapid destruction of skeletal muscle
\item Urine strips screen for hemoglobin/myoglobin by catalyzation of peroxide and orthotolidine \(\rightarrow\) orthotilidine to blue-green colour
\end{itemize}
\end{itemize}
\begin{itemize}
\item Microscopic examination of urine
\begin{itemize}
\item Sample preparation
\begin{itemize}
\item Mix urine in centrifuge tube
\item Centrifugate at 1,500-3,000 RPM for 5 mins
\item Decant supernatant and resuspend sediment
\item Place sample on slide
\end{itemize}
\item Sediments in urine
\begin{itemize}
\item Cells: erythrocytes (RBCs), leukocytes, epithelial cells
\item Casts: hyaline, erythrocyte, leukocyte, epithelial, granular, waxy, fatty, broad casts
\item Crystals: calcium oxalate, uric, triple phosphate, etc.
\item Other structures (e.g. microorganisms)
\end{itemize}
\item Hematuria
\begin{itemize}
\item Defined as RBCs in urine
\item High count indicates infection, trauma, tumors or renal stones
\item Disfigured/dysmorphic/bizarre/acanthocytic RBCs indicate kidney source (glomerulonephritis) e.g. kidney inflammation. Caused by crossing renal basement membrane
\item Small amounts of RBCs are seen in young + healthy people (0-5 RBCs/HPF)
\end{itemize}
\end{itemize}
\end{itemize}
\begin{center}\includegraphics[width=0.82\textwidth,height=0.68\textheight,keepaspectratio]{\reviewnotesimagepath d6-figure-principle-urine-specimen-small-amounts-rbcs-seen.png}\par\captionof{figure}{The illustration relates to Small amounts of RBCs are seen in young + healthy people (0-5 RBCs/HPF), within the topic of Principle of Urine Specimen collection, Urinalysis, and Urine Culture.}\labfig{appendix-figure-d6-principle-urine-specimen-small-amounts-rbcs-seen}\end{center}
\begin{itemize}
\item Classification of dysmorphic RBCs
\begin{itemize}
\item D1 cells \(\rightarrow\) ringlike shape w/ severe cytoplasm colour loss + membranous protrusion/blebs
\item D2 cells \(\rightarrow\) doughnut like shape w/ moderate cytoplasm colour loss and membranous protrusions/blebs
\item D3 cells \(\rightarrow\) doughnut like shape w/ mild cytoplasm colour loss w/o membranous protrusions/blebs
\end{itemize}
\end{itemize}
\begin{center}\includegraphics[width=0.82\textwidth,height=0.68\textheight,keepaspectratio]{\reviewnotesimagepath d6-figure-principle-urine-specimen-d3-cells-doughnut-like.png}\par\captionof{figure}{The illustration relates to D3 cells \(\rightarrow\) doughnut like shape w/ mild cytoplasm colour loss w/o membranous protrusions/blebs, within the topic of Principle of Urine Specimen collection, Urinalysis, and Urine Culture.}\labfig{appendix-figure-d6-principle-urine-specimen-d3-cells-doughnut-like}\end{center}
\begin{center}\includegraphics[width=0.82\textwidth,height=0.68\textheight,keepaspectratio]{\reviewnotesimagepath d6-figure-principle-urine-specimen-principle-urine-specimen-collection-08.png}\par\captionof{figure}{The illustration relates to Principle of Urine Specimen collection, Urinalysis, and Urine Culture, within the topic of Principle of Urine Specimen collection, Urinalysis, and Urine Culture.}\labfig{appendix-figure-d6-principle-urine-specimen-principle-urine-specimen-collection-08}\end{center}
\begin{itemize}
\item White blood cells
\begin{itemize}
\item WBC/bacteria presence considered abnormal \(\rightarrow\) possible UTI (e.g. cystitis, pyelonephritis)
\item WBCs are recruited by chemotactic signalling from damaged tissues, migrate to interstitium through capillary basement membrane \(\rightarrow\) renal tubule epithelium \(\rightarrow\) urine
\item Normal finding = 0-5 WBCs/HPF. Abnormal = pyuria
\item Also called ‘glitter cell’ because of granule brownian motion
\item Red arrow = macrophage
\end{itemize}
\end{itemize}
\begin{itemize}
\item Epithelium cells
\begin{itemize}
\item From various parts of urinary tract
\item Increased pathological condition at site of origin
\item Types: squamous, renal tubular, transitional (bladder, caudate, columnar, polygonal), oval fat bodies
\end{itemize}
\end{itemize}
\begin{center}\includegraphics[width=0.82\textwidth,height=0.68\textheight,keepaspectratio]{\reviewnotesimagepath d6-figure-principle-urine-specimen-types-squamous-renal-tubular.png}\par\captionof{figure}{The illustration relates to Types: squamous, renal tubular, transitional (bladder, caudate, columnar, polygonal), oval fat bodies, within the topic of Principle of Urine Specimen collection, Urinalysis, and Urine Culture.}\labfig{appendix-figure-d6-principle-urine-specimen-types-squamous-renal-tubular}\end{center}
\begin{center}\includegraphics[width=0.82\textwidth,height=0.68\textheight,keepaspectratio]{\reviewnotesimagepath d6-figure-principle-urine-specimen-principle-urine-specimen-collection-09.png}\par\captionof{figure}{The illustration relates to Principle of Urine Specimen collection, Urinalysis, and Urine Culture, within the topic of Principle of Urine Specimen collection, Urinalysis, and Urine Culture.}\labfig{appendix-figure-d6-principle-urine-specimen-principle-urine-specimen-collection-09}\end{center}
\subsubsection{Casts}
\begin{itemize}
\item Proteinaceous product of renal tubules formed in DCT and collecting ducts passed into urine
\item Rarely found in normal physiology \(\rightarrow\) positive finding indicates pathology
\item Cylinder shaped with impression of kidney tubules
\end{itemize}
\begin{center}\includegraphics[width=0.82\textwidth,height=0.68\textheight,keepaspectratio]{\reviewnotesimagepath d6-figure-principle-urine-specimen-cylinder-shaped-impression-kidney.png}\par\captionof{figure}{The illustration relates to Cylinder shaped with impression of kidney tubules, within the topic of Principle of Urine Specimen collection, Urinalysis, and Urine Culture.}\labfig{appendix-figure-d6-principle-urine-specimen-cylinder-shaped-impression-kidney}\end{center}
\begin{center}\includegraphics[width=0.82\textwidth,height=0.68\textheight,keepaspectratio]{\reviewnotesimagepath d6-figure-principle-urine-specimen-principle-urine-specimen-collection-10.png}\par\captionof{figure}{The illustration relates to Principle of Urine Specimen collection, Urinalysis, and Urine Culture, within the topic of Principle of Urine Specimen collection, Urinalysis, and Urine Culture.}\labfig{appendix-figure-d6-principle-urine-specimen-principle-urine-specimen-collection-10}\end{center}
\begin{itemize}
\item Hyaline cast \& Cylindroids: non pathological cast
\begin{itemize}
\item Most common type = solid Tamm-horsfall mucoprotein w/o inclusions (basic composition for all casts)
\item Found in low urine flow, concentrated urine/acidic environment
\item Found in dehydration, vigorous exercise
\item No colour, homogeneous, semitransparent
\end{itemize}
\item Cellular cast
\begin{itemize}
\item RBC cast \(\rightarrow\) indicate glomerular damage
\item WBC cast \(\rightarrow\) indicate inflammation, infection
\item Epithelial cast or RET cast \(\rightarrow\) indicates nephrotic syndrome, acute tubular necrosis, toxin ingestion, cytomegalovirus, viral hepatitis
\item Fatty cast \& Oval fat body cast \(\rightarrow\) formed by breakdown of endothelial cells
\end{itemize}
\item Granular cast \(\rightarrow\) “hyaline casts with granules” resulting from breakdown of cellular casts or aggregates in plasma protein (e.g. albumin) or immunoglobulin light chains – cigar–shaped and highly refractive
\begin{itemize}
\item Coarsely granular cast - initial stage of incorporation of cellular materials into the cast
\item Finely granular cast - degraded coarsely granular cast
\end{itemize}
\item Waxy cast - homogenous resembling paraffin \(\rightarrow\) short, broad, irregular ends. Deteriorated granular casts \(\rightarrow\) local oliguria or anuria in CKD, nephritic syndrome, renal failure
\begin{itemize}
\item Hb cast - contain Hb from degraded RBCs
\item Waxy cast
\item Mixed cast
\end{itemize}
\item Telescopic urine sediment - RBC + oval fat cell + WBC and all types of cast found in a sediment
\end{itemize}
\subsubsection{Crystaluria}
\begin{itemize}
\item Urinary sediment with distinct forms for each condition
\item Precipitated from excessive concentration/solubility decrease by temperature or pH
\item Significant in large numbers \(\rightarrow\) may cause calculi
\end{itemize}
\begin{center}\includegraphics[width=0.82\textwidth,height=0.68\textheight,keepaspectratio]{\reviewnotesimagepath d6-figure-principle-urine-specimen-significant-large-numbers-cause.png}\par\captionof{figure}{The illustration relates to Significant in large numbers \(\rightarrow\) may cause calculi, within the topic of Principle of Urine Specimen collection, Urinalysis, and Urine Culture.}\labfig{appendix-figure-d6-principle-urine-specimen-significant-large-numbers-cause}\end{center}
\begin{center}\includegraphics[width=0.82\textwidth,height=0.68\textheight,keepaspectratio]{\reviewnotesimagepath d6-figure-principle-urine-specimen-principle-urine-specimen-collection-11.png}\par\captionof{figure}{The illustration relates to Principle of Urine Specimen collection, Urinalysis, and Urine Culture, within the topic of Principle of Urine Specimen collection, Urinalysis, and Urine Culture.}\labfig{appendix-figure-d6-principle-urine-specimen-principle-urine-specimen-collection-11}\end{center}
\begin{center}\includegraphics[width=0.82\textwidth,height=0.68\textheight,keepaspectratio]{\reviewnotesimagepath d6-figure-principle-urine-specimen-principle-urine-specimen-collection-12.png}\par\captionof{figure}{The illustration relates to Principle of Urine Specimen collection, Urinalysis, and Urine Culture, within the topic of Principle of Urine Specimen collection, Urinalysis, and Urine Culture.}\labfig{appendix-figure-d6-principle-urine-specimen-principle-urine-specimen-collection-12}\end{center}
\begin{center}\includegraphics[width=0.82\textwidth,height=0.68\textheight,keepaspectratio]{\reviewnotesimagepath d6-figure-principle-urine-specimen-principle-urine-specimen-collection-13.png}\par\captionof{figure}{The illustration relates to Principle of Urine Specimen collection, Urinalysis, and Urine Culture, within the topic of Principle of Urine Specimen collection, Urinalysis, and Urine Culture.}\labfig{appendix-figure-d6-principle-urine-specimen-principle-urine-specimen-collection-13}\end{center}
\begin{center}\includegraphics[width=0.82\textwidth,height=0.68\textheight,keepaspectratio]{\reviewnotesimagepath d6-figure-principle-urine-specimen-principle-urine-specimen-collection-14.png}\par\captionof{figure}{The illustration relates to Principle of Urine Specimen collection, Urinalysis, and Urine Culture, within the topic of Principle of Urine Specimen collection, Urinalysis, and Urine Culture.}\labfig{appendix-figure-d6-principle-urine-specimen-principle-urine-specimen-collection-14}\end{center}
\begin{itemize}
\item Natural \& Acid crystal
\begin{itemize}
\item Associated with disease
\begin{itemize}
\item Cystine - protein catabolism byproduct - may indicate cystinuria (metabolic)
\item Tyrosine - byproduct of protein catabolism. Uncommon, indicate liver damage
\item Leucine - byproduct of protein catabolism, indicate liver damage
\item Blilirubin - seen in canine urine
\item Cholesterol - in acidic urine, accompany chyluria (from abdominal/thoracic lymph obstruction) \(\rightarrow\) possibly UTI/nephritis
\end{itemize}
\item No association with disease or dietary crystals:
\begin{itemize}
\item Calcium oxalate - found in acidic urine. Not significant but may form calculi
\item Uric acid - in acidic urine. May indicate presence of calculi. Found in gout and physiologically
\end{itemize}
\end{itemize}
\item Basidic crystal
\begin{itemize}
\item Triple phosphate
\item Ammonium biurate
\item Ammonium urate
\item Calcium carbonate
\item Calcium phosphate
\end{itemize}
\item Drug crystal
\item Identified drug: Doxycycline, sulfa
\item Unclassified drug
\item Miscellaneous
\begin{itemize}
\item Mucous
\item Amorphous
\item Sperm
\item Contaminant
\item Parasite
\item Budding yeast
\item Pseudohyphae
\end{itemize}
\end{itemize}
\subsubsection{Procedures}
\begin{itemize}
\item Urinalsys
\end{itemize}
\begin{center}\includegraphics[width=0.82\textwidth,height=0.68\textheight,keepaspectratio]{\reviewnotesimagepath d6-figure-principle-urine-specimen-urinalsys.png}\par\captionof{figure}{The illustration relates to Urinalsys, within the topic of Principle of Urine Specimen collection, Urinalysis, and Urine Culture.}\labfig{appendix-figure-d6-principle-urine-specimen-urinalsys}\end{center}
\begin{itemize}
\item Urine culture
\end{itemize}
\begin{center}\includegraphics[width=0.82\textwidth,height=0.68\textheight,keepaspectratio]{\reviewnotesimagepath d6-figure-principle-urine-specimen-urine-culture.png}\par\captionof{figure}{The illustration relates to Urine culture, within the topic of Principle of Urine Specimen collection, Urinalysis, and Urine Culture.}\labfig{appendix-figure-d6-principle-urine-specimen-urine-culture}\end{center}
\begin{center}\includegraphics[width=0.82\textwidth,height=0.68\textheight,keepaspectratio]{\reviewnotesimagepath d6-figure-principle-urine-specimen-principle-urine-specimen-collection-15.png}\par\captionof{figure}{The illustration relates to Principle of Urine Specimen collection, Urinalysis, and Urine Culture, within the topic of Principle of Urine Specimen collection, Urinalysis, and Urine Culture.}\labfig{appendix-figure-d6-principle-urine-specimen-principle-urine-specimen-collection-15}\end{center}
\subsubsection{Interpretation of UA, Renal Diseases}
\begin{itemize}
\item Urinary tract infections
\begin{itemize}
\item Urethritis - urethral infection
\item Cystitis - infection at bladder/ureter
\begin{itemize}
\item Urine positive for transitional epithelium, WBC, bacteria and simple squamous epithelium +- blood and nitrite
\item No casts in urine, cloudy
\end{itemize}
\item Pylonephritis - infection at renal pelvis
\begin{itemize}
\item Symptoms: dis/oliguria, lumbar pain, fever \& chills, nausea \& vomiting
\item Urine strip: proteinuria, nitrite +, UC +
\item Microscopic: high WBC, bacteria, WBC casts, caudate cells
\end{itemize}
\end{itemize}
\item Urinary tract obstruction
\begin{itemize}
\item Renal calculi/urolithiasis - calculi in bladder \(\leftarrow\) most are calcium
\item Nephrolithiasis - renal calculi
\end{itemize}
\item Glomerular disorder
\begin{itemize}
\item Glomerular damage
\item Glomerulonephritis
\begin{itemize}
\item Laboratory: high creatinine, BUN
\item Urine strip: positive protein
\item UA,  micro: smoky urine, high RBC w/ dysmorphia, RBC/Hb cast
\item Chronic expresses more casts of RBC
\end{itemize}
\end{itemize}
\end{itemize}
\begin{center}\includegraphics[width=0.82\textwidth,height=0.68\textheight,keepaspectratio]{\reviewnotesimagepath d6-figure-principle-urine-specimen-chronic-expresses-more-casts.png}\par\captionof{figure}{The illustration relates to Chronic expresses more casts of RBC, within the topic of Principle of Urine Specimen collection, Urinalysis, and Urine Culture.}\labfig{appendix-figure-d6-principle-urine-specimen-chronic-expresses-more-casts}\end{center}
\begin{itemize}
\item Nephrotic syndrome
\begin{itemize}
\item Laboratory: high conc. level of BUN, Crea., low conc. level of serum protein, hyperlipidemia, high level K+
\item Urine strip +ve protein
\item UA, micro: hematuria \& smoky urine, lipiduria, high count of cell \& cast, prominent hyaline cast, RBC with dysmorphic, \textbf{oval fat bodies \& fatty cast, cholesterol crystal}
\end{itemize}
\item Renal failure
\item End stage of renal failure
\begin{itemize}
\item < 10\% renal function, GFR < 30
\item Lab: high BUN, high creatine
\item UA: little of urine, broad waxy cast, broad granular cast, high count of cells and cast, present several types of clinical significant sediment
\end{itemize}
\item Mechanisms of glomerular injury
\begin{itemize}
\item Antibody deposition
\item Complement activation
\item Glomerulosclerosis
\item Cellular proliferation
\end{itemize}
\item Renal tubular damage
\end{itemize}
\begin{center}\includegraphics[width=0.82\textwidth,height=0.68\textheight,keepaspectratio]{\reviewnotesimagepath d6-figure-principle-urine-specimen-renal-tubular-damage.png}\par\captionof{figure}{The illustration relates to Renal tubular damage, within the topic of Principle of Urine Specimen collection, Urinalysis, and Urine Culture.}\labfig{appendix-figure-d6-principle-urine-specimen-renal-tubular-damage}\end{center}
\begin{itemize}
\item Renal tubular necrosis
\begin{itemize}
\item Blood biochemistry: electrolytes- level of sodium and potassium is low, high creatinine and high blood urea nitrogen (BUN)
\item Urinalysis: diluted urine, maintain sp.gr. at 1.010
\item Microscopic finding: Renal tubular epithelial [RTE] \& renal tubular epithelial cast [RTE cast]
\end{itemize}
\item Association of other organ disorders
\begin{itemize}
\item Neurogenic dysfunction \(\rightarrow\) incontinence
\item Urine spermatozoa in male
\end{itemize}
\end{itemize}
\begin{center}\includegraphics[width=0.82\textwidth,height=0.68\textheight,keepaspectratio]{\reviewnotesimagepath d6-figure-principle-urine-specimen-urine-spermatozoa-male.png}\par\captionof{figure}{The illustration relates to Urine spermatozoa in male, within the topic of Principle of Urine Specimen collection, Urinalysis, and Urine Culture.}\labfig{appendix-figure-d6-principle-urine-specimen-urine-spermatozoa-male}\end{center}
\begin{itemize}
\item Urinary system cancer / metastasis cancer: bladder cancer, metastasis \(\rightarrow\) hematuria (nonspecific), cellular atypia
\item Genetic disorders
\begin{itemize}
\item Polycystic kidney disease
\begin{itemize}
\item Grapelike kidney cyst
\item Lumbar, abdominal pain, hypertension, hematuria
\item Nonspecific urine sediment (hematuria)
\end{itemize}
\end{itemize}
\item Accidental / Events / Trauma
\end{itemize}
\begin{center}\includegraphics[width=0.82\textwidth,height=0.68\textheight,keepaspectratio]{\reviewnotesimagepath d6-figure-principle-urine-specimen-accidental-events-trauma.png}\par\captionof{figure}{The illustration relates to Accidental / Events / Trauma, within the topic of Principle of Urine Specimen collection, Urinalysis, and Urine Culture.}\labfig{appendix-figure-d6-principle-urine-specimen-accidental-events-trauma}\end{center}

\subsection{Pathology of KUB System}
\begin{center}\includegraphics[width=0.82\textwidth,height=0.68\textheight,keepaspectratio]{\reviewnotesimagepath d6-figure-pathology-kub-pathology-kub-system.png}\par\captionof{figure}{The illustration relates to Pathology of KUB System, within the topic of Pathology of KUB System.}\labfig{appendix-figure-d6-pathology-kub-pathology-kub-system}\end{center}
\begin{center}\includegraphics[width=0.82\textwidth,height=0.68\textheight,keepaspectratio]{\reviewnotesimagepath d6-figure-pathology-kub-pathology-kub-system-02.png}\par\captionof{figure}{The illustration relates to Pathology of KUB System, within the topic of Pathology of KUB System.}\labfig{appendix-figure-d6-pathology-kub-pathology-kub-system-02}\end{center}
\begin{center}\includegraphics[width=0.82\textwidth,height=0.68\textheight,keepaspectratio]{\reviewnotesimagepath d6-figure-pathology-kub-pathology-kub-system-03.png}\par\captionof{figure}{The illustration relates to Pathology of KUB System, within the topic of Pathology of KUB System.}\labfig{appendix-figure-d6-pathology-kub-pathology-kub-system-03}\end{center}

\subsection{Normal Flora of Genitourinary System}
\begin{itemize}
\item LOs
\begin{itemize}
\item Urinary, reproductive tract structure (physical, chemical, physiological) and influence on microbiota community
\item Predominant normal flora of genitourinary tract
\item Advantage, disadvantage of genitourinary tract normal flora
\end{itemize}
\end{itemize}
\subsubsection{Normal Flora of the Urinary Tract}
\begin{itemize}
\item Urobiome = microbiota of the urinary tract
\item Decreases upstream of external orifice
\item Defense mechanisms of GU tract
\begin{itemize}
\item Urine properties, flow
\item Antiadherence factors (uromodulin, Tamm-Horsfall protein)
\item Antimicrobial factors (e.g. nitric oxide, cathelicidins, defensins)
\item Cell exfoliation
\item Chemokines, cytokines
\item Neutrophils
\end{itemize}
\item Normal flora of urethra
\begin{itemize}
\item Coagulase-negative staphylococci
\item Corynebacterium spp.
\item Viridans streptococci
\item Anaerobic bacteria
\item Lactobacilli
\item Mycoplasma \& Ureaplasma
\end{itemize}
\item Uropathogens
\begin{itemize}
\item E. coli and other Enterobacteriaceae
\item Staphylococcus saprophyticus
\end{itemize}
\item Male urethra
\begin{itemize}
\item Distal 6 cm colonized by microbes
\item Most microbiota = skin microbiota
\item Sexually active = more diversity of microbiota community
\end{itemize}
\end{itemize}
\begin{center}\includegraphics[width=0.82\textwidth,height=0.68\textheight,keepaspectratio]{\reviewnotesimagepath d6-figure-normal-flora-genitourinary-sexually-active-more-diversity.png}\par\captionof{figure}{The illustration relates to Sexually active = more diversity of microbiota community, within the topic of Normal Flora of Genitourinary System.}\labfig{appendix-figure-d6-normal-flora-genitourinary-sexually-active-more-diversity}\end{center}
\begin{center}\includegraphics[width=0.82\textwidth,height=0.68\textheight,keepaspectratio]{\reviewnotesimagepath d6-figure-normal-flora-genitourinary-normal-flora-genitourinary-system.png}\par\captionof{figure}{The illustration relates to Normal Flora of Genitourinary System, within the topic of Normal Flora of Genitourinary System.}\labfig{appendix-figure-d6-normal-flora-genitourinary-normal-flora-genitourinary-system}\end{center}
\begin{center}\includegraphics[width=0.82\textwidth,height=0.68\textheight,keepaspectratio]{\reviewnotesimagepath d6-figure-normal-flora-genitourinary-normal-flora-genitourinary-system-02.png}\par\captionof{figure}{The illustration relates to Normal Flora of Genitourinary System, within the topic of Normal Flora of Genitourinary System.}\labfig{appendix-figure-d6-normal-flora-genitourinary-normal-flora-genitourinary-system-02}\end{center}
\begin{center}\includegraphics[width=0.82\textwidth,height=0.68\textheight,keepaspectratio]{\reviewnotesimagepath d6-figure-normal-flora-genitourinary-normal-flora-genitourinary-system-03.png}\par\captionof{figure}{The illustration relates to Normal Flora of Genitourinary System, within the topic of Normal Flora of Genitourinary System.}\labfig{appendix-figure-d6-normal-flora-genitourinary-normal-flora-genitourinary-system-03}\end{center}
\begin{itemize}
\item Interpretation of urine culture
\begin{itemize}
\item > 10\textasciicircum{}5 CFU/mL = likely UTI, bacteriuria
\item > 10\textasciicircum{}4 - > 105 CFU/mL = contamination or early infection
\item < 10\textasciicircum{}4 CFU/mL = negative
\end{itemize}
\item Alteration of urethral microbiota
\begin{itemize}
\item Antibiotics
\item Catheterization
\item Sexual activities (in male)
\end{itemize}
\item Urinary dysbiosis may drive urinary pathology
\end{itemize}
\begin{center}\includegraphics[width=0.82\textwidth,height=0.68\textheight,keepaspectratio]{\reviewnotesimagepath d6-figure-normal-flora-genitourinary-urinary-dysbiosis-drive-urinary.png}\par\captionof{figure}{The illustration relates to Urinary dysbiosis may drive urinary pathology, within the topic of Normal Flora of Genitourinary System.}\labfig{appendix-figure-d6-normal-flora-genitourinary-urinary-dysbiosis-drive-urinary}\end{center}
\subsubsection{Normal Flora of the Reproductive Tract}
\begin{itemize}
\item Inherent defense mechanisms
\end{itemize}
\begin{center}\includegraphics[width=0.82\textwidth,height=0.68\textheight,keepaspectratio]{\reviewnotesimagepath d6-figure-normal-flora-genitourinary-inherent-defense-mechanisms.png}\par\captionof{figure}{The illustration relates to Inherent defense mechanisms, within the topic of Normal Flora of Genitourinary System.}\labfig{appendix-figure-d6-normal-flora-genitourinary-inherent-defense-mechanisms}\end{center}
\begin{itemize}
\item In females, periovulatory (close to, during, after) period or in menopause, antimicrobial characteristics change
\item Hormonal influences on microbiota
\end{itemize}
\begin{center}\includegraphics[width=0.82\textwidth,height=0.68\textheight,keepaspectratio]{\reviewnotesimagepath d6-figure-normal-flora-genitourinary-hormonal-influences-microbiota.png}\par\captionof{figure}{The illustration relates to Hormonal influences on microbiota, within the topic of Normal Flora of Genitourinary System.}\labfig{appendix-figure-d6-normal-flora-genitourinary-hormonal-influences-microbiota}\end{center}
\begin{center}\includegraphics[width=0.82\textwidth,height=0.68\textheight,keepaspectratio]{\reviewnotesimagepath d6-figure-normal-flora-genitourinary-normal-flora-genitourinary-system-04.png}\par\captionof{figure}{The illustration relates to Normal Flora of Genitourinary System, within the topic of Normal Flora of Genitourinary System.}\labfig{appendix-figure-d6-normal-flora-genitourinary-normal-flora-genitourinary-system-04}\end{center}
\begin{center}\includegraphics[width=0.82\textwidth,height=0.68\textheight,keepaspectratio]{\reviewnotesimagepath d6-figure-normal-flora-genitourinary-normal-flora-genitourinary-system-05.png}\par\captionof{figure}{The illustration relates to Normal Flora of Genitourinary System, within the topic of Normal Flora of Genitourinary System.}\labfig{appendix-figure-d6-normal-flora-genitourinary-normal-flora-genitourinary-system-05}\end{center}
\begin{center}\includegraphics[width=0.82\textwidth,height=0.68\textheight,keepaspectratio]{\reviewnotesimagepath d6-figure-normal-flora-genitourinary-normal-flora-genitourinary-system-06.png}\par\captionof{figure}{The illustration relates to Normal Flora of Genitourinary System, within the topic of Normal Flora of Genitourinary System.}\labfig{appendix-figure-d6-normal-flora-genitourinary-normal-flora-genitourinary-system-06}\end{center}
\begin{center}\includegraphics[width=0.82\textwidth,height=0.68\textheight,keepaspectratio]{\reviewnotesimagepath d6-figure-normal-flora-genitourinary-normal-flora-genitourinary-system-07.png}\par\captionof{figure}{The illustration relates to Normal Flora of Genitourinary System, within the topic of Normal Flora of Genitourinary System.}\labfig{appendix-figure-d6-normal-flora-genitourinary-normal-flora-genitourinary-system-07}\end{center}
\begin{itemize}
\item Main flora = \emph{Lactobacillus spp.} (acid producing) - found in vagina, probiotics
\item Other flora
\begin{itemize}
\item Virus: human papilloma virus (HPV)
\item Fungi: \emph{Candida albicans}
\item Protozoa: Pathogenic \emph{Trichomonas vaginalis} \(\rightarrow\) possibly asymptomatic
\end{itemize}
\end{itemize}
\begin{center}\includegraphics[width=0.82\textwidth,height=0.68\textheight,keepaspectratio]{\reviewnotesimagepath d6-figure-normal-flora-genitourinary-protozoa-pathogenic-trichomonas-vaginalis.png}\par\captionof{figure}{The illustration relates to Protozoa: Pathogenic Trichomonas vaginalis \(\rightarrow\) possibly asymptomatic, within the topic of Normal Flora of Genitourinary System.}\labfig{appendix-figure-d6-normal-flora-genitourinary-protozoa-pathogenic-trichomonas-vaginalis}\end{center}
\begin{itemize}
\item Advantages of normal microbiota
\begin{itemize}
\item Competitive exclusion
\item pH control
\item Antimicrobial activity
\end{itemize}
\end{itemize}
\begin{center}\includegraphics[width=0.82\textwidth,height=0.68\textheight,keepaspectratio]{\reviewnotesimagepath d6-figure-normal-flora-genitourinary-antimicrobial-activity.png}\par\captionof{figure}{The illustration relates to Antimicrobial activity, within the topic of Normal Flora of Genitourinary System.}\labfig{appendix-figure-d6-normal-flora-genitourinary-antimicrobial-activity}\end{center}
\begin{center}\includegraphics[width=0.82\textwidth,height=0.68\textheight,keepaspectratio]{\reviewnotesimagepath d6-figure-normal-flora-genitourinary-normal-flora-genitourinary-system-08.png}\par\captionof{figure}{The illustration relates to Normal Flora of Genitourinary System, within the topic of Normal Flora of Genitourinary System.}\labfig{appendix-figure-d6-normal-flora-genitourinary-normal-flora-genitourinary-system-08}\end{center}
\begin{itemize}
\item Antibiotics can change the vaginal microbiota
\end{itemize}
\subsubsection{Dysbiosis}
\begin{itemize}
\item Decreased normal flora \(\rightarrow\) overgrowth of yeast, bacteria
\item Risk factors: race/ethnicity, genetic, diet, gut microbiome, sexual behavior, hygiene, contraceptive, smoking, stress, obesity
\item Alterations of sexual microbiota are caused by
\begin{itemize}
\item Sexual maturity (i.e., whether pre- or post- menarchal, whether pre- or post-menopausal)
\item Pregnancy
\item Stage of the reproductive cycle in menarchal individuals
\item Sexual activity – type, frequency, etc.
\item Use of contraceptive devices/treatments
\item Use of tampons
\item Hormone replacement therapy in post-menopausal individuals
\item Use of vaginal douches
\item Antibiotic use
\end{itemize}
\item Restoration of healthy microbiome \(\rightarrow\) pre, probiotics, hygiene practices, dietary adjustments, lifestyle and reproductive health, medical interventions
\item In immunocompromise, C. albicans may become a pathogen to the host
\end{itemize}
\begin{center}\includegraphics[width=0.82\textwidth,height=0.68\textheight,keepaspectratio]{\reviewnotesimagepath d6-figure-normal-flora-genitourinary-immunocompromise-albicans-become-pathogen.png}\par\captionof{figure}{The illustration relates to In immunocompromise, C. albicans may become a pathogen to the host, within the topic of Normal Flora of Genitourinary System.}\labfig{appendix-figure-d6-normal-flora-genitourinary-immunocompromise-albicans-become-pathogen}\end{center}
\begin{center}\includegraphics[width=0.82\textwidth,height=0.68\textheight,keepaspectratio]{\reviewnotesimagepath d6-figure-normal-flora-genitourinary-normal-flora-genitourinary-system-09.png}\par\captionof{figure}{The illustration relates to Normal Flora of Genitourinary System, within the topic of Normal Flora of Genitourinary System.}\labfig{appendix-figure-d6-normal-flora-genitourinary-normal-flora-genitourinary-system-09}\end{center}
\begin{center}\includegraphics[width=0.82\textwidth,height=0.68\textheight,keepaspectratio]{\reviewnotesimagepath d6-figure-normal-flora-genitourinary-normal-flora-genitourinary-system-10.png}\par\captionof{figure}{The illustration relates to Normal Flora of Genitourinary System, within the topic of Normal Flora of Genitourinary System.}\labfig{appendix-figure-d6-normal-flora-genitourinary-normal-flora-genitourinary-system-10}\end{center}

\subsection{Medication in KUB System}
\subsubsection{Diuretics}
\begin{itemize}
\item Purpose \(\rightarrow\) increase urine volume
\item Major classifications: carbonic anhydrase inhibitor, loop diuretics, thiazide diuretics, potassium sparing diuretics (e.g. aldosterone antagonists), osmotic diuretics, vasopressin antagonists
\end{itemize}
\begin{center}\includegraphics[width=0.82\textwidth,height=0.68\textheight,keepaspectratio]{\reviewnotesimagepath d6-figure-medication-kub-major-classifications-carbonic-anhydrase.png}\par\captionof{figure}{The illustration relates to Major classifications: carbonic anhydrase inhibitor, loop diuretics, thiazide diuretics, potassium sparing diuretics (e.g. aldosterone antagonists), osmotic diu, within the topic of Medication in KUB System.}\labfig{appendix-figure-d6-medication-kub-major-classifications-carbonic-anhydrase}\end{center}
\begin{center}\includegraphics[width=0.82\textwidth,height=0.68\textheight,keepaspectratio]{\reviewnotesimagepath d6-figure-medication-kub-medication-kub-system.png}\par\captionof{figure}{The illustration relates to Medication in KUB System, within the topic of Medication in KUB System.}\labfig{appendix-figure-d6-medication-kub-medication-kub-system}\end{center}
\begin{center}\includegraphics[width=0.82\textwidth,height=0.68\textheight,keepaspectratio]{\reviewnotesimagepath d6-figure-medication-kub-medication-kub-system-02.png}\par\captionof{figure}{The illustration relates to Medication in KUB System, within the topic of Medication in KUB System.}\labfig{appendix-figure-d6-medication-kub-medication-kub-system-02}\end{center}
\paragraph{Proximal tubule diuretics}
\begin{itemize}
\item Carbonic anhydrase inhibitor (Acetazol\textbf{amide})
\begin{itemize}
\item Inhibits carbonic anhydrases (CO2 \(\Leftrightarrow\) HCO3-) \(\rightarrow\) decreased bicarbonate reabsorption \(\rightarrow\) decreased H+ generation \(\rightarrow\) less Na+ reabsorption by N-H exchanger 3
\item Sodium-Glucose transporter 2 \(\rightarrow\) glucose reabsorption still performed
\item Na+ not conserved \(\rightarrow\) natriuresis, decreased bicarbonate resorption at PCT
\item Clinical indications
\begin{itemize}
\item Glaucoma (reduce the ocular pressure by lowering Na+)
\item Urinary alkalization (for urine tract disease)
\item Metabolic alkalosis (e.g. loop/thiazide diuretics)
\item Acute mountain sickness (counter respiratory alkalosis)
\end{itemize}
\item Toxic/adverse effects
\begin{itemize}
\item Hyperchloremic metabolic acidosis
\item Renal stones (phosphate, calcium increase in urine by bicarbonaturic response)
\item Renal potassium wasting, increased Na delivered to collecting tubules
\end{itemize}
\item Drugs for glaucoma: Dichlorphenamide, Methazolamide
\end{itemize}
\item Sodium-glucose cotransporter 2 inhibitors
\begin{itemize}
\item Drug name: dapagli\textbf{flozin}, canagliflozin, empagliflozin,
\item ertugliflozin, and ipragliflozin
\item Clinical indications
\begin{itemize}
\item T2DM
\item CKD + proteinuria
\item Heart failure
\end{itemize}
\item Toxicity/adverse effects
\begin{itemize}
\item Genital fungal infection
\item Urinary tract infection
\item Minimal electrolyte effects
\end{itemize}
\end{itemize}
\end{itemize}
\paragraph{Loop Diuretics}
\begin{itemize}
\item Main transporter into interstitium = NA+/K+/2Cl- cotransporter type 2 (NKCC2) - K+ moves into lumen due to potential gradient
\item Inhibit NKCC luminal reabsorption of TAL, increase Mg2+ and Ca2+ excretion
\item Drug name: \textbf{Furosemide}, Bumetanide, Torsemide, Ethacrynic acid
\item Highest potency
\item Limited clearance by glomerulus (protein bound)
\item Clinical indications
\begin{itemize}
\item Acute pulmonary edema and other edematous conditions.
\item Hyperkalemia
\item Acute renal failure to increase urine flow
\item Anion overdose
\end{itemize}
\item Toxicity/adverse effects
\begin{itemize}
\item Hypokalemic metabolic alkalosis
\item Ototoxicity
\item Hyperuricemia
\item Hypomagnesemia
\item Allergic and their reaction: sulfonamide (except ethacrynic)
\end{itemize}
\item Relative potencies
\end{itemize}
\begin{center}\includegraphics[width=0.82\textwidth,height=0.68\textheight,keepaspectratio]{\reviewnotesimagepath d6-figure-medication-kub-relative-potencies.png}\par\captionof{figure}{The illustration relates to Relative potencies, within the topic of Medication in KUB System.}\labfig{appendix-figure-d6-medication-kub-relative-potencies}\end{center}
\paragraph{Thiazide Diuretics}
\begin{itemize}
\item Main transporter = Na+/Cl- cotransporter
\item Drug class affecting: thiazide diuretics (\textbf{Hydrochlorothiazide})
\item Action: block Na+/Cl- contransporter (NCC) \(\rightarrow\) inhibit reabsorption from luminal DCT
\item Clinical indications
\begin{itemize}
\item Hypertension
\item Heart failure
\item Nephrolithiasis due to idiopathic hypercalciuria
\item Nephrogenic DI
\end{itemize}
\item Toxicity
\begin{itemize}
\item Hypokalemic metabolic alkalosis
\item Impaired carbohydrate tolerance (hyperpolarization of beta cells, inhibiting insulin release)
\item Hyperlipidemia
\item Hyponatremia (hypovolemia increase ADH, impair diluting capacity of the kidney, increased thirst)
\item Impaired uric acid metabolism and gout
\item Allergic reaction ‘sulfonamide group’
\end{itemize}
\item Drugs in the class (1 = non-thiazide, but similar action)
\end{itemize}
\begin{center}\includegraphics[width=0.82\textwidth,height=0.68\textheight,keepaspectratio]{\reviewnotesimagepath d6-figure-medication-kub-drugs-class-non-thiazide.png}\par\captionof{figure}{The illustration relates to Drugs in the class (1 = non-thiazide, but similar action), within the topic of Medication in KUB System.}\labfig{appendix-figure-d6-medication-kub-drugs-class-non-thiazide}\end{center}
\paragraph{Collecting Duct Diuretics}
\begin{itemize}
\item K+ sparing
\begin{itemize}
\item Main channel = ENaC (epithelium Na channel) \(\leftarrow\) aldosterone sensitive resorption of Na+ in collecting duct
\item Drug classes: mineralocorticoid inhibitor (\textbf{spironolactone}, eplerenone), ENaC inhibitor (\textbf{amiloride}, triamterene)
\item Prevent K+ secretion from principal cell of collecting duct
\item Spironolactone = competitive inhibitor of aldosterone \(\rightarrow\) can lead to gynecomastia, decreased libido
\item Direct EnAC inhibitors: prevent Na+ reabsorption in collecting tubules
\item Indications:
\begin{itemize}
\item Mineralocorticoid excess or hyperaldosteronism
\item Secondary hyperaldosteronism by heart failure, hepatic cirrhosis, nephrotic syndrome, condition that diminished effective intravascular volume
\item Amiloride: Liddle syndrome (autosomal dominant, activation of ENaC), spironolactone (lack efficacy)
\end{itemize}
\item Adverse effects
\begin{itemize}
\item Hypercalcemia (especially if: kidney disease, using RAAS blocker drug)
\item Hyperchloremic metabolic acidosis
\item Type IV renal tubular acidosis (inhibited H+, K+ secretion)
\item Gynecomastia
\item Acute renal failure (triamterene, indomethacin)
\item Kidney stones (triamterene)
\end{itemize}
\item K+ sparing diuretic list (1 = approved only in HT)
\end{itemize}
\end{itemize}
\begin{center}\includegraphics[width=0.82\textwidth,height=0.68\textheight,keepaspectratio]{\reviewnotesimagepath d6-figure-medication-kub-sparing-diuretic-list-approved.png}\par\captionof{figure}{The illustration relates to K+ sparing diuretic list (1 = approved only in HT), within the topic of Medication in KUB System.}\labfig{appendix-figure-d6-medication-kub-sparing-diuretic-list-approved}\end{center}
\begin{itemize}
\item Vasopressin antagonists
\begin{itemize}
\item Vasopressin receptor mediated expression of aquaporin-2
\item Main drug class: vasopressin receptor antagonist (Tol\textbf{vaptan}, Conivaptan)
\item Indication
\begin{itemize}
\item Syndrome of inappropriate ADH
\item Other causes of elevated ADH: heart failure
\item Autosomal dominant polycystic kidney disease
\end{itemize}
\item Toxicity / adverse effect
\begin{itemize}
\item Nephrogenic diabetes insipidus
\item Renal failure
\item Hypotension
\end{itemize}
\end{itemize}
\end{itemize}
\paragraph{Osmotic Diuretics}
\begin{itemize}
\item \textbf{Mannitol }\(\rightarrow\) action on proximal tubule, descending limb of loop, oppose action of ADH in collecting duct
\item Clinical indication: reduce intracranial, intraocular pressure
\item Toxicity
\begin{itemize}
\item Extracellular volume expansion
\item Dehydration, hyperkalemia, hypernatremia
\item Hyponatremia
\item Acute renal failure
\end{itemize}
\end{itemize}
\paragraph{Effects of Diuretics}
\begin{itemize}
\item pH and electrolytes
\end{itemize}
\begin{center}\includegraphics[width=0.82\textwidth,height=0.68\textheight,keepaspectratio]{\reviewnotesimagepath d6-figure-medication-kub-ph-electrolytes.png}\par\captionof{figure}{The illustration relates to pH and electrolytes, within the topic of Medication in KUB System.}\labfig{appendix-figure-d6-medication-kub-ph-electrolytes}\end{center}
\begin{itemize}
\item Renal excretion
\end{itemize}
\begin{center}\includegraphics[width=0.82\textwidth,height=0.68\textheight,keepaspectratio]{\reviewnotesimagepath d6-figure-medication-kub-renal-excretion.png}\par\captionof{figure}{The illustration relates to Renal excretion, within the topic of Medication in KUB System.}\labfig{appendix-figure-d6-medication-kub-renal-excretion}\end{center}
\subsubsection{IV Fluids}
\begin{center}\includegraphics[width=0.82\textwidth,height=0.68\textheight,keepaspectratio]{\reviewnotesimagepath d6-figure-medication-kub-iv-fluids.png}\par\captionof{figure}{The illustration relates to IV Fluids, within the topic of Medication in KUB System.}\labfig{appendix-figure-d6-medication-kub-iv-fluids}\end{center}
\begin{center}\includegraphics[width=0.82\textwidth,height=0.68\textheight,keepaspectratio]{\reviewnotesimagepath d6-figure-medication-kub-medication-kub-system-03.png}\par\captionof{figure}{The illustration relates to Medication in KUB System, within the topic of Medication in KUB System.}\labfig{appendix-figure-d6-medication-kub-medication-kub-system-03}\end{center}
\begin{itemize}
\item Osmolality (mOsm/kg H2O) = 2 [Na+] + glucose /18 + BUN /2.8
\item Effective osmolality (mOsm/kg H2O) = 2 [Na+] + glucose /18
\item Normal osmolarity range 280-295 mOsm/kg H2O
\end{itemize}
\begin{center}\includegraphics[width=0.82\textwidth,height=0.68\textheight,keepaspectratio]{\reviewnotesimagepath d6-figure-medication-kub-normal-osmolarity-range-280.png}\par\captionof{figure}{The illustration relates to Normal osmolarity range 280-295 mOsm/kg H2O, within the topic of Medication in KUB System.}\labfig{appendix-figure-d6-medication-kub-normal-osmolarity-range-280}\end{center}
\begin{center}\includegraphics[width=0.82\textwidth,height=0.68\textheight,keepaspectratio]{\reviewnotesimagepath d6-figure-medication-kub-medication-kub-system-04.png}\par\captionof{figure}{The illustration relates to Medication in KUB System, within the topic of Medication in KUB System.}\labfig{appendix-figure-d6-medication-kub-medication-kub-system-04}\end{center}
\subsubsection{Chronic Kidney Disease}
\begin{itemize}
\item CKD = abnormal kidney structure/function > 3 months persistence \(\rightarrow\) usually irreversible, measured by reduction in GFR
\end{itemize}
\begin{center}\includegraphics[width=0.82\textwidth,height=0.68\textheight,keepaspectratio]{\reviewnotesimagepath d6-figure-medication-kub-ckd-abnormal-kidney-structure.png}\par\captionof{figure}{The illustration relates to CKD = abnormal kidney structure/function > 3 months persistence \(\rightarrow\) usually irreversible, measured by reduction in GFR, within the topic of Medication in KUB System.}\labfig{appendix-figure-d6-medication-kub-ckd-abnormal-kidney-structure}\end{center}
\begin{itemize}
\item CKD complications:
\begin{itemize}
\item Fluid, electrolyte, acid-base imbalances (edema, hyperkalemia K > 5.5 mEq/L, metabolic acidosis)
\item Anemia (Hb < 13 g/dL male, 12 female) \(\leftarrow\) iron supplement, erythropoiesis stimulators
\end{itemize}
\end{itemize}
\begin{center}\includegraphics[width=0.82\textwidth,height=0.68\textheight,keepaspectratio]{\reviewnotesimagepath d6-figure-medication-kub-anemia-hb-13-dl.png}\par\captionof{figure}{The illustration relates to Anemia (Hb < 13 g/dL male, 12 female) \(\leftarrow\) iron supplement, erythropoiesis stimulators, within the topic of Medication in KUB System.}\labfig{appendix-figure-d6-medication-kub-anemia-hb-13-dl}\end{center}
\begin{itemize}
\item Mineral bone disease (hyperphosphatemia, hypocalcemia, vit D deficiency)
\end{itemize}
\begin{center}\includegraphics[width=0.82\textwidth,height=0.68\textheight,keepaspectratio]{\reviewnotesimagepath d6-figure-medication-kub-mineral-bone-disease-hyperphosphatemia.png}\par\captionof{figure}{The illustration relates to Mineral bone disease (hyperphosphatemia, hypocalcemia, vit D deficiency), within the topic of Medication in KUB System.}\labfig{appendix-figure-d6-medication-kub-mineral-bone-disease-hyperphosphatemia}\end{center}
\begin{itemize}
\item Uremia
\item Treatment
\begin{itemize}
\item Reduction in dietary potassium, review medications causing hyperkalemia
\item Medication
\end{itemize}
\item Medication
\end{itemize}
\begin{center}\includegraphics[width=0.82\textwidth,height=0.68\textheight,keepaspectratio]{\reviewnotesimagepath d6-figure-medication-kub-medication.png}\par\captionof{figure}{The illustration relates to Medication, within the topic of Medication in KUB System.}\labfig{appendix-figure-d6-medication-kub-medication}\end{center}
\begin{itemize}
\item Oral phosphate binder
\end{itemize}
\begin{center}\includegraphics[width=0.82\textwidth,height=0.68\textheight,keepaspectratio]{\reviewnotesimagepath d6-figure-medication-kub-oral-phosphate-binder.png}\par\captionof{figure}{The illustration relates to Oral phosphate binder, within the topic of Medication in KUB System.}\labfig{appendix-figure-d6-medication-kub-oral-phosphate-binder}\end{center}
\begin{center}\includegraphics[width=0.82\textwidth,height=0.68\textheight,keepaspectratio]{\reviewnotesimagepath d6-figure-medication-kub-medication-kub-system-05.png}\par\captionof{figure}{The illustration relates to Medication in KUB System, within the topic of Medication in KUB System.}\labfig{appendix-figure-d6-medication-kub-medication-kub-system-05}\end{center}
\begin{center}\includegraphics[width=0.82\textwidth,height=0.68\textheight,keepaspectratio]{\reviewnotesimagepath d6-figure-medication-kub-medication-kub-system-06.png}\par\captionof{figure}{The illustration relates to Medication in KUB System, within the topic of Medication in KUB System.}\labfig{appendix-figure-d6-medication-kub-medication-kub-system-06}\end{center}
\begin{itemize}
\item Alkalizing agents
\end{itemize}
\begin{center}\includegraphics[width=0.82\textwidth,height=0.68\textheight,keepaspectratio]{\reviewnotesimagepath d6-figure-medication-kub-alkalizing-agents.png}\par\captionof{figure}{The illustration relates to Alkalizing agents, within the topic of Medication in KUB System.}\labfig{appendix-figure-d6-medication-kub-alkalizing-agents}\end{center}
\begin{center}\includegraphics[width=0.82\textwidth,height=0.68\textheight,keepaspectratio]{\reviewnotesimagepath d6-figure-medication-kub-medication-kub-system-07.png}\par\captionof{figure}{The illustration relates to Medication in KUB System, within the topic of Medication in KUB System.}\labfig{appendix-figure-d6-medication-kub-medication-kub-system-07}\end{center}

\section{Reproductive Anatomy and Histology}

\subsection{Histology of Male Reproductive System}
LOs:

\begin{itemize}
\item Explain the histology of the structures in testes
\item Explain the spermatogenesis
\item Identify the characteristics of the excurrent duct system
\item Discriminate the histology of accessory glands
\item Identify the histology of penis
\end{itemize}
\begin{center}\includegraphics[width=0.82\textwidth,height=0.68\textheight,keepaspectratio]{\reviewnotesimagepath d6-figure-histology-male-reproductive-identify-histology-penis.png}\par\captionof{figure}{The illustration relates to Identify the histology of penis, within the topic of Histology of Male Reproductive System.}\labfig{appendix-figure-d6-histology-male-reproductive-identify-histology-penis}\end{center}
\begin{center}\includegraphics[width=0.82\textwidth,height=0.68\textheight,keepaspectratio]{\reviewnotesimagepath d6-figure-histology-male-reproductive-histology-male-reproductive-system.png}\par\captionof{figure}{The illustration relates to Histology of Male Reproductive System, within the topic of Histology of Male Reproductive System.}\labfig{appendix-figure-d6-histology-male-reproductive-histology-male-reproductive-system}\end{center}
\subsubsection{Components of Testes}
\begin{itemize}
\item Seminiferous tubules
\item Sertoli cells - “supporting cells” for sperm responding to FSH
\item Spermatogonia (basal) \(\rightarrow\) spermatocyte \(\rightarrow\) spermatid \(\rightarrow\) sperm
\item Spermatogenesis \(\rightarrow\) spermiogenesis (+ motility)
\end{itemize}
\begin{center}\includegraphics[width=0.82\textwidth,height=0.68\textheight,keepaspectratio]{\reviewnotesimagepath d6-figure-histology-male-reproductive-spermatogenesis-spermiogenesis-motility.png}\par\captionof{figure}{The illustration relates to Spermatogenesis \(\rightarrow\) spermiogenesis (+ motility), within the topic of Histology of Male Reproductive System.}\labfig{appendix-figure-d6-histology-male-reproductive-spermatogenesis-spermiogenesis-motility}\end{center}
\begin{center}\includegraphics[width=0.82\textwidth,height=0.68\textheight,keepaspectratio]{\reviewnotesimagepath d6-figure-histology-male-reproductive-histology-male-reproductive-system-02.png}\par\captionof{figure}{The illustration relates to Histology of Male Reproductive System, within the topic of Histology of Male Reproductive System.}\labfig{appendix-figure-d6-histology-male-reproductive-histology-male-reproductive-system-02}\end{center}
\subsubsection{Vas Deferens}
\begin{itemize}
\item Later stage: longitudinal \(\rightarrow\) circular \(\rightarrow\) longitudinal muscle layers
\end{itemize}

\subsection{Histology of Female Reproductive System}
LOs

\begin{itemize}
\item Explain the histology of ovary
\begin{itemize}
\item Ogenesis
\end{itemize}
\item Explain the histology of the genital tract
\begin{itemize}
\item Uterine tube
\item Uterus
\item Vagina
\end{itemize}
\item Explain the histology of the breasts
\item Primordial germ cells \(\rightarrow\) oogonia enter the long prophase of a first meiotic division \(\rightarrow\) arrested in prophase I: called “primary oocyte” and surrounded by follicular cells
\item Primary oocyte + follicular cells = primordial follicle
\item Primary \(\rightarrow\) secondary \(\rightarrow\) graafian \(\rightarrow\) rupture to corpus luteum + oocyte
\begin{itemize}
\item Key figures to determine the staging: oocyte, follicular cells, surrounding stroma, location
\end{itemize}
\end{itemize}
\begin{center}\includegraphics[width=0.82\textwidth,height=0.68\textheight,keepaspectratio]{\reviewnotesimagepath d6-figure-histology-female-reproductive-key-figures-determine-staging.png}\par\captionof{figure}{The illustration relates to Key figures to determine the staging: oocyte, follicular cells, surrounding stroma, location, within the topic of Histology of Female Reproductive System.}\labfig{appendix-figure-d6-histology-female-reproductive-key-figures-determine-staging}\end{center}
\begin{center}\includegraphics[width=0.82\textwidth,height=0.68\textheight,keepaspectratio]{\reviewnotesimagepath d6-figure-histology-female-reproductive-histology-female-reproductive-system.png}\par\captionof{figure}{The illustration relates to Histology of Female Reproductive System, within the topic of Histology of Female Reproductive System.}\labfig{appendix-figure-d6-histology-female-reproductive-histology-female-reproductive-system}\end{center}
\begin{center}\includegraphics[width=0.82\textwidth,height=0.68\textheight,keepaspectratio]{\reviewnotesimagepath d6-figure-histology-female-reproductive-histology-female-reproductive-system-02.png}\par\captionof{figure}{The illustration relates to Histology of Female Reproductive System, within the topic of Histology of Female Reproductive System.}\labfig{appendix-figure-d6-histology-female-reproductive-histology-female-reproductive-system-02}\end{center}

\section{Reproductive Physiology and Urology}

\subsection{Male Reproductive Physiology}
\begin{itemize}
\item LOs:
\begin{itemize}
\item Structure, function of testes (including key components: seminiferous tubules, sertoli cells, leydig cells)
\item Spermatogenesis process and influencing factors (+ hormones)
\item Role of hormones in male reproductive system (hypothalamic-pituitary axis: GnRH, LH, FSH, testosterone)
\item Endocrine function of testes, contribution to physical, psychological development
\item Age-related hormonal changes (decreased testosterone) and health impact (andropause, testosterone deficiency)
\end{itemize}
\end{itemize}
\subsubsection{Male reproductive system overview}
\begin{itemize}
\item Consists of primary (testes) + accessory (seminal vesicles, prostate gland, urethra, penis) sex organs
\item External genitalia = scrotum, penis, urethra (remainder is internal)
\item Functions of male reproductive system
\begin{itemize}
\item Spermatogenesis
\item Maintenance of reproductive tract (sperm transport, semen production)
\item Development of secondary sexual characteristics (body hair growth, voice deepening, libido)
\end{itemize}
\end{itemize}
\subsubsection{Functional anatomy of testes}
\begin{itemize}
\item Covering layers (anterior, lateral surface covered by all 3 layers, posterior surface covered by T albuginea only)
\begin{itemize}
\item Tunica vasculosa = inner vascular layer for providing blood to testes
\item Tunica albuginea = dense connective tissue \(\rightarrow\) provides structure, divides testes into lobules
\item Tunica vaginalis = outermost layer from peritoneum \(\rightarrow\) protection, reduces friction
\end{itemize}
\item Parenchyma (functional parts)
\begin{itemize}
\item Lobules of testes: seminiferous tubules, site of spermatogenesis \(\rightarrow\) rete testis (collect sperm from sem tubules) \(\rightarrow\) efferent ductules \(\rightarrow\) epididymis (stores sperm, maturation) \(\rightarrow\) vas deferens \(\rightarrow\) ejaculatory duct
\item Leydig cells (adjacent to seminiferous tubules): produce testosterone (male sexual characteristics, sperm production)
\end{itemize}
\end{itemize}
\begin{itemize}
\item Seminiferus tubule (contd.)
\end{itemize}
\begin{center}\includegraphics[width=0.82\textwidth,height=0.68\textheight,keepaspectratio]{\reviewnotesimagepath d6-figure-male-reproductive-physiology-seminiferus-tubule-contd.png}\par\captionof{figure}{The illustration relates to Seminiferus tubule (contd.), within the topic of Male Reproductive Physiology.}\labfig{appendix-figure-d6-male-reproductive-physiology-seminiferus-tubule-contd}\end{center}
\begin{itemize}
\item Thread–like convoluted tubular structures producing spermatozoa/sperms. 400-600 seminiferous tubules in each testis
\item Each tubule = 30-70 cm long, diameter 150-300 µm
\item Sertoli cells: support, nourish developing sperm cells. Produce inhibin to regulate FSH secretion (and inhibit spermatogonia production)
\item Spermatogenic cells: become mature sperm cells
\item Interstitial stroma: space between seminiferous tubules
\begin{itemize}
\item Leydig cells: testosterone–producing
\item Blood vessels
\end{itemize}
\item Lamina propria (CT surrounding seminiferous tubules)
\begin{itemize}
\item Myoid cells: contractile cells, move sperm and fluid
\end{itemize}
\end{itemize}
\begin{center}\includegraphics[width=0.82\textwidth,height=0.68\textheight,keepaspectratio]{\reviewnotesimagepath d6-figure-male-reproductive-physiology-myoid-cells-contractile-cells.png}\par\captionof{figure}{The illustration relates to Myoid cells: contractile cells, move sperm and fluid, within the topic of Male Reproductive Physiology.}\labfig{appendix-figure-d6-male-reproductive-physiology-myoid-cells-contractile-cells}\end{center}
\begin{center}\includegraphics[width=0.82\textwidth,height=0.68\textheight,keepaspectratio]{\reviewnotesimagepath d6-figure-male-reproductive-physiology-male-reproductive-physiology.png}\par\captionof{figure}{The illustration relates to Male Reproductive Physiology, within the topic of Male Reproductive Physiology.}\labfig{appendix-figure-d6-male-reproductive-physiology-male-reproductive-physiology}\end{center}
\begin{center}\includegraphics[width=0.82\textwidth,height=0.68\textheight,keepaspectratio]{\reviewnotesimagepath d6-figure-male-reproductive-physiology-male-reproductive-physiology-03.png}\par\captionof{figure}{The illustration relates to Male Reproductive Physiology, within the topic of Male Reproductive Physiology.}\labfig{appendix-figure-d6-male-reproductive-physiology-male-reproductive-physiology-03}\end{center}
\begin{itemize}
\item Spermatogenic cells
\begin{itemize}
\item Aka germ cells
\item Precursor cells to spermatozoa. Located between sertoli cells
\item Arranged in an orderly manner, 4-8 layers thick
\item Spermatogonia (primitive stage) in children \(\rightarrow\) puberty hormones trigger spermatogenesis to mature spermatozoa
\item Stages of cell: spermatogonium \(\rightarrow\) primary spermatocyte \(\rightarrow\) secondary spermatocyte \(\rightarrow\) spermatid
\end{itemize}
\item Sertoli cells
\begin{itemize}
\item Support development of spermatogenic cells in seminiferous tubules
\item Aka sustentacular/nurse cells
\item Structure = large, tall, irregular columnar cells from basement membrane to lumen of seminiferous tubules
\item Germ cells attach to sertoli cells by cytoplasmic connection \(\rightarrow\) germ cells released at maturity into lumen
\item Functions
\begin{itemize}
\item Support, nourish spermatogenic cells
\item Maintain blood–testis barrier
\item Hormonal regulation \(\rightarrow\) produce inhibin on response to FSH \(\rightarrow\) regulate sperm production
\item Secretion of androgen–binding protein: maintain high local testosterone levels to aid spermatogenesis
\item Embryo: production of Müllerian inhibiting factor \(\rightarrow\) regression of Müllerian duct in male fetus, allow male reproductive development
\item Phagocytosis of excess cytoplasm in spermatogenesis
\end{itemize}
\end{itemize}
\item The blood–testis barrier
\begin{itemize}
\item Protect developing sperm cells from immune attack
\item Selectively permeable for essential substances (hormones, nutrients), blocking other large molecules
\item Regulation of microenvironment (provide optimal conditions for spermatogenesis
\item Prevent antigenic exposure (sperm proteins entering circulation)
\item Barrier damage
\begin{itemize}
\item Causes: trauma, infection (e.g. mumps), radiation, toxins
\item Effect: sperm enter bloodstream \(\rightarrow\) immune response, autoantibodies \(\rightarrow\) infertile/sterile
\end{itemize}
\end{itemize}
\end{itemize}
\subsubsection{Functions of Testes}
\begin{itemize}
\item Gametogenesis (spermatogenesis)
\begin{itemize}
\item Proliferation \(\rightarrow\) growth \(\rightarrow\) maturation \(\rightarrow\) transformation
\end{itemize}
\end{itemize}
\begin{center}\includegraphics[width=0.82\textwidth,height=0.68\textheight,keepaspectratio]{\reviewnotesimagepath d6-figure-male-reproductive-physiology-proliferation-growth-maturation-transformation.png}\par\captionof{figure}{The illustration relates to Proliferation \(\rightarrow\) growth \(\rightarrow\) maturation \(\rightarrow\) transformation, within the topic of Male Reproductive Physiology.}\labfig{appendix-figure-d6-male-reproductive-physiology-proliferation-growth-maturation-transformation}\end{center}
\begin{center}\includegraphics[width=0.82\textwidth,height=0.68\textheight,keepaspectratio]{\reviewnotesimagepath d6-figure-male-reproductive-physiology-male-reproductive-physiology-04.png}\par\captionof{figure}{The illustration relates to Male Reproductive Physiology, within the topic of Male Reproductive Physiology.}\labfig{appendix-figure-d6-male-reproductive-physiology-male-reproductive-physiology-04}\end{center}
\begin{itemize}
\item Process of spermatogenesis
\begin{itemize}
\item Controls: FSH, testosterone (increase spermatogonia proliferation, increase sertoli androgen binding protein)
\item Proliferation of spermatogonia (stem cells) \(\rightarrow\) meiosis (spermatocytes \(\rightarrow\) spermatids) \(\rightarrow\) spermiogenesis: mature, develop into sperm \(\rightarrow\) spermiation: release of mature sperm
\item Major influences
\begin{itemize}
\item Sertoli cells (see above)
\item Hormones, including: follicle–stimulating hormone, growth hormone, testosterone, inhibin + activin, estrogen, luteinizing hormone
\end{itemize}
\end{itemize}
\end{itemize}
\begin{center}\includegraphics[width=0.82\textwidth,height=0.68\textheight,keepaspectratio]{\reviewnotesimagepath d6-figure-male-reproductive-physiology-hormones-including-follicle-stimulating.png}\par\captionof{figure}{The illustration relates to Hormones, including: follicle–stimulating hormone, growth hormone, testosterone, inhibin + activin, estrogen, luteinizing hormone, within the topic of Male Reproductive Physiology.}\labfig{appendix-figure-d6-male-reproductive-physiology-hormones-including-follicle-stimulating}\end{center}
\begin{center}\includegraphics[width=0.82\textwidth,height=0.68\textheight,keepaspectratio]{\reviewnotesimagepath d6-figure-male-reproductive-physiology-male-reproductive-physiology-05.png}\par\captionof{figure}{The illustration relates to Male Reproductive Physiology, within the topic of Male Reproductive Physiology.}\labfig{appendix-figure-d6-male-reproductive-physiology-male-reproductive-physiology-05}\end{center}
\begin{center}\includegraphics[width=0.82\textwidth,height=0.68\textheight,keepaspectratio]{\reviewnotesimagepath d6-figure-male-reproductive-physiology-male-reproductive-physiology-06.png}\par\captionof{figure}{The illustration relates to Male Reproductive Physiology, within the topic of Male Reproductive Physiology.}\labfig{appendix-figure-d6-male-reproductive-physiology-male-reproductive-physiology-06}\end{center}
\begin{itemize}
\item Temperature \(\rightarrow\) increased temperature prevents spermatogenesis (common in cryptorchidism). Testes in scrotum are approx 2C cooler than body temperature; temperature maintained by pampiniform plexus of blood vessels and cremasteric muscle (move testes away/toward body)
\end{itemize}
\begin{center}\includegraphics[width=0.82\textwidth,height=0.68\textheight,keepaspectratio]{\reviewnotesimagepath d6-figure-male-reproductive-physiology-temperature-increased-temperature-prevents.png}\par\captionof{figure}{The illustration relates to Temperature \(\rightarrow\) increased temperature prevents spermatogenesis (common in cryptorchidism). Testes in scrotum are approx 2C cooler than body temperature; temperatu, within the topic of Male Reproductive Physiology.}\labfig{appendix-figure-d6-male-reproductive-physiology-temperature-increased-temperature-prevents}\end{center}
\begin{itemize}
\item Nutrition
\item Genetic, environmental factors
\item Hormonal secretions by testis
\begin{itemize}
\item Testosterone (large quantities)
\begin{itemize}
\item Major androgen of leydig cells
\item Synthesis increased by LH (by stimulation of cholesterol desmolase in the first pathway step)
\end{itemize}
\item Dihydrotestosterone (more active form)
\begin{itemize}
\item Converted form of testosterone. Occurs in accessory sex organs (e.g. prostate)
\item 5 alpha reductase inhibitors (like finasteride) can treat benign prostatic hyperplasia by inhibiting testosterone conversion
\end{itemize}
\item Androstendione
\item Activin, inhibin (no androgenic effect, regulatory)
\item Female sex hormones (estrogen, progesterone)
\end{itemize}
\item Testosterone (contd.)
\begin{itemize}
\item Normal secretion rate (adult male) 13.9–31.33 µmol/dL, small amounts in female from ovary and some from adrenal gland
\item 98\% bound form—inactive (65\% β-globulin in gonadal/sex steroid binding globulin, 33\% to albumin), 2\% free form
\item Normal plasma testosterone
\begin{itemize}
\item 300-1,000 ng/dL (10.4-34.7 nmol/L) in adult male
\item 30–70 ng/dL (1.04–2.43 nmol/L) in adult female
\end{itemize}
\item Regulation of testosterone secretion
\begin{itemize}
\item Hypothalamus \(\rightarrow\) LH releasing hormone \(\rightarrow\) anterior pituitary \(\rightarrow\) LH
\item LH or interstitial cell stimulating hormone in leydig cell \(\rightarrow\) proportional testosterone secretion
\item Negative feedback (testosterone inhibits LH releasing hormone secretion in hypothalamus)
\end{itemize}
\end{itemize}
\end{itemize}
\begin{center}\includegraphics[width=0.82\textwidth,height=0.68\textheight,keepaspectratio]{\reviewnotesimagepath d6-figure-male-reproductive-physiology-negative-feedback-testosterone-inhibits.png}\par\captionof{figure}{The illustration relates to Negative feedback (testosterone inhibits LH releasing hormone secretion in hypothalamus), within the topic of Male Reproductive Physiology.}\labfig{appendix-figure-d6-male-reproductive-physiology-negative-feedback-testosterone-inhibits}\end{center}
\begin{center}\includegraphics[width=0.82\textwidth,height=0.68\textheight,keepaspectratio]{\reviewnotesimagepath d6-figure-male-reproductive-physiology-male-reproductive-physiology-07.png}\par\captionof{figure}{The illustration relates to Male Reproductive Physiology, within the topic of Male Reproductive Physiology.}\labfig{appendix-figure-d6-male-reproductive-physiology-male-reproductive-physiology-07}\end{center}
\begin{itemize}
\item Functions of testosterone
\begin{itemize}
\item Fetus: sex differentiation, development of accessory sexual organs, descent of testis
\item Adult: sex organ effects, including increasing size of penis, scrotum, testes (8x between puberty onset and age 20)
\end{itemize}
\item Secondary sexual characteristics determined by testosterone
\end{itemize}
\begin{center}\includegraphics[width=0.82\textwidth,height=0.68\textheight,keepaspectratio]{\reviewnotesimagepath d6-figure-male-reproductive-physiology-secondary-sexual-characteristics-determined.png}\par\captionof{figure}{The illustration relates to Secondary sexual characteristics determined by testosterone, within the topic of Male Reproductive Physiology.}\labfig{appendix-figure-d6-male-reproductive-physiology-secondary-sexual-characteristics-determined}\end{center}
\begin{center}\includegraphics[width=0.82\textwidth,height=0.68\textheight,keepaspectratio]{\reviewnotesimagepath d6-figure-male-reproductive-physiology-male-reproductive-physiology-08.png}\par\captionof{figure}{The illustration relates to Male Reproductive Physiology, within the topic of Male Reproductive Physiology.}\labfig{appendix-figure-d6-male-reproductive-physiology-male-reproductive-physiology-08}\end{center}
\begin{itemize}
\item Androgen (contd.)
\begin{itemize}
\item Androgen is secreted by interstitial cells of leydig in testes
\item Leydig cells are found in infants and adult males: few are present in childhood
\item Secretion of testosterone occurs mainly in newborns and following puberty
\end{itemize}
\end{itemize}
\begin{itemize}
\item Androgens are also secreted by zona reticularis in adrenal cortex. However, no significant physiological effect (small quantities)
\begin{itemize}
\item Testosterone
\item Androstenedione
\item Dehydroepiandrosterone
\end{itemize}
\item Presence of female sex hormones in males
\begin{itemize}
\item Estrogen typically 12-34 pg/mL
\item Sources of estrogen in males
\begin{itemize}
\item Adrenal cortex (small quantity)
\item Testes (20\%) \(\rightarrow\) formed from androgens in sertoli cells by aromatase enzyme
\item Other organs (esp. liver) - 80\% of circulating levels
\end{itemize}
\item Progesterone in males typically 0.3 ng/mL
\end{itemize}
\end{itemize}
\subsubsection{Semen and Sperm}
\begin{center}\includegraphics[width=0.82\textwidth,height=0.68\textheight,keepaspectratio]{\reviewnotesimagepath d6-figure-male-reproductive-physiology-semen-sperm.png}\par\captionof{figure}{The illustration relates to Semen and Sperm, within the topic of Male Reproductive Physiology.}\labfig{appendix-figure-d6-male-reproductive-physiology-semen-sperm}\end{center}
\begin{center}\includegraphics[width=0.82\textwidth,height=0.68\textheight,keepaspectratio]{\reviewnotesimagepath d6-figure-male-reproductive-physiology-male-reproductive-physiology-09.png}\par\captionof{figure}{The illustration relates to Male Reproductive Physiology, within the topic of Male Reproductive Physiology.}\labfig{appendix-figure-d6-male-reproductive-physiology-male-reproductive-physiology-09}\end{center}
\begin{itemize}
\item Capacitation process
\begin{itemize}
\item prepares sperm for fertilization (increases motility, ability to penetrate zona pellucida)
\item Occurs in female reproductive tract (uterus, oviducts)
\item Key changes: removal of seminal plasma proteins coating sperm, increased membrane fluidity (better interaction with oocyte), hyperactivation of sperm motility
\item Promote acrosome reaction (for fertilization)
\end{itemize}
\end{itemize}
\begin{center}\includegraphics[width=0.82\textwidth,height=0.68\textheight,keepaspectratio]{\reviewnotesimagepath d6-figure-male-reproductive-physiology-promote-acrosome-reaction-fertilization.png}\par\captionof{figure}{The illustration relates to Promote acrosome reaction (for fertilization), within the topic of Male Reproductive Physiology.}\labfig{appendix-figure-d6-male-reproductive-physiology-promote-acrosome-reaction-fertilization}\end{center}
\begin{itemize}
\item Age–related hormonal changes in males
\end{itemize}
\begin{center}\includegraphics[width=0.82\textwidth,height=0.68\textheight,keepaspectratio]{\reviewnotesimagepath d6-figure-male-reproductive-physiology-age-related-hormonal-changes.png}\par\captionof{figure}{The illustration relates to Age–related hormonal changes in males, within the topic of Male Reproductive Physiology.}\labfig{appendix-figure-d6-male-reproductive-physiology-age-related-hormonal-changes}\end{center}
\begin{itemize}
\item Fetal life - development of internal/external structures, regression of Müllerian ducts
\begin{itemize}
\item hCG (Human Chorionic Gonadotropin) + LH \(\rightarrow\) Leydig cells \(\rightarrow\) testosterone \(\rightarrow\) Wolffian duct + 5-alpha reductase
\item Testosterone \(\rightarrow\) DHT \(\rightarrow\) urogenital sinus, genital organ formation (genitual tubercle \(\rightarrow\) scrotum, penis, prostate)
\item Sertoli cells \(\rightarrow\) Müllerian inhibiting factor \(\rightarrow\) regression of female structures
\item Testosterone presence \(\rightarrow\) development of Wolffian ducts \(\rightarrow\) male internal structures (epididymis, vas deferens, seminal vesicles)
\end{itemize}
\end{itemize}
\begin{center}\includegraphics[width=0.82\textwidth,height=0.68\textheight,keepaspectratio]{\reviewnotesimagepath d6-figure-male-reproductive-physiology-testosterone-presence-development-wolffian.png}\par\captionof{figure}{The illustration relates to Testosterone presence \(\rightarrow\) development of Wolffian ducts \(\rightarrow\) male internal structures (epididymis, vas deferens, seminal vesicles), within the topic of Male Reproductive Physiology.}\labfig{appendix-figure-d6-male-reproductive-physiology-testosterone-presence-development-wolffian}\end{center}
\begin{itemize}
\item Childhood \(\rightarrow\) following birth, LH, testosterone drop, remain low until puberty
\item Puberty \(\rightarrow\) LH pulse amplitude increases \(\rightarrow\) increase in mean LH (especially during sleep) \(\rightarrow\) leydig cell stimulation, testosterone secretion
\item Adult \(\rightarrow\) LH \(\rightarrow\) testosterone = LH and testosterone parallel each other
\item Aging adult \(\rightarrow\) testosterone, inhibin levels decrease gradually with age (age 70, approximately 60-70\% testosterone secretion) \(\rightarrow\) increased LH, FSH
\end{itemize}

\subsection{Female Reproductive Physiology}
\begin{itemize}
\item LO:
\begin{itemize}
\item Female reproductive organs (primary, accessory sex organs)
\item Sexual life in females (first, second third periods)
\item Ovary (follicles, hormones)
\item Menstrual cycle
\item Physiological changes of hormones during pregnancy
\item Hormonal control of mammary glands
\item Age-related hormonal changes in females
\end{itemize}
\item Female Reproductive Organs
\end{itemize}
\begin{itemize}
\item Primary sex organs: ovaries \(\rightarrow\) produce eggs/ova, secrete female sex hormones (estrogen, progesterone)
\item Accessory sex organs: genital ducts (fallopian tubes, uterus, cervix, vagina) and external genitalia (labia majora/minora, clitoris)
\end{itemize}
\begin{center}\includegraphics[width=0.82\textwidth,height=0.68\textheight,keepaspectratio]{\reviewnotesimagepath d6-figure-female-reproductive-physiology-accessory-sex-organs-genital.png}\par\captionof{figure}{The illustration relates to Accessory sex organs: genital ducts (fallopian tubes, uterus, cervix, vagina) and external genitalia (labia majora/minora, clitoris), within the topic of Female Reproductive Physiology.}\labfig{appendix-figure-d6-female-reproductive-physiology-accessory-sex-organs-genital}\end{center}
\begin{center}\includegraphics[width=0.82\textwidth,height=0.68\textheight,keepaspectratio]{\reviewnotesimagepath d6-figure-female-reproductive-physiology-female-reproductive-physiology.png}\par\captionof{figure}{The illustration relates to Female Reproductive Physiology, within the topic of Female Reproductive Physiology.}\labfig{appendix-figure-d6-female-reproductive-physiology-female-reproductive-physiology}\end{center}
\subsubsection{The Ovary}
\begin{itemize}
\item Two main functions: gametogenesis, endocrine
\item Medulla/zona vasculosa
\begin{itemize}
\item Deeper portion of ovary w/ stroma of loose CT
\item Contains blood vessels, lymphatics, nerve fibers and smooth muscle fibers
\end{itemize}
\item Cortex
\begin{itemize}
\item Outer, broader portion of ovary
\item Consists of compact cellular layers interrupted at hilum (medulla continuous to mesovarium)
\item Lined with germinal epithelium with tunica albuginea underneath
\item Possesses structures
\begin{itemize}
\item Glandular structures: ovarian follicles at different stages)
\item Connective tissue cells
\item Interstitial cells (cluster of epithelial cells w/ fine lipid granules) formed from theca interna \(\rightarrow\) produce androgens
\end{itemize}
\end{itemize}
\end{itemize}
\begin{center}\includegraphics[width=0.82\textwidth,height=0.68\textheight,keepaspectratio]{\reviewnotesimagepath d6-figure-female-reproductive-physiology-interstitial-cells-cluster-epithelial.png}\par\captionof{figure}{The illustration relates to Interstitial cells (cluster of epithelial cells w/ fine lipid granules) formed from theca interna \(\rightarrow\) produce androgens, within the topic of Female Reproductive Physiology.}\labfig{appendix-figure-d6-female-reproductive-physiology-interstitial-cells-cluster-epithelial}\end{center}
\begin{itemize}
\item Ovarian follicles
\begin{itemize}
\item In utero, the outer cortex contains germinal epithelium (from germinal ridges) \(\rightarrow\) gives rise to primordial ova
\item Primordial ova move towards inner cortex \(\rightarrow\) granulosa cells from ovary stroma surround ova
\item Primordial ovum + granulosa cells = primordial follicle
\item 6 million primordial follicles (month 7-8) \(\rightarrow\) 1M at birth \(\rightarrow\) 300-400K at puberty \(\rightarrow\) 450 dominant follicles ovulate
\end{itemize}
\end{itemize}
\begin{center}\includegraphics[width=0.82\textwidth,height=0.68\textheight,keepaspectratio]{\reviewnotesimagepath d6-figure-female-reproductive-physiology-million-primordial-follicles-month.png}\par\captionof{figure}{The illustration relates to 6 million primordial follicles (month 7-8) \(\rightarrow\) 1M at birth \(\rightarrow\) 300-400K at puberty \(\rightarrow\) 450 dominant follicles ovulate, within the topic of Female Reproductive Physiology.}\labfig{appendix-figure-d6-female-reproductive-physiology-million-primordial-follicles-month}\end{center}
\begin{center}\includegraphics[width=0.82\textwidth,height=0.68\textheight,keepaspectratio]{\reviewnotesimagepath d6-figure-female-reproductive-physiology-female-reproductive-physiology-02.png}\par\captionof{figure}{The illustration relates to Female Reproductive Physiology, within the topic of Female Reproductive Physiology.}\labfig{appendix-figure-d6-female-reproductive-physiology-female-reproductive-physiology-02}\end{center}
\begin{center}\includegraphics[width=0.82\textwidth,height=0.68\textheight,keepaspectratio]{\reviewnotesimagepath d6-figure-female-reproductive-physiology-female-reproductive-physiology-04.png}\par\captionof{figure}{The illustration relates to Female Reproductive Physiology, within the topic of Female Reproductive Physiology.}\labfig{appendix-figure-d6-female-reproductive-physiology-female-reproductive-physiology-04}\end{center}
\begin{itemize}
\item Ovarian hormones
\begin{itemize}
\item Ovary secretes estrogen (C18 steroid), progesterone (C21 steroid), inhibin, relaxin, androgens (small quantities)
\item Estrogen = derived from androgens: androstenedione from theca interna \(\rightarrow\) migrate to granulosa cells \(\rightarrow\) estrogen conversion by aromatase
\item In non–pregnant state
\begin{itemize}
\item Hormones secreted largely by theca interna of ovarian follicles
\item Some secretion by ovarian corpus luteum
\item Predominant secretion at the follicular phase before ovulation
\end{itemize}
\item In pregnancy
\begin{itemize}
\item Estrogen secreted by placenta
\end{itemize}
\end{itemize}
\end{itemize}
\begin{center}\includegraphics[width=0.82\textwidth,height=0.68\textheight,keepaspectratio]{\reviewnotesimagepath d6-figure-female-reproductive-physiology-estrogen-secreted-placenta.png}\par\captionof{figure}{The illustration relates to Estrogen secreted by placenta, within the topic of Female Reproductive Physiology.}\labfig{appendix-figure-d6-female-reproductive-physiology-estrogen-secreted-placenta}\end{center}
\begin{itemize}
\item Estrogen (contd.)
\begin{itemize}
\item Present in 3 forms in plasma \(\rightarrow\) E2:E1:E3 = 10:5:1
\begin{itemize}
\item β-estradiol (E2)
\item Estrone (E1)
\item Estriol (E3)
\end{itemize}
\item Estrogen half life = 30-60 mins
\item Rate of estradiol secretion depends on phase of menstrual cycle
\begin{itemize}
\item 36 µg/dL peak before ovulation. \textbf{[Accuracy check: Verify the unit; the source may have confused pg/mL with µg/dL.] }
\item 250 µg/dL during midluteal phase. \textbf{[Accuracy check: Verify the unit; the source may have confused pg/mL with µg/dL.] }
\item Decline during menopause
\item Men about 50 µg/dL (184 nmol/dL). \textbf{[Accuracy check: Verify both concentration units; the values appear inconsistent with one another and may reflect a unit transcription error.] }
\end{itemize}
\item Transport of estradiol \(\rightarrow\) 98\% bound form (38\% bound to β-globulin called gonadal steroid binding globulin/sex steroid binding globulin, 60\% bound to albumin), 2\% free form
\end{itemize}
\end{itemize}
\begin{center}\includegraphics[width=0.82\textwidth,height=0.68\textheight,keepaspectratio]{\reviewnotesimagepath d6-figure-female-reproductive-physiology-transport-estradiol-98-bound.png}\par\captionof{figure}{The illustration relates to Transport of estradiol \(\rightarrow\) 98\% bound form (38\% bound to β-globulin called gonadal steroid binding globulin/sex steroid binding globulin, 60\% bound to albumin), 2\%, within the topic of Female Reproductive Physiology.}\labfig{appendix-figure-d6-female-reproductive-physiology-transport-estradiol-98-bound}\end{center}
\begin{itemize}
\item Functions of estrogen
\end{itemize}
\begin{center}\includegraphics[width=0.82\textwidth,height=0.68\textheight,keepaspectratio]{\reviewnotesimagepath d6-figure-female-reproductive-physiology-functions-estrogen.png}\par\captionof{figure}{The illustration relates to Functions of estrogen, within the topic of Female Reproductive Physiology.}\labfig{appendix-figure-d6-female-reproductive-physiology-functions-estrogen}\end{center}
\begin{center}\includegraphics[width=0.82\textwidth,height=0.68\textheight,keepaspectratio]{\reviewnotesimagepath d6-figure-female-reproductive-physiology-female-reproductive-physiology-06.png}\par\captionof{figure}{The illustration relates to Female Reproductive Physiology, within the topic of Female Reproductive Physiology.}\labfig{appendix-figure-d6-female-reproductive-physiology-female-reproductive-physiology-06}\end{center}
\begin{center}\includegraphics[width=0.82\textwidth,height=0.68\textheight,keepaspectratio]{\reviewnotesimagepath d6-figure-female-reproductive-physiology-female-reproductive-physiology-07.png}\par\captionof{figure}{The illustration relates to Female Reproductive Physiology, within the topic of Female Reproductive Physiology.}\labfig{appendix-figure-d6-female-reproductive-physiology-female-reproductive-physiology-07}\end{center}
\begin{center}\includegraphics[width=0.82\textwidth,height=0.68\textheight,keepaspectratio]{\reviewnotesimagepath d6-figure-female-reproductive-physiology-female-reproductive-physiology-08.png}\par\captionof{figure}{The illustration relates to Female Reproductive Physiology, within the topic of Female Reproductive Physiology.}\labfig{appendix-figure-d6-female-reproductive-physiology-female-reproductive-physiology-08}\end{center}
\begin{itemize}
\item Progesterone (cont.)
\begin{itemize}
\item Sources
\begin{itemize}
\item In non-pregnant: small quantities by theca interna cells during follicular phase, large quantities by corpus luteum during secretory phase, small quantity by corpus luteum
\item Pregnant: large quantities by corpus luteum (1st trimester), remainder secreted by placenta
\end{itemize}
\item Transport in blood
\begin{itemize}
\item 98\% bound form (18\% bound to corticosteroid binding globulin, 80\% to albumin)
\item 2\% free form
\end{itemize}
\item Half life = 4-5 minutes
\item Rate of secretion
\begin{itemize}
\item Men = 0.3 ng/mL (1 nmol/L)
\item Women = 0.9 ng/mL (3 mmol/L) during follicular phase, increased in late follicle phase
\end{itemize}
\item Pattern of secretion in women
\begin{itemize}
\item Theca cells produce pregnenolone \(\rightarrow\) granulosa cells convert to progesterone
\item Late follicular phase = increased secretion of progesterone
\item Luteal phase, production by corpus luteum \(\rightarrow\) large quantities, marked increase to 18 ng/mL (60 nmol/L)
\end{itemize}
\item Functions of progesterone
\end{itemize}
\end{itemize}
\begin{center}\includegraphics[width=0.82\textwidth,height=0.68\textheight,keepaspectratio]{\reviewnotesimagepath d6-figure-female-reproductive-physiology-functions-progesterone.png}\par\captionof{figure}{The illustration relates to Functions of progesterone, within the topic of Female Reproductive Physiology.}\labfig{appendix-figure-d6-female-reproductive-physiology-functions-progesterone}\end{center}
\begin{center}\includegraphics[width=0.82\textwidth,height=0.68\textheight,keepaspectratio]{\reviewnotesimagepath d6-figure-female-reproductive-physiology-female-reproductive-physiology-09.png}\par\captionof{figure}{The illustration relates to Female Reproductive Physiology, within the topic of Female Reproductive Physiology.}\labfig{appendix-figure-d6-female-reproductive-physiology-female-reproductive-physiology-09}\end{center}
\subsubsection{Normal Menstrual cycle}
\begin{itemize}
\item Ovarian activity cycles repeating approximately every month
\begin{itemize}
\item Mean of 28 days, but only 15\% of women exactly 28 days
\item Range of normal = 21-35 days with 3-8 day menses duration
\item Ovulation occurs on day 14 (36 hours after onset of mid-cycle LH surge)
\item Typical volume loss 30 mL (normal range 20-70 mL)
\item Abnormal frequency
\begin{itemize}
\item Frequent menstruation = polymenorrhoea (< 21 days)
\item Infrequent menstruation = oligomenorrhoea (> 35 days)
\item Absent = amenorrhoea (either primary or secondary)
\end{itemize}
\item Abnormal duration
\begin{itemize}
\item Prolonged = menorrhagia (> 7 days)
\item Shortened = hypomenorrhoea (< 3 days)
\end{itemize}
\item Abnormal volume
\begin{itemize}
\item Heavy = hypermenorrhoea (> 70 mL)
\item Light = hypomenorrhoea (< 20 mL)
\end{itemize}
\end{itemize}
\item First menstruation = menarche. Usually occurs between 12 - 15
\item Mean age of menopause = 51.4
\item Phases of menstrual cycle
\begin{itemize}
\item First day of cycle = first day of menses
\item Ovarian cycle: follicular, ovulation, luteal, menses
\end{itemize}
\end{itemize}
\begin{center}\includegraphics[width=0.82\textwidth,height=0.68\textheight,keepaspectratio]{\reviewnotesimagepath d6-figure-female-reproductive-physiology-ovarian-cycle-follicular-ovulation.png}\par\captionof{figure}{The illustration relates to Ovarian cycle: follicular, ovulation, luteal, menses, within the topic of Female Reproductive Physiology.}\labfig{appendix-figure-d6-female-reproductive-physiology-ovarian-cycle-follicular-ovulation}\end{center}
\begin{itemize}
\item Follicular/proliferative/preovulatory phase
\begin{itemize}
\item Day 1-14 \(\rightarrow\) estrogen dominates \(\rightarrow\) renewal of endometrium
\item FSH slightly elevated = granulosa cells proliferate, increased estrogen secretion in follicles \(\rightarrow\) 20 primordial follicles activate
\item One follicle (dominant follicle) secretes most estradiol, promoting growth and FSH sensitivity \(\rightarrow\) the dominant follicle continues development (others lack sufficient FSH \(\rightarrow\) only contribute androgen \(\rightarrow\) become atretic/die)
\end{itemize}
\item Ovulation
\begin{itemize}
\item Approximately day 14
\item LH surge induces ovulation
\item Always 14 days before end of cycle (fixed time). \textbf{[Accuracy check: Avoid presenting luteal-phase timing as invariant; verify how cycle variability is described in the source reference.] }
\item Estradiol synthesis burst (at end of follicular phase) \(\rightarrow\) positive feedback for FSH, LH secretion \(\rightarrow\) LH surge \(\rightarrow\) ovulation
\item LH > FSH because inhibin prevents FSH release by anterior pituitary
\item Estrogen levels decrease after ovulation \(\rightarrow\) rise again at luteal phase
\item Cervical mucus increases, viscosity lowers and increases sperm permeability
\end{itemize}
\end{itemize}
\begin{center}\includegraphics[width=0.82\textwidth,height=0.68\textheight,keepaspectratio]{\reviewnotesimagepath d6-figure-female-reproductive-physiology-cervical-mucus-increases-viscosity.png}\par\captionof{figure}{The illustration relates to Cervical mucus increases, viscosity lowers and increases sperm permeability, within the topic of Female Reproductive Physiology.}\labfig{appendix-figure-d6-female-reproductive-physiology-cervical-mucus-increases-viscosity}\end{center}
\begin{itemize}
\item Luteal phase
\begin{itemize}
\item approximately 2 weeks from day 14-28)
\item Dominated by progesterone, lower estrogen \(\rightarrow\) quinescent phase for endometrium, prepare uterus for implantation
\item Ruptured graafian follicle \(\rightarrow\) ovum released \(\rightarrow\) follicle filled with blood (slow clotting) \(\rightarrow\) corpus hemorrhagicum \(\rightarrow\) corpus luteum \(\rightarrow\) cavity closes, heals gradually \(\rightarrow\) blood replaced with serous fluid w/fibrin \(\rightarrow\) corpus luteum diameter 15 mm
\item Corpus luteum/yellow body: glandular yellow body from ruptured graafian follicle after release of ovum \(\rightarrow\) contributes to hormone release \(\rightarrow\) degenerates after end of cycle
\end{itemize}
\end{itemize}
\begin{center}\includegraphics[width=0.82\textwidth,height=0.68\textheight,keepaspectratio]{\reviewnotesimagepath d6-figure-female-reproductive-physiology-corpus-luteum-yellow-body.png}\par\captionof{figure}{The illustration relates to Corpus luteum/yellow body: glandular yellow body from ruptured graafian follicle after release of ovum \(\rightarrow\) contributes to hormone release \(\rightarrow\) degenerates after end, within the topic of Female Reproductive Physiology.}\labfig{appendix-figure-d6-female-reproductive-physiology-corpus-luteum-yellow-body}\end{center}
\begin{itemize}
\item Corpus luteum–morphology (contd.)
\begin{itemize}
\item In corpus luteum, luteinization occurs: granulosa cells, theca interna cells \(\rightarrow\) granulosa/theca lutein cells
\item Granulosa lutein cells: have fine lipid granules, yellowish pigment granules (pigment \(\rightarrow\) corpus luteum colour)
\item Theca lutein cells: only lipid granules
\item Follicular cavity of corpus luteum = reduced size, irregular outline, filled with serous fluid + blood clots
\item Without fertilization, corpus luteum reaches max size at 1 week \(\rightarrow\) degenerates to corpus luteum menstrualis/spurium (smaller, involuted/shriveled cell) \(\rightarrow\) corpus albicans (by luteolysis process)
\item With fertilization \(\rightarrow\) corpus luteum increases size (20-30 mm diameter) \(\rightarrow\) becomes true corpus luteum of pregnancy/corpus luteum graviditatis (verum) \(\rightarrow\) help maintain pregnancy \(\rightarrow\) degrades after 3-4 months
\end{itemize}
\item Corpus luteum–function (contd.)
\begin{itemize}
\item Secretion of endocrine hormone, help maintain pregnancy
\item Major secretions (influenced by LH): a lot of progesterone (by granulosa luteal cells), some estrogen (by theca lutein cells)
\item Help maintain pregnancy, first 3 months before placenta secretes hormones
\item If corpus luteum becomes inactive before 3rd month \(\rightarrow\) abortion
\end{itemize}
\end{itemize}
\begin{center}\includegraphics[width=0.82\textwidth,height=0.68\textheight,keepaspectratio]{\reviewnotesimagepath d6-figure-female-reproductive-physiology-if-corpus-luteum-becomes.png}\par\captionof{figure}{The illustration relates to If corpus luteum becomes inactive before 3rd month \(\rightarrow\) abortion, within the topic of Female Reproductive Physiology.}\labfig{appendix-figure-d6-female-reproductive-physiology-if-corpus-luteum-becomes}\end{center}
\begin{itemize}
\item Menses: withdrawal of hormonal support \(\rightarrow\) endometrium necrosis, menstruation
\end{itemize}
\subsubsection{The Uterine Cycle}
\begin{center}\includegraphics[width=0.82\textwidth,height=0.68\textheight,keepaspectratio]{\reviewnotesimagepath d6-figure-female-reproductive-physiology-uterine-cycle.png}\par\captionof{figure}{The illustration relates to The Uterine Cycle, within the topic of Female Reproductive Physiology.}\labfig{appendix-figure-d6-female-reproductive-physiology-uterine-cycle}\end{center}
\begin{itemize}
\item Menstrual phase:
\begin{itemize}
\item Following ovulation, if no pregnancy
\item On day 26/27 reduction in estrogen, progesterone \(\rightarrow\) endometrium is shed/desquamated \(\rightarrow\) expelled through vagina w/ blood + tissue fluid
\item Duration 4-5 days = menstrual period \(\rightarrow\) only thin (1 mm) layer of endometrium remains
\end{itemize}
\item Proliferative phase
\begin{itemize}
\item Estrogen from ovary \(\rightarrow\) rapid proliferation of endometrium in uterus \(\rightarrow\) endometrium reaches 3-4 mm at end
\item Epithelium first reappears on endometrium in 4-7 days
\item Uterine glands and blood vessels develop
\item Ovulation occurs at day 14 by LH \(\rightarrow\) secretory phase
\end{itemize}
\item Secretory phase
\begin{itemize}
\item Between 15th and 28th day of menstrual cycle (between ovulation and first day of next cycle)
\item Corpus luteum secretes progesterone, estrogen
\begin{itemize}
\item Estrogen leads to proliferation of uterine cells, endometrium thickening
\item Progesterone causes enlarged endometrial stroma, further growth of glands
\item Progesterone leads to endometrial gland secretory function (after completing changes)
\end{itemize}
\end{itemize}
\end{itemize}
\begin{center}\includegraphics[width=0.82\textwidth,height=0.68\textheight,keepaspectratio]{\reviewnotesimagepath d6-figure-female-reproductive-physiology-progesterone-leads-endometrial-gland.png}\par\captionof{figure}{The illustration relates to Progesterone leads to endometrial gland secretory function (after completing changes), within the topic of Female Reproductive Physiology.}\labfig{appendix-figure-d6-female-reproductive-physiology-progesterone-leads-endometrial-gland}\end{center}
\begin{itemize}
\item Implantation of the fertilized egg
\begin{itemize}
\item Day 0: fertilization \(\rightarrow\) day 5: loss of zona pellucida, transport to uterus \(\rightarrow\) day 7-10: implantation in uterine wall
\item Components of fertilized egg
\begin{itemize}
\item Trophoblast: layer surrounding blastocyst, nourish embryo and later form placenta
\item Inner cell mass: located as a cluster at the blastocyst \(\rightarrow\) develops into embryo, isolated from uterine environment by trophoblast
\end{itemize}
\end{itemize}
\end{itemize}
\subsubsection{Functions of the Placenta}
\begin{itemize}
\item Nutrition: transfers nutrients and oxygen to the fetus, transfer waste products/CO2 to the mother
\end{itemize}
\begin{center}\includegraphics[width=0.82\textwidth,height=0.68\textheight,keepaspectratio]{\reviewnotesimagepath d6-figure-female-reproductive-physiology-nutrition-transfers-nutrients-oxygen.png}\par\captionof{figure}{The illustration relates to Nutrition: transfers nutrients and oxygen to the fetus, transfer waste products/CO2 to the mother, within the topic of Female Reproductive Physiology.}\labfig{appendix-figure-d6-female-reproductive-physiology-nutrition-transfers-nutrients-oxygen}\end{center}
\begin{itemize}
\item Endocrine
\begin{itemize}
\item Placenta secretes hormones by syncytiotrophoblast of chorionic villi (outer layer of the placenta)
\item Human chorionic gonadotrophin (hCG)
\begin{itemize}
\item Appears after implantation, indicates pregnancy
\item Resembles LH \(\rightarrow\) maintain the integrity of corpus luteum, promotes continued progesterone secretion (w/ hCG, corpus luteum persists 3-4 months before decline) \(\rightarrow\) placenta starts secreting progesterone, estrogen instead
\item Pregnancy tests detect hCG by antibodies
\end{itemize}
\item Human placental lactogen (hPL)
\begin{itemize}
\item Activate and prepare mammary glands for milk production
\item Mammary glands stimulated by combined hPL, estrogen, progesterone + maternal hormone–growth hormone, prolactin, thyroid hormones
\end{itemize}
\item Relaxin
\begin{itemize}
\item Secreted by placenta, corpus luteum
\item Function: increase flexibility of pubic symphysis (expand during delivery), dilation of cervix (fetus enter vaginal canal), delay labour until late pregnancy
\end{itemize}
\item Estrogen, progesterone
\begin{itemize}
\item Maintains endometrial lining
\item 3rd trimester \(\rightarrow\) estrogen production accelerates \(\rightarrow\) stimulate labour and delivery
\end{itemize}
\item Immune functions
\begin{itemize}
\item Allows passthrough of IgG antibodies \(\rightarrow\) protection of fetus
\item Evade immune attack from mother’s immune system
\item Blood reservoir for fetus
\end{itemize}
\end{itemize}
\end{itemize}
\subsubsection{Hormone Changes During Pregnancy}
\begin{itemize}
\item Fertilization \(\rightarrow\) corpus luteum rescued from regression by human chorionic gonadotrophin (produced by placenta)
\item First trimester \(\rightarrow\) corpus luteum produces estradiol and progesterone, serum hCG doubles every 2 days and peaks at about 10 weeks
\end{itemize}
\begin{center}\includegraphics[width=0.82\textwidth,height=0.68\textheight,keepaspectratio]{\reviewnotesimagepath d6-figure-female-reproductive-physiology-first-trimester-corpus-luteum.png}\par\captionof{figure}{The illustration relates to First trimester \(\rightarrow\) corpus luteum produces estradiol and progesterone, serum hCG doubles every 2 days and peaks at about 10 weeks, within the topic of Female Reproductive Physiology.}\labfig{appendix-figure-d6-female-reproductive-physiology-first-trimester-corpus-luteum}\end{center}
\begin{itemize}
\item Second, third trimesters
\begin{itemize}
\item Progesterone produced by placenta
\item Estrogens \(\rightarrow\) produced by interaction of fetal adrenal glands and placenta (fetal adrenal gland \(\rightarrow\) dehydroepiandrosterone–sulfate \(\rightarrow\) hydroxylation in fetal liver \(\rightarrow\) intermediates transferred to placenta \(\rightarrow\) enzymes remove sulfate, aromatize to estrogens (major = estriol)
\item Human placental lactogen produced throughout pregnancy \(\rightarrow\) similar action to growth hormone, prolactin
\end{itemize}
\item Parturition (childbirth)
\begin{itemize}
\item Throughout pregnancy, progesterone increases threshold for uterine contraction
\item Near term, estrogen : progesterone ratio increases \(\rightarrow\) increased sensitivity to contractile stimuli
\item Initiating event of childbirth = unknown (oxytocin stimulates contractions, but does not increase before childbirth)
\end{itemize}
\item Lactation
\begin{itemize}
\item Estrogen, progesterone stimulate growth + development of breasts during pregnancy
\begin{itemize}
\item Estrogen \(\rightarrow\) ductile system branching
\item Progesterone \(\rightarrow\) alveolar gland development
\item Placental lactogen \(\rightarrow\) breast development
\item Prolactin (ant. pituitary) \(\rightarrow\) milk production (inhibited by placental progesterone)
\end{itemize}
\item Prolactin increases steadily during pregnancy (estrogen increases secretion of prolactin from anterior pituitary gland)
\item Lactation inhibited during pregnancy because estrogen, progesterone block action of prolactin
\item Lactation is maintained by suckling (stimulates oxytocin, prolactin)
\item Ovulation suppressed with continued lactation because
\begin{itemize}
\item Prolactin inhibits hypothalamic GnRH secretion, GnRH actions on anterior pituitary \(\rightarrow\) prevents LH, FSH secretion
\item Antagonist of LH and FSH on ovaries
\end{itemize}
\end{itemize}
\end{itemize}
\begin{center}\includegraphics[width=0.82\textwidth,height=0.68\textheight,keepaspectratio]{\reviewnotesimagepath d6-figure-female-reproductive-physiology-antagonist-lh-fsh-ovaries.png}\par\captionof{figure}{The illustration relates to Antagonist of LH and FSH on ovaries, within the topic of Female Reproductive Physiology.}\labfig{appendix-figure-d6-female-reproductive-physiology-antagonist-lh-fsh-ovaries}\end{center}
\begin{center}\includegraphics[width=0.82\textwidth,height=0.68\textheight,keepaspectratio]{\reviewnotesimagepath d6-figure-female-reproductive-physiology-female-reproductive-physiology-10.png}\par\captionof{figure}{The illustration relates to Female Reproductive Physiology, within the topic of Female Reproductive Physiology.}\labfig{appendix-figure-d6-female-reproductive-physiology-female-reproductive-physiology-10}\end{center}
\subsubsection{Menopause}
\begin{itemize}
\item Permanent cessation of menstruation (typical age 45-55)
\item Measurable with menopause test
\end{itemize}
\begin{center}\includegraphics[width=0.82\textwidth,height=0.68\textheight,keepaspectratio]{\reviewnotesimagepath d6-figure-female-reproductive-physiology-measurable-menopause-test.png}\par\captionof{figure}{The illustration relates to Measurable with menopause test, within the topic of Female Reproductive Physiology.}\labfig{appendix-figure-d6-female-reproductive-physiology-measurable-menopause-test}\end{center}
\begin{itemize}
\item Causes of menopause
\begin{itemize}
\item Natural menopause
\begin{itemize}
\item Reduced ovarian response to FSH, LH
\item Altered GnRH release from hypothalamus
\item Pituitary gland becomes less responsive to GnRH
\end{itemize}
\item Induced menopause
\begin{itemize}
\item Hysterectomy
\item Radiotherapy
\end{itemize}
\end{itemize}
\item Phases of menopause
\end{itemize}
\begin{center}\includegraphics[width=0.82\textwidth,height=0.68\textheight,keepaspectratio]{\reviewnotesimagepath d6-figure-female-reproductive-physiology-phases-menopause.png}\par\captionof{figure}{The illustration relates to Phases of menopause, within the topic of Female Reproductive Physiology.}\labfig{appendix-figure-d6-female-reproductive-physiology-phases-menopause}\end{center}
\begin{center}\includegraphics[width=0.82\textwidth,height=0.68\textheight,keepaspectratio]{\reviewnotesimagepath d6-figure-female-reproductive-physiology-female-reproductive-physiology-11.png}\par\captionof{figure}{The illustration relates to Female Reproductive Physiology, within the topic of Female Reproductive Physiology.}\labfig{appendix-figure-d6-female-reproductive-physiology-female-reproductive-physiology-11}\end{center}

\subsection{Common diseases in urology, male genital organs}
\begin{itemize}
\item Outline
\begin{itemize}
\item Congenital disorders
\item Trauma
\item Infection, inflammation disorders
\item Neoplastic disorders and Tumor-like conditions
\end{itemize}
\end{itemize}
\subsubsection{Congenital disorders}
\begin{itemize}
\item Hypospadias and epispadias
\begin{itemize}
\item Abnormal opening of urethra (hypospadias below and epispadias above)
\item Rare, assoc w/ bladder exstrophy (incomplete development of urethra in utero)
\item Occurs in 1:100 male children, increased risk if familial history
\end{itemize}
\end{itemize}
\begin{center}\includegraphics[width=0.82\textwidth,height=0.68\textheight,keepaspectratio]{\reviewnotesimagepath d6-figure-common-diseases-urology-occurs-100-male-children.png}\par\captionof{figure}{The illustration relates to Occurs in 1:100 male children, increased risk if familial history, within the topic of Common diseases in urology, male genital organs.}\labfig{appendix-figure-d6-common-diseases-urology-occurs-100-male-children}\end{center}
\begin{center}\includegraphics[width=0.82\textwidth,height=0.68\textheight,keepaspectratio]{\reviewnotesimagepath d6-figure-common-diseases-urology-common-diseases-urology-male.png}\par\captionof{figure}{The illustration relates to Common diseases in urology, male genital organs, within the topic of Common diseases in urology, male genital organs.}\labfig{appendix-figure-d6-common-diseases-urology-common-diseases-urology-male}\end{center}
\begin{itemize}
\item Treatment: check for normal urine flow, refer urology and delay circumcision (for reconstructive surgery)
\item Phimosis and paraphimosis
\begin{itemize}
\item Phymosis = orifice of prepuce too small for retraction over penis. Tx = circumcision
\item Paraphimosis = foreskin trapped behind penis
\begin{itemize}
\item Risk - frequent catheterization w/o reducing foreskin \(\rightarrow\) paraphimosis, forcible retraction of phimosis, frequent sexual activity
\item Treatment (if ineffective, refer to surgery) = wrap swelling in gauze, apply increasing pressure to squeeze fluid out (10-15 mins of pressure) \(\rightarrow\) replace foreskin \(\rightarrow\) circumcision later
\end{itemize}
\end{itemize}
\item Undescended testis
\begin{itemize}
\item Failure of testes to descend into scrotum (but still in normal path, undeveloped)
\item Most common genital disorder identified at birth
\item Premature infant = 33\%, full term birth = 2-4\%, at 1 year 1\%
\item Alternate Dx: ectopic testes: testes deviate from normal path of decent
\end{itemize}
\end{itemize}
\begin{center}\includegraphics[width=0.82\textwidth,height=0.68\textheight,keepaspectratio]{\reviewnotesimagepath d6-figure-common-diseases-urology-alternate-dx-ectopic-testes.png}\par\captionof{figure}{The illustration relates to Alternate Dx: ectopic testes: testes deviate from normal path of decent, within the topic of Common diseases in urology, male genital organs.}\labfig{appendix-figure-d6-common-diseases-urology-alternate-dx-ectopic-testes}\end{center}
\begin{itemize}
\item Hydrocele
\begin{itemize}
\item Enlarged scrotum
\item Usu. spontaneous resolution (if present > 1 yr, refer to surgery)
\item Caused by collection of fluid in tunica vaginalis
\item Congenital cause = patent processus, peritoneal fluid decent
\item Primary = large, hard tense and slow occurrence w/o defined cause (age > 40)
\item Secondary = rapid development, small, w/o underlying cause (age 20-40)
\end{itemize}
\item Spermatocele
\end{itemize}
\begin{center}\includegraphics[width=0.82\textwidth,height=0.68\textheight,keepaspectratio]{\reviewnotesimagepath d6-figure-common-diseases-urology-spermatocele.png}\par\captionof{figure}{The illustration relates to Spermatocele, within the topic of Common diseases in urology, male genital organs.}\labfig{appendix-figure-d6-common-diseases-urology-spermatocele}\end{center}
\begin{itemize}
\item Asymptomatic, usu. incidental finding
\item Painless mobile swelling posterio–superiorly
\item Assoc w/ heavy scrotum, dull discomfort
\item Causes: depend on age
\begin{itemize}
\item Children: rare, usu assoc w/ congenital genitourinary abnormality, infection with gram -ve rods
\item Sexually active M < 35: assoc chlamydia trachomatis, Neisseria
\end{itemize}
\end{itemize}
\subsubsection{Inflammation/Infection}
\begin{itemize}
\item Orchiditis
\begin{itemize}
\item Swelling, pain, tenderness, erythema and shiny overlying skin
\item Usu spread from epididymitis, hematogenous or post–viral
\item Viral: Mumps, coxsackie, echovirus, parvovirus
\item Bacterial: brucellosis
\end{itemize}
\item Treatment of epididimo–orchitis
\begin{itemize}
\item Bacterial etiology
\begin{itemize}
\item < 35 y/o and sexually active: antibiotic for STD (esp. gonorrhea, chlamydia) \(\rightarrow\) Ceftriaxone, doxycycline, azithromycin
\item > 35 y/o: bacteria etiology requires added cover for other gram -ve bacteria \(\rightarrow\) fluoroquinolone (not used in gonorrhea), Trimethoprim/sulfamethoxazole \(\rightarrow\) no orchidectomy
\end{itemize}
\end{itemize}
\end{itemize}
\subsubsection{Trauma}
\begin{itemize}
\item Kydney
\begin{itemize}
\item S\&S: flank pain, gross hematuria, BP drop
\item Investigation: KUB film, U/S, IVP (intravenous pyelogram = contrast film at set times after injection), CT KUB
\end{itemize}
\end{itemize}
\begin{center}\includegraphics[width=0.82\textwidth,height=0.68\textheight,keepaspectratio]{\reviewnotesimagepath d6-figure-common-diseases-urology-investigation-kub-film-ivp.png}\par\captionof{figure}{The illustration relates to Investigation: KUB film, U/S, IVP (intravenous pyelogram = contrast film at set times after injection), CT KUB, within the topic of Common diseases in urology, male genital organs.}\labfig{appendix-figure-d6-common-diseases-urology-investigation-kub-film-ivp}\end{center}
\begin{itemize}
\item Tx: grades 1-3 conservative, 4 up to patient, 5 emerg surgery
\item Ureter
\begin{itemize}
\item S\&S: urine leakage, periotonitis., hydronephrosis
\item Investigation: KUB film, U/S, IVP (interval IV pyelogram), CT KUB
\item Positive finding = contrast CT with finding of extravasation (urine leak) into medial perineal space, absence of contrast in distal ureter (if transected)
\item Additional: retrotrade pyelogram, one-shot IVP
\item Typical mechanism: surgical injury
\item Treatment: emergent repair, double J stent in ureter
\end{itemize}
\item Bladder
\begin{itemize}
\item S\&S: lower abdominal pain, peritonitis, gross hematuria
\item Investigation CT cystogram/fluoroscopic cystogram w/ 200-300 mL dilute contrast
\item Finding: either intraperitoneum or extraperitoneum bladder injury
\item Treatment
\begin{itemize}
\item Intraperitoneal = emergency surgical repair
\item Extraperitoneal = prolonged urethral catheterization
\end{itemize}
\end{itemize}
\end{itemize}
\begin{center}\includegraphics[width=0.82\textwidth,height=0.68\textheight,keepaspectratio]{\reviewnotesimagepath d6-figure-common-diseases-urology-extraperitoneal-prolonged-urethral-catheterization.png}\par\captionof{figure}{The illustration relates to Extraperitoneal = prolonged urethral catheterization, within the topic of Common diseases in urology, male genital organs.}\labfig{appendix-figure-d6-common-diseases-urology-extraperitoneal-prolonged-urethral-catheterization}\end{center}
\begin{itemize}
\item Urethra
\end{itemize}
\begin{center}\includegraphics[width=0.82\textwidth,height=0.68\textheight,keepaspectratio]{\reviewnotesimagepath d6-figure-common-diseases-urology-urethra.png}\par\captionof{figure}{The illustration relates to Urethra, within the topic of Common diseases in urology, male genital organs.}\labfig{appendix-figure-d6-common-diseases-urology-urethra}\end{center}
\begin{itemize}
\item Mechanism of injury
\begin{itemize}
\item Posterior urethral injury = commonly pelvic fractures \(\rightarrow\) stretch/contusion to complete disruption
\item Anterior urethral injury = road traffic accidents, falls, straddle injury
\item Iatrogenic = endoscope trauma \(\rightarrow\) most common cause of stricture
\end{itemize}
\item S\&S: blood at urethral meatus (exit point of urine), hematuria, difficulty/inability to void, distended + palpable bladder, high-riding prostate, displaced fracture at pubic rami, pelvic hematoma, butterfly bruise at peritoneum, blood on digital rectal examination
\item Treatment (location dependant)
\begin{itemize}
\item Anterior urethra \(\rightarrow\) foley catheter attempt \(\rightarrow\) suprapubic cystostomy
\item Posterior urethra \(\rightarrow\) suprapubic cystostomy
\end{itemize}
\end{itemize}
\subsubsection{Neoplasms of KUB System}
\begin{itemize}
\item Benign prostate hyperplasia
\begin{itemize}
\item S\&S: poor flow (decreased force/caliber of stream), hesitancy/straining, terminal dribbling, intermittent stream, incomplete evacuation of bladder, double voiding, post void dribbling, urine retention (acute or chronic)
\item Irritation–associated S\&S: urgency, frequency (+nocturia), urge incontinence (enuresis, no control of voiding)
\item Complications of BPH
\begin{itemize}
\item Obstructive uropathy
\item Bladder hypertrophy
\item Trabeculation
\item Diverticuli
\item Hydroureter (bilateral)
\item Hydronephrosis
\item Lithiasis/stone
\item Secondary infection
\item NOT risk for prostatic carcinoma
\end{itemize}
\item Treatment
\begin{itemize}
\item Alpha blocker - relaxes prostate, bladder neck smooth muscle by blocking sympathetic activity
\item 5-alpha reductase inhibitors - improve symptoms, increase flow and prevent BPH by reducing hormonal prostate enlargement (testosterone, dihydroxytestosterone)
\item Retain Foley catheter \(\rightarrow\) fail med/complication? \(\rightarrow\) surgery
\end{itemize}
\end{itemize}
\item Benign renal mass
\begin{itemize}
\item Solid = angiomyolipoma
\begin{itemize}
\item Investigations: U/S KUB, CT KUB
\item Tx F/U by U/S or CT
\end{itemize}
\item Cystic = renal cyst
\end{itemize}
\end{itemize}
\begin{center}\includegraphics[width=0.82\textwidth,height=0.68\textheight,keepaspectratio]{\reviewnotesimagepath d6-figure-common-diseases-urology-cystic-renal-cyst.png}\par\captionof{figure}{The illustration relates to Cystic = renal cyst, within the topic of Common diseases in urology, male genital organs.}\labfig{appendix-figure-d6-common-diseases-urology-cystic-renal-cyst}\end{center}
\begin{center}\includegraphics[width=0.82\textwidth,height=0.68\textheight,keepaspectratio]{\reviewnotesimagepath d6-figure-common-diseases-urology-common-diseases-urology-male-02.png}\par\captionof{figure}{The illustration relates to Common diseases in urology, male genital organs, within the topic of Common diseases in urology, male genital organs.}\labfig{appendix-figure-d6-common-diseases-urology-common-diseases-urology-male-02}\end{center}
\subsubsection{Common Malignancies of KUB System}
\begin{itemize}
\item Urothelial cancer
\begin{itemize}
\item Renal pelvis, ureter, bladder (most common)
\item Detection by gross, painless hematuria (85\% of patients)
\item Of all gross hematuria \(\rightarrow\) 50\% with demonstrable cause, 20\% urologic malignancy, 12\% bladder tumor
\item In microscopic hematuria: cystoscopy, upper tract imaging, urine cytology
\item Bladder cancers
\begin{itemize}
\item Superficial (papillary) growth
\item Invasive - sessile mass
\item Risk factors: male (4:1), age 60-80, smoking (3x risk, worse progression but reduced after quitting), occupation (paint, print, iron/Al processing, gas/tar manufacture, aromatic amine), pelvic radiation chronic inflammation (bladder stone, long term foley, snail fever (parasite in urinary tract), some drugs (phenacetin, cyclophosphamide)
\end{itemize}
\end{itemize}
\end{itemize}
\begin{center}\includegraphics[width=0.82\textwidth,height=0.68\textheight,keepaspectratio]{\reviewnotesimagepath d6-figure-common-diseases-urology-risk-factors-male-age.png}\par\captionof{figure}{The illustration relates to Risk factors: male (4:1), age 60-80, smoking (3x risk, worse progression but reduced after quitting), occupation (paint, print, iron/Al processing, gas/tar manu, within the topic of Common diseases in urology, male genital organs.}\labfig{appendix-figure-d6-common-diseases-urology-risk-factors-male-age}\end{center}
\begin{itemize}
\item Treatments of CA bladder
\begin{itemize}
\item Stage 0, 1: transurethral resection, chemotherapy, immunotherapy
\item Stage 2, 3: radical cystectomy, partial cystectomy + radiation, chemotherapy
\item Stage 4: no cure (possible chemotherapy)
\end{itemize}
\item Renal cell carcinoma
\begin{itemize}
\item S\&S: haematuria (gross = 40\%), flank pain (40\%), palpable mass flank/abdomen (25\%), weight loss, fever (20\%), hypertension, hypercalcemia, night sweats, malaise, varicocele/scrotal varicose veins (obstructed testicular vein), weight loss/constitutional (33\%), paraneoplastics (30\%)
\item Diagnosis
\begin{itemize}
\item Usu indidental (50\%)
\item Presenting as metastasis = 30\% \(\rightarrow\) bone pain, cough
\item Gross hematuria, flank pain, abdomen mass \(\rightarrow\) patients with all symptoms < 10\%
\end{itemize}
\item Investigation: US, CT abdomen
\item Tx: radical/partial nephrectomy
\end{itemize}
\item Penile cancer
\begin{itemize}
\item TX: low grade (1, 2) \(\rightarrow\) organ preserved, maintain sexual quality
\begin{itemize}
\item Limited excision, Mohs surgery, laser ablation
\item Amputation for large/invasive lesions \(\rightarrow\) rapid tumor control
\end{itemize}
\end{itemize}
\item Testicular cancer
\begin{itemize}
\item Risk factors: cryptorchidism (failure of testes to descend), germ cell tumor, intratubular germ cell neoplasia, family history, HIV (seminoma risk), Klinefelter syndrome (mediastinal germ cell tumor), down syndrome, androgen insensitivity/mixed gonadal dysgenesis
\item Epidemiology:
\begin{itemize}
\item Age: 2-4, 20-40 (germ cell tumor), > 50 (lymphoma)
\item Common in young white, less in african american
\item Incidence: high in denmark, norway, switzerland; low in eastern europe, asia
\end{itemize}
\item Investigation: AFP (high in pure embryonic carcinoma., teratoma, yolk sac tumor, but not pure choriocarcinoma- placenta/trophoblast cells), B-HCG (embryonal/choriocarcinoma - age 40-60 , seminoma - age 5-10), LDH (advanced germ cell tumor \(\rightarrow\) seminoma)
\item Treatment: radical orchidectomy (NO biopsy)
\end{itemize}
\item Prostate cancer
\end{itemize}
\begin{center}\includegraphics[width=0.82\textwidth,height=0.68\textheight,keepaspectratio]{\reviewnotesimagepath d6-figure-common-diseases-urology-prostate-cancer.png}\par\captionof{figure}{The illustration relates to Prostate cancer, within the topic of Common diseases in urology, male genital organs.}\labfig{appendix-figure-d6-common-diseases-urology-prostate-cancer}\end{center}
\begin{itemize}
\item Diagnosis: digital rectal exam nodular hard prostate, PSA > 4.0 ng/mL, transrectal ultrasound biopsy
\item Tx: radical prostatectomy, hormone therapy, bilateral orchidectomy (reduce testosterone, DHT), radiation/chemotherapy
\end{itemize}
\subsubsection{Other Causes}
\begin{itemize}
\item Torsion testis
\end{itemize}
\begin{center}\includegraphics[width=0.82\textwidth,height=0.68\textheight,keepaspectratio]{\reviewnotesimagepath d6-figure-common-diseases-urology-torsion-testis.png}\par\captionof{figure}{The illustration relates to Torsion testis, within the topic of Common diseases in urology, male genital organs.}\labfig{appendix-figure-d6-common-diseases-urology-torsion-testis}\end{center}
\begin{itemize}
\item Treatment: surgical intervention (for cases where suspected torsion/can’t exclude torsion)
\item Dx diag: torsed appendix testes
\begin{itemize}
\item Mimics testicular torsion
\item PE + for blue dot sign
\item Normal testes positionMay be tender/inflamed testes + point tender at appendix testes
\item Investigate with u/s to access perfusion
\item Tx: by symptom, conservative + urologic assessment
\end{itemize}
\end{itemize}
\begin{center}\includegraphics[width=0.82\textwidth,height=0.68\textheight,keepaspectratio]{\reviewnotesimagepath d6-figure-common-diseases-urology-tx-symptom-conservative-urologic.png}\par\captionof{figure}{The illustration relates to Tx: by symptom, conservative + urologic assessment, within the topic of Common diseases in urology, male genital organs.}\labfig{appendix-figure-d6-common-diseases-urology-tx-symptom-conservative-urologic}\end{center}
\begin{itemize}
\item Acute urinary retention
\begin{itemize}
\item Unable to empty bladder
\end{itemize}
\end{itemize}
\begin{center}\includegraphics[width=0.82\textwidth,height=0.68\textheight,keepaspectratio]{\reviewnotesimagepath d6-figure-common-diseases-urology-unable-empty-bladder.png}\par\captionof{figure}{The illustration relates to Unable to empty bladder, within the topic of Common diseases in urology, male genital organs.}\labfig{appendix-figure-d6-common-diseases-urology-unable-empty-bladder}\end{center}
\begin{itemize}
\item Causes
\begin{itemize}
\item Men: prostate enlargement, urethral stricture
\item Women: pelvic organ prolapse, urethral stricture, neurogenic bladder
\item Both genders: sex, hematuria, drugs, pain, pelvic organ surgery, fractured pelvis, diabetic, multiple sclerosis
\end{itemize}
\item Treatment = emergency drainage by: urethral catheter, suprapubic puncture, suprapubic cystostomy (for foreign body removal/prolonged drainage e.g. ruptured urethra/severe stricture)
\item Obstructive uropathy
\begin{itemize}
\item Structural/functional abnormality restricting urine flow \(\rightarrow\) urine backflow to kidney at any point in urinary tract
\item Back pressure causes hydronephrosis (dilation of renal calyx) \(\rightarrow\) obstructive nephropathy (functional/anatomic renal damage)
\item Etiology (by age)
\begin{itemize}
\item Children: urethral valve, stricture/stenosis at urethrovesicle or ureteropelvic junction
\item Young adults: calculus
\item Older adults: BPH, prostate cancer, retroperitoneal/pelvic tumors, calculi
\end{itemize}
\end{itemize}
\end{itemize}
\begin{center}\includegraphics[width=0.82\textwidth,height=0.68\textheight,keepaspectratio]{\reviewnotesimagepath d6-figure-common-diseases-urology-older-adults-bph-prostate.png}\par\captionof{figure}{The illustration relates to Older adults: BPH, prostate cancer, retroperitoneal/pelvic tumors, calculi, within the topic of Common diseases in urology, male genital organs.}\labfig{appendix-figure-d6-common-diseases-urology-older-adults-bph-prostate}\end{center}

\section{Medication in Pregnancy}

\subsection{Medications in Pregnancy}
\subsubsection{Drugs in Pregnancy}
\begin{itemize}
\item Key concepts
\begin{itemize}
\item Drug induced teratogenicity = concern, but low risk (only account for 4-5\% of the total risk, compared to 65\% polygenic/multifactorial)
\item Risk of birth defects is highest in organogenesis (first trimester day 14 until 11th week after LMP)
\item Folic acid, multivitamins \(\rightarrow\) prevent congenital abnormalities (0.4 to 0.8 mg of folic acid daily for first 3 months, 4 mg if high risk)
\item Most drugs safe during breastfeeding
\item When possible, Tx using non–pharmacologic agents
\item If Rx needed in pregnancy, evaluate risk–benifit
\end{itemize}
\item 85\% of pregnancies \(\rightarrow\) nausea, vomiting. Treatment:
\begin{itemize}
\item Nonpharmacologic: lifestyle (avoid triggers, rest), change diet (small, frequent meals, fluid restrict, blander food)
\item First line med = \textbf{pyridoxine} (B6) += doxylamine
\item Secondline: \textbf{dimenhydrinate}, metoclopramide, diphenhydramine, or ondansetron
\end{itemize}
\item Constipation \(\rightarrow\) treatment:
\begin{itemize}
\item Nonpharmacologic: high fiber, fluids, regular exercise, avoid prolonged time on toilet
\item Hemorrhoids: soak in sitz bath, apply ice \(\rightarrow\) acetaminophen, topical analgesics
\item Bulk forming laxatives (psylium, calcium polycarbophil) - firstline Rx
\end{itemize}
\item Anemia in pregnancy = Hb < 11 g/dL \(\rightarrow\) fatigue, palpitation, infection, decreased exercise tolerance + fetal risk (premature, low birth weight, perinatal death)
\item Nasal congestion/cough treatment
\begin{itemize}
\item Nonpharmacologic: rest, fluid, humidified air, nasal saline + acetaminophen for fever
\item NO ORAL DECONGESTANTS. \textbf{[Accuracy check: This is an absolute medication restriction; specify the drug, trimester, and clinical context or verify against current pregnancy guidance.] }
\item Can use: corticosteroids (chronic), 1st/2nd generation antihistamine
\end{itemize}
\item Preterm labour
\begin{itemize}
\item Tx: antenatal corticosteroids (dexamethasone) 1 dose to promote surfactant development in fetal lung
\end{itemize}
\item Labor pain: tocolytics
\begin{itemize}
\item Nifedipine => best benefit-to-risk ratio, ADR: headache, flushing, dizziness, hypotension
\item Magnesium sulfate => unproven efficacy, ADR: maternal pulmonary edema, cardiac arrest
\item Terbutaline (beta2 agonist) => can prolong pregnancy but cannot decrease neonatal morbidity, ADR: cardiac arrhythmia, palpitation, chest pain
\item Indomethacin (NSAIDs) => can prolong pregnancy but cannot decrease
\end{itemize}
\item For induction of labour
\begin{itemize}
\item Indication
\begin{itemize}
\item Membrane rupture
\item Fetal growth restriction
\item Uteroplacental insufficiency as in diabetes, pre-eclampsia or eclampsia
\item Gestation beyond 42 weeks
\item Fetal lung must be sufficiently mature. \textbf{[Accuracy check: Verify this as a general prerequisite for induction; the indication and gestational age determine the clinical decision.] }
\end{itemize}
\item Contraindication \(\rightarrow\) abnormal fetal position, evidence of fetal distress, placental abnormalities, previous uterine surgery
\item Drugs
\begin{itemize}
\item Prostaglandins orally or vaginally stimulate uterine contraction => used in postpartum hemorrhage (PPH)
\item PGE1, PGE2, PGF2a => can facilitate labor by promoting ripening and dilation of the cervix
\item Examples of drugs => dinoprostone (PGE2), misoprostol (PGE1)
\item + Oxytocin or Methylergonovine is used for induction of labor, augmentation of labor, and prophylaxis/treatment of PPH T1/2 = 5 min. \textbf{[Accuracy check: The source groups methylergonovine with labor-induction agents; verify this carefully against current prescribing information and separate intrapartum from postpartum use.] }
\end{itemize}
\end{itemize}
\item Medication in lactation
\begin{itemize}
\item Ideal = mother breastfeed for 1 year (75\% start at birth, 17\% start at 6 months). \textbf{[Accuracy check: Verify the recommendation’s duration and source; breastfeeding guidance may differ across organizations and has changed over time.] }
\item Enhance lactation by: optimize technique
\item Dopamine antagonist (increase prolactin) \(\rightarrow\) ADR: fatigue, abdomen pain, depression, extrapyramidal symptoms in mother
\end{itemize}
\end{itemize}
\subsubsection{Other Content}
\begin{itemize}
\item Pregnancy category of drugs. \textbf{[Accuracy check: The image appears to use the former A/B/C/D/X letter-category system; verify against current pregnancy-labeling standards before using these categories clinically.] }
\end{itemize}
\begin{center}\includegraphics[width=0.82\textwidth,height=0.68\textheight,keepaspectratio]{\reviewnotesimagepath d6-figure-medications-pregnancy-pregnancy-category-drugs.png}\par\captionof{figure}{The illustration relates to Pregnancy category of drugs, within the topic of Medications in Pregnancy.}\labfig{appendix-figure-d6-medications-pregnancy-pregnancy-category-drugs}\end{center}
\begin{itemize}
\item Organogenesis \(\rightarrow\) highest susceptibility to complications (structural abnormalities, spontaneous abortion)
\item Congenital defects
\begin{itemize}
\item Neural tube defects (spina bifida, anencephaly - fail to close \(\rightarrow\) neural tissue necrosis, encephalocele - brain hernia): prevented by folic acid
\item Limb defects (amelia - absent, meromelia - absent part, phocomelia - poorly formed) \(\rightarrow\) causes: genetic (autosomal recessive), medication (thalidomide), antiemetic
\item Cleft lip/palette \(\rightarrow\) causes: genetic (autosomal recessive), maternal smoking, medications: anticonvulsants, benzodiazepines (sleeping pill)
\end{itemize}
\end{itemize}
\begin{center}\includegraphics[width=0.82\textwidth,height=0.68\textheight,keepaspectratio]{\reviewnotesimagepath d6-figure-medications-pregnancy-cleft-lip-palette-causes.png}\par\captionof{figure}{The illustration relates to Cleft lip/palette \(\rightarrow\) causes: genetic (autosomal recessive), maternal smoking, medications: anticonvulsants, benzodiazepines (sleeping pill), within the topic of Medications in Pregnancy.}\labfig{appendix-figure-d6-medications-pregnancy-cleft-lip-palette-causes}\end{center}
\begin{center}\includegraphics[width=0.82\textwidth,height=0.68\textheight,keepaspectratio]{\reviewnotesimagepath d6-figure-medications-pregnancy-medications-pregnancy.png}\par\captionof{figure}{The illustration relates to Medications in Pregnancy, within the topic of Medications in Pregnancy.}\labfig{appendix-figure-d6-medications-pregnancy-medications-pregnancy}\end{center}
\begin{center}\includegraphics[width=0.82\textwidth,height=0.68\textheight,keepaspectratio]{\reviewnotesimagepath d6-figure-medications-pregnancy-medications-pregnancy-02.png}\par\captionof{figure}{The illustration relates to Medications in Pregnancy, within the topic of Medications in Pregnancy.}\labfig{appendix-figure-d6-medications-pregnancy-medications-pregnancy-02}\end{center}
\end{document}